\documentclass[10pt]{article}
\usepackage[utf8]{inputenc}
\usepackage[T1]{fontenc}
\usepackage{graphicx}
\usepackage[export]{adjustbox}
\graphicspath{ {./images/} }
\usepackage{hyperref}
\hypersetup{colorlinks=true, linkcolor=blue, filecolor=magenta, urlcolor=cyan,}
\urlstyle{same}
\usepackage{amsmath}
\usepackage{amsfonts}
\usepackage{amssymb}
\usepackage[version=4]{mhchem}
\usepackage{stmaryrd}
\usepackage{bbold}
\usepackage{mathrsfs}
\usepackage{caption}
\usepackage{multirow}

%New command to display footnote whose markers will always be hidden
\let\svthefootnote\thefootnote
\newcommand\blfootnotetext[1]{%
  \let\thefootnote\relax\footnote{#1}%
  \addtocounter{footnote}{-1}%
  \let\thefootnote\svthefootnote%
}
%Overriding the \footnotetext command to hide the marker if its value is `0`
\let\svfootnotetext\footnotetext
\renewcommand\footnotetext[2][?]{%
  \if\relax#1\relax%
    \ifnum\value{footnote}=0\blfootnotetext{#2}\else\svfootnotetext{#2}\fi%
  \else%
    \if?#1\ifnum\value{footnote}=0\blfootnotetext{#2}\else\svfootnotetext{#2}\fi%
    \else\svfootnotetext[#1]{#2}\fi%
  \fi
}
\DeclareUnicodeCharacter{22C5}{\ifmmode\cdot\else{$\cdot$}\fi}
\DeclareUnicodeCharacter{0394}{\ifmmode\Delta\else{$\Delta$}\fi}
\DeclareUnicodeCharacter{22EE}{\ifmmode\vdots\else{$\vdots$}\fi}
\title{Chapter 19: Censored Data, Sample Selection, and Attrition}

\author{}
\date{}

\begin{document}
\maketitle

\subsection*{19.1 Introduction}
In previous chapters we assumed that we can obtain a random sample from the population of interest. For example, in Part II, where we studied models linear in the parameters, we assumed that data on the dependent variable, the explanatory variables, and instrumental variables can be obtained by means of random samplingwhether in a cross section or panel data context. In earlier chapters of Part IV we studied various nonlinear models for response variables that are limited in some way. Chapter 15 extensively considered binary response models, and we saw that the most commonly used models imply nonconstant partial effects. The same is true for corner solution responses in Chapter 17 and those for count and fractional responses in Chapter 18. It is critical to understand that our reason for looking beyond linear models in those chapters is to obtain functional forms that are more realistic than models that are linear in parameters.

In this chapter we turn to several missing data problems. It is critical-even more so than in the previous chapters-to distinguish between assumptions placed on the population model and assumptions about how the data were generated. Under random sampling, the interesting issues concern the population model and assumptions we make about distributional features in the population. With nonrandom sampling, we must take particular care in stating assumptions about the population and separately stating assumptions on the sampling scheme.

We first study the general problem of data censoring. With data censoring we can still randomly sample units from the relevant population, but we face the problem that one or more of the variables is censored: we only observe the response over a certain range-sometimes a very limited range and sometimes a broad range. We cover several examples of data censoring in Section 19.2.

In addition to data censoring, we also treat the general problem of sample selection. Section 19.3 begins with a general discussion of examples of missing data schemes, and Section 19.4 establishes when the missing data problem can be ignored without resulting in inconsistent estimators. With sample selection problems, we may or may not be able to randomly sample units from the population. The case of data truncation, which we study in Section 19.5, is a situation where we do not randomly draw units from the population; rather, we randomly sample from a subpopulation defined in terms of one or more of the observed variables. Naturally, the population parameters cannot always be identified with such sampling schemes, but they can be under suitable assumptions.

Another sample selection problem is incidental truncation, where certain variables are observed only if other variables take on particular values. In such cases, we often\\
can randomly sample units, but we have data missing on key variables (and, unlike in the data censoring case, we have no information on the outcome, or even a range of possible outcomes, of the missing data). We treat this case in Sections 19.6 and 19.7.

Most of the methods that allow for sample selection to be systematically correlated with unobservables rely on linearity or some other simple response function (such as exponential). An alternative approach is inverse probability weighting, which can be applied to general missing data problems if we have good enough predictors of selection so that, conditional on those predictors, selection is appropriately exogenous. We cover inverse probability weighting in the context of M-estimation in Section 19.8. Section 19.9 covers sample selection, including the specific problem of attrition, in panel data applications.

Before we treat specific censoring and selection schemes, it is important to understand the notational conventions in this chapter. As in previous chapters, we continue to let $y$ denote the response variable in the population of interest. Therefore, we will write population models with $y$ as the dependent variable. The distinction between the underlying response variable of interest, and what we can observe, is critical throughout this chapter. For example, in Section 19.2 we consider the case where we can randomly sample from the population but we only observe a censored version of $y$. If $w$ denotes the censored outcome, then for a randomly drawn unit $i$ we observe $w_{i}$, but this may not equal $y_{i}$.

For models with endogenous explanatory variables, we continue to use the conventions of the previous chapters, letting, for example, $y_{1}$ and $y_{2}$ denote the endogenous variables in the underlying population (where $y_{1}$ is typically the response variable and $y_{2}$ is an endogenous explanatory variable). In principle, each of these could be censored, and we use $w_{1}$ and $w_{2}$, respectively, to denote the censored outcomes.

\subsection*{19.2 Data Censoring}
Traditionally, data censoring problems are treated within the frameworks of Chapters 15-17-alongside binary responses, multinomial responses, and corner solution responses. The benefit of a parallel treatment is that it economizes on the presentation, but it has significant costs: empirical researchers tend to let the two very different issues of functional form for limited dependent variables and data censoring problems blend together. Thus, although the statistical models used for limited dependent variables and censored dependent variables are similar, their interpretation, as well as the way one views and reacts to violations of standard assumptions, is very different. Because the statistical tools for handling censored data are very similar to\\
specifying and estimating models for limited dependent variables, our coverage of the estimation details can be terse.

As an example, consider the problem of top coding, where a variable is reported only up to a specified ceiling. For outcomes above the ceiling, all we know is that the outcome was above the ceiling. Common examples are survey data on wealth and income. In order to elicit responses from wealthy people, some surveys only ask about the amount of wealth up to a given threshold, allowing wealthy people to simply indicate if their wealth is above the threshold. Probably the underlying population model in this case is a standard linear model, but that is a separate issue. To emphasize this point, we might also use top coding when collecting data on charitable contributions. As we saw in Chapter 17, in a large population the charitable contributions is best described as a corner solution outcome, with corner at zero. A sensible population model for charitable contributions is the Tobit model, or perhaps one of the two-part models we covered in Section 17.6. Specification of this population model is separate from the data collection scheme. If we top code contributions at, say, $\$ 10,000$ (where contributions are measured in $\$ 1,000$ s), then the reported contributions data will appear to have two corners: one at zero and one at 10 . Of course, these are very different in nature: when we observe a zero, we know that the person had zero charitable contributions. If we observe 10 , we only know that the contributions were at least $\$ 10,000$. If the top coding were instead at $\$ 20,000$, nothing would happen to the population distribution of contributions; it still has a corner at zero. But the top coding changes the upper corner, and by increasing the censoring value we observe more of the population distribution of contributions. As we will see in Section 19.2.3, the statistical model for estimating the parameters when a Type I Tobit variable is top coded is equivalent to a two-limit Tobit; but the underlying model is the standard Type I Tobit model.

Although we will touch on more complicated situations, such as the Tobit model just described, our main focus on this chapter is on a very familiar model: the linear model from introductory econometrics. Let $y$ denote the variable of interest, and assume that it follows a standard linear model in the population:


\begin{equation*}
y=\mathbf{x} \boldsymbol{\beta}+u \tag{19.1}
\end{equation*}


$\mathrm{E}(u \mid \mathbf{x})=0$,\\
where $\mathbf{x}$ is $1 \times K$ with first-element unity. Under assumptions (19.1) and (19.2), if we have random draws $\left(\mathbf{x}_{i}, y_{i}\right)$ from the population, then OLS is consistent and $\sqrt{N}$ asymptotically normal for the parameters of interest, $\boldsymbol{\beta}$.

The problem we study in this chapter occurs when we observe only a censored version of $y$. In the top coding example, suppose that wealth is measured in\\
thousands of dollars and is top coded at $\$ 200,000$. Then we can define the censored version of wealth (for any unit that can be drawn from the population) as $w= \min (y, 200)$. Of course, for each random draw $i, w_{i}=\min \left(y_{i}, 200\right)$. If we are presented with the data set on $\mathbf{x}_{i}$ and $w_{i}$, where $w_{i}$ is called "wealth," we should notice that the maximum value of "wealth" in the sample is 200 , with a nontrivial fraction of observations at exactly 200. Because there is no behavorial reason to see a focal point for wealth at 200-let alone, to observe no values greater than 200-we would recognize that the wealth variable has been top coded at 200 .

\subsection*{19.2.1 Binary Censoring}
We first cover the case where the censoring of the underlying response variable is extreme. As an example, suppose we want to model willingness to pay (WTP) for a proposed public project. Assume that the underlying model is as in equations (19.1) and (19.2) with $y=w t p$. When we draw family (say) $i$ from the population, we would like to observe $\left(\mathbf{x}_{i}, w t p_{i}\right)$; if we did so for all $i$, we would estimate $\boldsymbol{\beta}$ by OLS. But willingess to pay can be difficult to elicit, and reported amounts might be noisy. Instead, suppose that each family is presented with a cost of the project, $r_{i}$. Presented with this cost, the household either says it is in favor of the project or not. Thus, along with $\mathbf{x}_{i}$ and $r_{i}$, we observe the binary response\\
$w_{i}=1\left[y_{i}>r_{i}\right]$,\\
where we assume, for now, that the chance that $y_{i}$ equals $r_{i}$ is zero.\\
What is the most natural way to proceed to estimate $\boldsymbol{\beta}$ ? If we impose some strong assumptions on the underlying population and the nature of $r_{i}$, then we can proceed with maximum likelihood. In particular, assume\\
$u_{i} \mid \mathbf{x}_{i}, r_{i} \sim \operatorname{Normal}\left(0, \sigma^{2}\right)$.

Assumption (19.4) implies that $y_{i}=\mathbf{x}_{i} \boldsymbol{\beta}+u_{i}$ actually satisfies the classical linear model (CLM). It also requires that $r_{i}$ is independent of $y_{i}$ conditional on $\mathbf{x}_{i}$, that is,\\
$\mathrm{D}\left(y_{i} \mid \mathbf{x}_{i}, r_{i}\right)=\mathrm{D}\left(y_{i} \mid \mathbf{x}_{i}\right)$.

Assumption (19.5) is satisfied if $r_{i}$ is randomized or if $r_{i}$ is chosen as a function of $\mathbf{x}_{i}$, or some combination of these alternatives.

Given assumption (19.4),


\begin{align*}
\mathrm{P}\left(w_{i}\right. & \left.=1 \mid \mathbf{x}_{i}, r_{i}\right)=\mathrm{P}\left(y_{i}>r_{i} \mid \mathbf{x}_{i}, r_{i}\right)=\mathrm{P}\left[u_{i} / \sigma>\left(r_{i}-\mathbf{x}_{i} \boldsymbol{\beta}\right) / \sigma \mid \mathbf{x}_{i}, r_{i}\right] \\
& =1-\Phi\left[\left(r_{i}-\mathbf{x}_{i} \boldsymbol{\beta}\right) / \sigma\right]=\Phi\left[\left(\mathbf{x}_{i} \boldsymbol{\beta}-r_{i}\right) / \sigma\right] . \tag{19.6}
\end{align*}


We can see that the binary response, which indicates whether unit $i$ is in favor of the project at cost $r_{i}$, follows a probit model with parameter vector $\boldsymbol{\beta} / \sigma$ on $\mathbf{x}_{i}$ and $-1 / \sigma$ on $r_{i}$. (In almost all applications $\mathbf{x}_{i}$ would include an intercept, and we allow that possibility here.) Therefore, all the parameters, including $\sigma$, are identified provided $\mathbf{x}_{i}$ is not perfectly collinear and $r_{i}$ varies across $i$ in a way not perfectly linearly related to $\mathbf{x}_{i}$. Given the data censoring, the maximum likelihood estimators (MLEs) are the asymptotically efficient estimators of $\boldsymbol{\beta}$ and $\sigma$ (or $\sigma^{2}$ ). Because the underlying population model is linear, we are interested in the slopes, $\beta_{j}$. In the next section we show that the binary censoring problem is a special case of interval censoring with unitspecific thresholds.

The costs of the binary censoring scheme are potentially severe. If we could observe $y_{i}$, specifying $\mathrm{E}\left(y_{i} \mid \mathbf{x}_{i}\right)=\mathbf{x}_{i} \boldsymbol{\beta}$ would suffice for consistent estimation of $\boldsymbol{\beta}$; in fact, we could just specify a linear projection and use OLS. With censoring, we must add more, and assumption (19.4) implies that the underlying model satisfies the CLM. It is in this setting where discussions of the deleterious effects of nonnormality and heteroskedasticity when using probit models make sense. In Chapter 15 we focused on the case where the binary response is the variable we want to explain, in which case we are interested in estimating the partial effects on the response probability. In that setting, heteroskedasticity and nonnormality in the error of the latent variable model change the functional form of the quantity of interest, and so the relevant issue concerns how those two problems affect the estimated partial effects (and the implications they have for a standard probit model, or even a linear probability model, for estimating the partial effects). With data censoring, we are interested only in the parameters in the underlying linear model. Consequently, it is now legitimate to be concerned about the effects on the parameter estimates of heteroskedasticity or nonnormality in the underlying linear model.

In Section 15.7.6 we discussed various ways of estimating parameters up to scale without placing strong restrictions on $\mathbf{D}\left(u_{i} \mid \mathbf{x}_{i}, r_{i}\right)$. Those estimators can be used in the present context. For example, if the distribution of $\left(\mathbf{x}_{i}, r_{i}, y_{i}\right)$ implies linear conditional expectations for all elements conditional on $y_{i}$, then the Chung and Goldberger (1984) results can be used for OLS. Ruud's $(1983,1986)$ findings can be applied when $u_{i}$ is independent of $\left(\mathbf{x}_{i}, r_{i}\right)$ if the distribution of $u_{i}$ is misspecified. Manski's $(1975,1988)$ maximum score estimator requires only symmetry of $\mathrm{D}\left(u_{i} \mid \mathbf{x}_{i}, r_{i}\right)$ (around zero), and Horowitz's (1992) smoothed version is more convenient for inference. But in every case these methods only estimate the slope coefficients up to a common scale factor (and the intercept cannot be estimated at all); therefore we cannot learn the magnitude of the effect of any element of $\mathbf{x}$ on willingness to pay, nor do we have a way of predicting willingness to pay for given values of the covariates.

A linear model for willingness to pay is not ideal because $w t p \geq 0$. If $w t p$ is zero for some subset of the population, a sensible population model is the Type I Tobit:\\
$y=w t p=\max (0, \mathbf{x} \boldsymbol{\beta}+u)$,\\
under assumption (19.4). If we had a random sample, we would use Type I Tobit MLE to estimate $\boldsymbol{\beta}$ and $\sigma^{2}$ (naturally, $r_{i}$ would not come into play), and we would use the MLEs to estimate the means, say, from a Type I Tobit. Interestingly, if we have binary censoring, the estimation procedure is identical to that outlined for a linear model for $w t p$, provided that the $r_{i}$ are all strictly positive. Because we do not observe $y_{i}$, we cannot distinguish between equations (19.1) and (19.7) when $r_{i}>0$. But if we believe that $y$ is zero for a nontrivial fraction of the population, any calculations should reflect that belief by using the Type I Tobit formulas for estimating partial effects.

One way to possibly determine whether WTP is ever zero in the population is to set some $r_{i}$ to zero so that the outcome $w_{i}=0$ means $w t p_{i}=0$. Or, we might set some $r_{i}$ sufficiently "close" to zero so that $w_{i}=0$ practically means $w t p_{i}=0$. Of course, this approach requires a particular survey design before the data have been collected.

If we really think $y=w t p>0$ in the population-say, everyone in a city would be willing to pay at least a small amount for a new park-then the underlying population model should be something like


\begin{equation*}
y=\exp (\mathbf{x} \boldsymbol{\beta}+u) \tag{19.8}
\end{equation*}


under assumption (19.4). Then, $\log (y)=\mathbf{x} \boldsymbol{\beta}+u$, so that the previous analysis applies but with $r_{i}$ replaced with $\log \left(r_{i}\right)$ (and we are back to assuming $r_{i}>0$ ). Interestingly, the change in the functional form for $\mathrm{P}\left(w_{i}=1 \mid \mathbf{x}_{i}, r_{i}\right)$ from equation (19.6) to


\begin{equation*}
\mathrm{P}\left(w_{i}=1 \mid \mathbf{x}_{i}, r_{i}\right)=\Phi\left[\left(\mathbf{x}_{i} \boldsymbol{\beta}-\log \left(r_{i}\right)\right) / \sigma\right] \tag{19.9}
\end{equation*}


provides a way to distinguish between assumptions (19.7) and (19.8) as the underlying population models when $r_{i}>0$ for all $i$.

Naturally, the previous models apply to cases other than willingness to pay. But if $y$ cannot take on negative values, the linear model we started with in equation (19.1) is not ideal. Assumption (19.7) leads to the same estimation method but implies a Tobit form for $\mathrm{E}(y \mid \mathbf{x})$. If we think $y>0$ in the population, then equation (19.8) is more attractive, and we can use equations (19.6) and (19.9) to distinguish between them based on values of the log-likelihood functions. In fact, because we know the coefficient on $r_{i}$ in the first case, and $\log \left(r_{i}\right)$ in the second, must be nonzero, we can use Vuoung's (1989) model selection test to choose between them; see Section 13.11.2.

\subsection*{19.2.2 Interval Coding}
We now consider the linear model (19.1) in a scenario where the continuous, quantitative outcome, $y$, is only recorded to fall into a particular interval. In this case we say we have interval-coded data (or interval-censored data). We are still interested in the population regression $\mathrm{E}(y \mid \mathbf{x})=\mathbf{x} \boldsymbol{\beta}$. Let $r_{1}<r_{2}<\cdots<r_{J}$ denote the known interval limits; these are specified as part of the survey design. For example, rather than asking individuals to report actual annual income, they report the interval that their income falls into.

Under the normality assumption in equation (19.4), we can estimate $\boldsymbol{\beta}$ and $\sigma^{2}$. Not surprisingly, the structure of the problem is similar to the ordered probit model we covered in Section 16.3. In fact, we can define


\begin{align*}
w=0 & \text { if } y \leq r_{1} \\
w=1 & \text { if } r_{1}<y \leq r_{2} \\
\vdots &  \tag{19.10}\\
w=J & \text { if } y>r_{J}
\end{align*}


and easily obtain the conditional probabilities $\mathrm{P}(w=j \mid \mathbf{x})$ for $j=0,1, \ldots, J$. The $\log$ likelihood for a random draw $i$ is


\begin{align*}
l_{i}(\boldsymbol{\beta}, \sigma)= & 1\left[w_{i}=0\right] \log \left\{\Phi\left[\left(r_{1}-\mathbf{x}_{i} \boldsymbol{\beta}\right) / \sigma\right]\right\}+1\left[w_{i}=1\right] \log \left\{\Phi \left[\left(r_{2}-\mathbf{x}_{i} \boldsymbol{\beta}\right) / \sigma\right.\right. \\
& \left.-\Phi\left[\left(r_{1}-\mathbf{x}_{i} \boldsymbol{\beta}\right) / \sigma\right]\right\}+\cdots+1\left[w_{i}=J\right] \log \left\{1-\Phi\left[\left(r_{J}-\mathbf{x}_{i} \boldsymbol{\beta}\right) / \sigma\right]\right\} \tag{19.11}
\end{align*}


The maximum likelihood estimators, $\hat{\boldsymbol{\beta}}$ and $\hat{\sigma}^{2}$, are often called interval regression estimators, with the understanding that the underlying population distribution is homoskedastic normal.

Although equation (19.11) looks a lot like the log likelihood for the ordered probit model, there is an important difference: in ordered probit, the cut points are parameters to estimate, and the parameters $\boldsymbol{\beta}$ do not measure interesting partial effects. With interval regression, the interval endpoints are given (or are themselves data), and $\boldsymbol{\beta}$ contains the partial effects of interest. In particular, as in the case of binary censoring, when we obtain the interval regression estimates, we interpret the $\hat{\boldsymbol{\beta}}$ as if we had been able to run the regression $y_{i}$ on $\mathbf{x}_{i}, i=1, \ldots, N$. Imposing the assumptions of the classical linear model allows us to estimate the parameters in the distribution $\mathrm{D}(y \mid \mathbf{x})$, even though the data are interval censored.

Sometimes in applications of interval regression the observed, censored variable, $w$, is set to some value within the interval that contains $y$. For example, if $y$ is wealth,\\
we might set $w$ to the midpoint of the interval that $y$ falls into. (Of course, we have to use some other rule if $y<r_{1}$ or $y>r_{J}$.) Provided the definition of $w$ determines the proper interval, the maximum likelihood estimators of $\boldsymbol{\beta}$ and $\sigma$ will be the same.

When $w$ is defined to have the same units as $y$, it is tempting to ignore the grouping of the data and just to run an OLS regression of $w_{i}$ on $\mathbf{x}_{i}, i=1, \ldots, N$. Naturally, such a procedure is generally inconsistent for $\boldsymbol{\beta}$. Nevertheless, the results of Chung and Goldberger apply: if $\mathrm{E}(\mathbf{x} \mid y)$ is linear in $y$, a linear regression can estimate the slope coefficients up to a common scale factor.

Sometimes the interval limits change across $i$, a possibility that causes no problems if we assume the limits are exogenous in the following sense:\\
$\mathrm{D}\left(y_{i} \mid \mathbf{x}_{i}, r_{i 1}, \ldots, r_{i J}\right)=\mathrm{D}\left(y_{i} \mid \mathbf{x}_{i}\right)$.

In the binary censoring example from the previous section, assumption (19.12) holds because the one limit value ( $J=1$ ) is randomly assigned. Generally, the limits can be a function of $\mathbf{x}_{i}$ (because these are being conditioned on). The resulting log likelihood is exactly as in equation (19.11) with $r_{j}$ replaced with $r_{i j}$. Some econometrics packages have "interval regression" commands that allow one to specify the lower and upper endpoints for each unit $i$, allowing for unit-specific endpoints.

Because of the underlying normality assumption, we can use the Rivers and Vuong (1988) and Smith and Blundell (1986) control function approach to test and correct for endogeneity of explanatory variables. The underlying model is the standard linear model\\
$y_{1}=\mathbf{z}_{1} \boldsymbol{\delta}_{1}+\alpha_{1} y_{2}+u_{1}$,\\
and we observe the censored variable, $w_{1}$. Given the linear reduced form $y_{2}= \mathbf{z} \boldsymbol{\delta}_{2}+v_{2}$, we proceed as before: just add the first-stage residuals, $\hat{v}_{2}$, to the interval regression model, along with $\left(\mathbf{z}_{1}, y_{2}\right)$. Of course, we are interested in $\alpha_{1}$ and $\boldsymbol{\delta}_{1}$, along with the coefficient on $\hat{v}_{2}$ to determine whether $y_{2}$ is in fact endogenous. Unfortunately, such an approach only works when $y_{2}$ is not censored. It is very difficult to account for interval censoring of $y_{2}$ along with that for $y_{1}$.

Modifying interval regression for linear, unobserved-effects panel data models is straightforward, provided we are willing to rely on the Chamberlain-Mundlak device. We would write


\begin{equation*}
y_{i t}=\mathbf{x}_{i t} \boldsymbol{\beta}+c_{i}+u_{i t}, \quad t=1, \ldots, T \tag{19.14}
\end{equation*}


$c_{i}=\psi+\overline{\mathbf{x}}_{i} \boldsymbol{\xi}+a_{i}$,\\
where all unobservables have normal distributions and the interval limits, $\left\{r_{i t j}: j=\right. 1,2, \ldots, J\}$, can vary by $i$ and $t$. Estimation can be carried out under the assumption\\
of serial independence in $\left\{u_{i t}: t=1, \ldots, T\right\}$, so that the log likelihood has a random effects structure, or without imposing any assumption on the serial dependence (which leads to pooled estimation of the type we covered in Chapters 15 and 16).

\subsection*{19.2.3 Censoring from Above and Below}
Two kinds of censoring are common: one is seen when the value of a variable is observed only when it is below a known cap, the other when it is above a known floor (or, in some cases, both). Consider first the case of right censoring (censoring from above). For each $i$, a censoring threshold, $r_{i}$, is observed. When we randomly draw a unit from the population, we observe the explanatory variables, $\mathbf{x}_{i}$. However, rather than observing the outcome on $y_{i}$, we effectively observe\\
$w_{i}=\min \left(y_{i}, r_{i}\right)$.

If the underlying population distribution for $y$ is continuous, the probability that $y_{i}=r_{i}$ is zero, and so, if $w_{i}=r_{i}$, we know that the observation is censored; if $w_{i}<r_{i}$, we know that $w_{i}=y_{i}$; that is, we observe $y_{i}$. If $y_{i}$ is (partially) discrete, there can be positive probability that $y_{i}=r_{i}$. As a practical matter, when the censoring points change across $i$, it is helpful in data sets to define a binary variable indicating whether an observation is censored.

For concreteness, we focus on the case where the population distribution of $y_{i}$ is continuous. Let $f(y \mid \mathbf{x} ; \boldsymbol{\theta})$ denote the conditional density. Under $\mathrm{D}\left(y_{i} \mid \mathbf{x}_{i}, r_{i}\right)= \mathrm{D}\left(y_{i} \mid \mathbf{x}_{i}\right)$, we can easily obtain the density of $w_{i}$ conditional on $\left(\mathbf{x}_{i}, r_{i}\right)$ because, for $w<r_{i}$,\\
$\mathrm{P}\left(w_{i} \leq w \mid \mathbf{x}_{i}, r_{i}\right)=\mathrm{P}\left(y_{i} \leq w \mid \mathbf{x}_{i}\right)=F\left(w \mid \mathbf{x}_{i} ; \boldsymbol{\theta}\right)$,\\
where $F\left(\cdot \mid \mathbf{x}_{i} ; \boldsymbol{\theta}\right)$ is the cdf of $y_{i}$ conditional on $\mathbf{x}_{i}$. Therefore, the probability density of $w_{i}$ given $\left(\mathbf{x}_{i}, r_{i}\right)$ is simply $f\left(w \mid \mathbf{x}_{i} ; \boldsymbol{\theta}\right)$ for $w<r_{i}$, that is, for values strictly less than the censoring point. Further,\\
$\mathrm{P}\left(w_{i}=r_{i} \mid \mathbf{x}_{i}, r_{i}\right)=\mathrm{P}\left(y_{i} \geq r_{i} \mid \mathbf{x}_{i}, r_{i}\right)=1-F\left(r_{i} \mid \mathbf{x}_{i} ; \boldsymbol{\theta}\right)$.\\
It follows that we can write the probability density of $w_{i}$ given $\left(\mathbf{x}_{i}, r_{i}\right)$ as\\
$g\left(w \mid \mathbf{x}_{i}, r_{i} ; \boldsymbol{\theta}\right)=\left[f\left(w \mid \mathbf{x}_{i} ; \boldsymbol{\theta}\right)\right]^{1\left[w<r_{i}\right]}\left[1-F\left(r_{i} \mid \mathbf{x}_{i} ; \boldsymbol{\theta}\right)\right]^{1\left[w=r_{i}\right]}$.

The log-likelihood function for a random draw $i$ (where we do not bother to distinguish between the "true" value of theta and a generic value) is\\
$\log \left[g\left(w \mid \mathbf{x}_{i}, r_{i} ; \boldsymbol{\theta}\right)\right]=1\left[w_{i}<r_{i}\right] \log \left[f\left(w_{i} \mid \mathbf{x}_{i} ; \boldsymbol{\theta}\right)\right]+1\left[w_{i}=r_{i}\right] \log \left[1-F\left(r_{i} \mid \mathbf{x}_{i} ; \boldsymbol{\theta}\right)\right]$,\\
and we sum this expression across all $i$ to obtain the log likelihood for the entire sample. Some authors prefer to write the log likelihood separately for the uncensored and censored observations, but equation (19.18) is actually preferred because it shows that the units are being randomly sampled from a population. Here, we are drawing $\left(w_{i}, \mathbf{x}_{i}, r_{i}\right)$. In the vast majority of cases, the conditions sufficient for MLE to be well behaved covered in Chapter 13 hold for censored estimation because the model $f(y \mid \mathbf{x} ; \boldsymbol{\theta})$ is smooth in $\boldsymbol{\theta}$.

An interesting feature of equation (19.18) is that we only need to observe the censoring point, $r_{i}$, for censored observations. (However, we do need to know which observations are censored and which are not.) This feature of the MLE approach is useful because sometimes in applications to duration models-see Chapter 22-the censoring value is reported only for observations that are actually censored.

In the leading case, $y$ follows a classical linear model in the population of interest, that is,\\
$\mathrm{D}(y \mid \mathbf{x})=\operatorname{Normal}\left(\mathbf{x} \boldsymbol{\beta}, \sigma^{2}\right)$,\\
in which case we have what is typically called the censored normal regression model. (In the econometrics literature, this is sometimes called the censored Tobit model or Type I Tobit model, but we are reserving those names for the case of corner solution responses; see Chapter 17.) The log likelihood for the censored normal regression model is\\
$l_{i}(\boldsymbol{\theta})=1\left[w_{i}<r_{i}\right] \log \left\{\sigma^{-1} \phi\left[\left(w_{i}-\mathbf{x}_{i} \boldsymbol{\beta}\right) / \sigma\right]\right\}+1\left[w_{i}=r_{i}\right] \log \left\{1-\Phi\left[\left(w_{i}-\mathbf{x}_{i} \boldsymbol{\beta}\right) / \sigma\right]\right\}$.

Many standard econometrics packages estimate this model with little computational difficulty. Often, a transformation, such as taking the natural log, is needed to make equation (19.19) a reasonable assumption. In cases where no such transformation is available that ensures normality, the more general formulation in equation (19.18) can be used to obtain the log likelihood.

Distinguishing between the underlying population model and the censoring scheme can lead to some perhaps surprising implications for econometric practice. For example, suppose that survey data are collected on charitable contributions, where contributions are censored at a fixed cap—let us say $\$ 10,000$ a year for concreteness. In the population, there is no natural upper bound for charitable contributions, so one knows immediately that the observed pileup at $\$ 10,000$ in the survey data is due to built-in data censoring. It is proper to say that charitable contributions are "censored from above at $\$ 10,000$." But the pileup at zero is a different matter: it is due to the fact that some fraction of the population will have zero charitable contributions\\
in a particular year. Therefore, an appropriate course of action is to treat charitable contributions in the population as a corner solution response, with a corner at zero. (We might use a Tobit model, or one of the two-part models discussed in Section 17.6.) Unlike the censoring from above at $\$ 10,000$, it makes no sense to say charitable contributions are also "censored from below at zero." There is a corner at zero, but it is not due to data censoring.

Interestingly, in the situation just described, if we assume a Type I Tobit model in the population for charitable contributions, the log-likelihood function for the rightcensored charitable contributions is identical to that for the two-limit Tobit model discussed in Section 17.7. Estimation can be done rather straightforwardly by specifying zero and 10,000 as the lower and upper bounds, respectively. The practically important issue concerns what we do with the estimates. In fact, after obtaining the estimates, all calculations-of response probabilities, expected values, and so onshould be based on the Type I Tobit model (ignoring the right censoring). We are interested in features of $\mathrm{D}(y \mid \mathbf{x})$ in the population, where $y$ is actual charitable contributions. Therefore, after accounting for the top coding in estimation by using a two-limit Tobit, we revert to the Type I Tobit for all statements about charitable contributions. The parameters themselves are of interest insofar as they allow us to compute partial effects on probabilities, means, and medians. But the corner at $\$ 10,000$ plays no role in such calculations. Similar comments hold for two-part models for corner solutions.

A more subtle example occurs when, say, an upper limit is imposed by law. Consider a stylized case where individuals may contribute no more than $15 \%$ of their income to retirement plans. In the population of working people, some individuals will contribute zero, some will contribute at the 15 percent upper limit, and many will contribute a percentage strictly between zero and 15. If we are interested in the effect of explanatory variables-say, taking courses on retirement savings-on the contribution under the current legal regime, a two-limit Tobit model makes sense for the contribution percentage. If we want to know, say, the mean or median difference in the rate between those who have and have not had a retirement savings course, we would use the formulas for the two-limit Tobit model. However, one might want to know the effects of covariates on the contribution percentage in the absence of institutional constraints. Then we would be back to the previous situation: the corner at zero is a corner that arises from utility maximization, but the corner at 15 is imposed-in this case, by law rather than the data collection scheme. In this case, one would use formulas for a Type I Tobit with corner at zero, even though the estimates necessarily are obtained from a two-limit Tobit. In this application, an additional calculation may be of interest: what would be the effect on the average\\
contribution rate if the limit were increased? We can use the formulas for a two-limit Tobit in this case to compute the appropriate derivative.

Not surprisingly, endogenous explanatory variables with censored data can be handled using methods very similar to those for the Type I Tobit model. For example, suppose the population model is\\
$y_{1}=\mathbf{z}_{1} \boldsymbol{\delta}_{1}+\alpha_{1} y_{2}+u_{1}$,\\
where $\mathbf{D}\left(u_{1} \mid \mathbf{z}\right)$ is $\operatorname{Normal}\left(0, \sigma_{1}^{2}\right)$. However, the data are right censored. If the reduced form for $y_{2}$ is $y_{2}=\mathbf{z} \boldsymbol{\delta}_{1}+v_{2}$, where ( $u_{1}, v_{2}$ ) is independent of $\mathbf{z}$ and bivariate normal, we can apply the Smith and Blundell (1986) approach to account for the right censoring. This assumes that $y_{2}$ is not censored, so that the first step is OLS of $y_{2}$ on $\mathbf{z}$ using a random sample. The residuals, $\hat{v}_{2}$, are added to the censored normal regression in the second stage. Of course, because the underlying population model is linear, we are interested in $\alpha_{1}$ and $\boldsymbol{\delta}_{1}$. Joint MLE is possible, too, and would be more efficient and avoid the problem of inference after two-step estimation. (Of course, the bootstrap can be applied here by randomly drawing units from the sample. Remember, we have a random sample of units. We simply include the first-step estimation and censored normal estimation within each bootstrap iteration.)

With enough normality, it should not be surprising that the Chamberlain-Mundlak device can be used in the context of right and left censoring. Again, for simplicity consider the linear model $y_{i t}=\mathbf{x}_{i t} \boldsymbol{\beta}+c_{i}+u_{i t}$, where right censoring is at $r_{i t}$, which can change across $i$ and $t$. Natural exogeneity assumptions, along with convenient distributional assumptions, are


\begin{equation*}
\mathrm{D}\left(u_{i t} \mid \mathbf{x}_{i}, r_{i 1}, \ldots, r_{i T}, c_{i}\right)=\mathrm{D}\left(u_{i t}\right)=\operatorname{Normal}\left(0, \sigma_{u}^{2}\right), \quad t=1, \ldots, T \tag{19.22}
\end{equation*}


$\mathrm{D}\left(c_{i} \mid \mathbf{x}_{i}, r_{i 1}, \ldots, r_{i T}\right)=\mathrm{D}\left(c_{i} \mid \mathbf{x}_{i}\right)=\operatorname{Normal}\left(\psi+\overline{\mathbf{x}}_{i} \boldsymbol{\xi}, \sigma_{a}^{2}\right)$.

As usual, these assumptions mean we can write\\
$y_{i t}=\psi+\mathbf{x}_{i t} \boldsymbol{\beta}+\overline{\mathbf{x}}_{i} \boldsymbol{\xi}+v_{i t}, \quad \mathrm{D}\left(v_{i t} \mid \mathbf{x}_{i}, \mathbf{r}_{i}\right)=\operatorname{Normal}\left(0, \sigma_{a}^{2}+\sigma_{u}^{2}\right)$,\\
where $\mathbf{r}_{i}=\left(r_{i 1}, \ldots, r_{i T}\right)$ is the vector of censoring values for unit $i$. Now, we can just apply pooled censored normal regression, with censoring points $r_{i t}$, and consistently estimate $\psi, \boldsymbol{\beta}, \boldsymbol{\xi}$, and $\sigma_{v}^{2}=\sigma_{a}^{2}+\sigma_{u}^{2}$. Because this is a partial likelihood method, we need to make the inference robust to serial correlation. Generally, we cannot separately identify $\sigma_{a}^{2}$ and $\sigma_{u}^{2}$ unless we make the further assumption that $\left\{u_{i t}: t=1, \ldots\right.$, $T\}$ is serially independent, in which case we can use a correlated random effects (CRE) likelihood approach similar in structure to the CRE Tobit model. Naturally,\\
with the underlying population model linear, we are mainly interested in $\boldsymbol{\beta}$ and appropriate inference concerning $\boldsymbol{\beta}$.

All the methods just discussed require a full distributional assumption. (In the panel data case, this statement applies to $\mathrm{D}\left(y_{i t} \mid \mathbf{x}_{i}, c_{i}\right)$, not to a joint distribution across $t$.) Certainly it is of interest to relax such assumptions when possible. (In Section 19.8, we will discuss ways of using probability weights to relax distributional assumptions.) Powell's (1984) censored least absolute deviations (CLAD) estimator can be applied to censored data without putting strong restrictions on $\mathrm{D}(y \mid \mathbf{x})$. To see how, begin with a linear model for the conditional median,\\
$\operatorname{Med}(y \mid \mathbf{x})=\mathbf{x} \boldsymbol{\beta}$.\\
(Of course, this may or may not be the conditional mean. Powell's approach applies to the conditional median.) Again, assume that our random sample consists of $\left(\mathbf{x}_{i}, r_{i}, w_{i}\right)$, where $w_{i}=\min \left(y_{i}, r_{i}\right)$. LAD can be applied to the censoring case because the median passes through the min function:\\
$\operatorname{Med}\left(w_{i} \mid \mathbf{x}_{i}, r_{i}\right)=\operatorname{Med}\left[\min \left(y_{i} \mid \mathbf{x}_{i}, r_{i}\right), r_{i}\right]=\min \left[\operatorname{Med}\left(y_{i} \mid \mathbf{x}_{i}\right), r_{i}\right]=\min \left(\mathbf{x}_{i} \boldsymbol{\beta}, r_{i}\right)$,\\
where the second equality holds under the assumption that $\operatorname{Med}\left(y_{i} \mid \mathbf{x}_{i}, r_{i}\right)= \operatorname{Med}\left(y_{i} \mid \mathbf{x}_{i}\right)$-which means censoring is exogenous with respect to $y_{i}$ in the conditional median sense. As before, the censoring values can be related to $\mathbf{x}_{i}$. Given equation (19.26), we can use LAD to estimate $\boldsymbol{\beta}$, resulting in the CLAD estimator:\\
$\min _{\mathbf{b}} \sum_{i=1}^{N}\left|w_{i}-\min \left(\mathbf{x}_{i} \mathbf{b}, r_{i}\right)\right|$.

As discussed in Chapter 17, Powell (1984) shows that CLAD is consistent and $\sqrt{N}$ asymptotically normal. His paper contains formulas for estimating the asymptotic variance of the CLAD estimator.

One subtle point concerning the CLAD estimator when applied to censored data is that it requires that the censoring value, $r_{i}$, be available even when the observation is not right censored. This data requirement is not much of an issue in top-coding cases, especially when the same value is used. (For example, if wealth is top coded at $\$ 500,000$, that information is known, and $r_{i}=500,000$ for all $i$.) However, in some duration problems-which we treat explicitly in Chapter 22-only $w_{i}$ is observed. That is, along with a censoring indicator, we observe either $y_{i}$ or $r_{i}$, but not both. Recall that the maximum likelihood estimator can be applied in situations where $r_{i}$ is\\
not always observed: we only need to observe $r_{i}$ when the observation is actually censored.

\subsection*{19.3 Overview of Sample Selection}
We now turn to estimation when a sample from subset of the population is used to estimate the unknown parameters. The term selected sample is generally used to describe a sample that is not randomly drawn from the underlying population. As mentioned in Section 19.1, there are a variety of selection mechanisms that result in selected samples. Some mechanisms are due to sample design, while others are due to the behavior of the units being sampled, including nonresponse on survey questions and attrition from social programs.

Before we launch into specifics, there is an important general point to remember: sample selection can only be an issue once the population of interest has been carefully specified. If we propose a model for a subset of a larger population, it is proper to proceed by obtaining a random sample from that subpopulation and then using the standard econometric methods that we have covered thus far. That we do not have a random sample from the larger population does not affect our ability to consistently estimate the parameters of the model for the subpopulation.

As a specific example that often leads to confusion, consider the conditional lognormal hurdle model that we discussed in Section 17.6.2. In that model, the distribution of $\log (y)$, conditional on $y>0$ and the covariates $\mathbf{x}$, is normally distributed with a linear conditional mean and constant variance. Therefore, the parameters in the model describing the $y>0$ subpopulation, $\boldsymbol{\beta}$ and $\sigma^{2}$, are consistently estimated using MLE (which is OLS in this case) on the subsample with $y_{i}>0$. Using a linear model for $\log (y)$ on the subsample $y_{i}>0$ does not work in the Exponential Type II Tobit model because that model implies that $\log (y)$ conditional on $y>0$ is not $\log$ normal, nor does it have a linear conditional mean. Some authors prefer to view the failure of the standard $\log (y)$ regression in the ET2T model as a sample selection problem, but that labeling misses the key point: in hurdle models, we are presumably interested in features of $\mathbf{D}(y \mid \mathbf{x})$, and the only issue is whether we have those features correctly specified. In the ET2T model, one often sees discussion of "sample selection bias" in estimating the vector $\boldsymbol{\beta}$ in the formulation $y=1[\mathbf{x} \boldsymbol{\gamma}+v>0] \exp (\mathbf{x} \boldsymbol{\beta}+u)$. But $\boldsymbol{\beta}$ does not by itself provide partial effects of any conditional mean involving $y$, and so, as we discussed in Section 17.6.3, focusing on estimates of $\boldsymbol{\beta}$ in the ET2T model is inappropriate. (In the lognormal hurdle model, $\boldsymbol{\beta}$ indexes the semielasticities and elasticities of $\mathrm{E}(y \mid \mathbf{x}, y>0)$, which makes the $\beta_{j}$ of direct interest.) By contrast,\\
in a sample selection context with a linear regression for the underlying population, the focus on the single set of regression parameters is entirely appropriate, as we wll see in the next several sections.

As a second example, suppose $y$ is a fractional response that takes on the values zero and one with positive probability, and takes on a range of values strictly between zero and one. One possibility is to model $\mathrm{P}(y=0 \mid \mathbf{x})$ and $\mathrm{P}(y=1 \mid \mathbf{x})$ along with $\mathrm{E}(y \mid \mathbf{x}, 0<y<1)$. Suppose the latter is $\mathrm{E}(y \mid \mathbf{x}, 0<y<1)=\Phi(\mathbf{x} \boldsymbol{\beta})$. Then, to consistently estimate $\boldsymbol{\beta}$ we can apply any of the consistent estimators for fractional responses in Section 18.6.1 to the subsample with $0<y_{i}<1$. We do not introduce a sample selection problem by ignoring the data with outcomes $y_{i}=0$ or $y_{i}=1$.

Now that we know some contexts where sample selection is not an issue, we provide some examples where nonrandom sampling is relevant and can (but does not always) cause serious problems.

Example 19.1 (Saving Function): Suppose we wish to estimate a saving function for all families in a given country, and the population saving function is\\
saving $=\beta_{0}+\beta_{1}$ income $+\beta_{2}$ age $+\beta_{3}$ married $+\beta_{4}$ kids $+u$,\\
where age is the age of the household head and the other variables are self-explanatory. However, we only have access to a survey that included families whose household head was 45 years of age or older. This restricted sampling raises a sample selection issue because we are interested in the saving function for all families, but we can obtain a random sample only for a subset of the population.

Example 19.2 (Truncation Based on Wealth): We are interested in estimating the effect of worker eligibility in a particular pension plan (for example, a 401(k) plan) on family wealth. Let the population model be\\
wealth $=\beta_{0}+\beta_{1}$ plan $+\beta_{2}$ educ $+\beta_{3}$ age $+\beta_{4}$ income $+u$,\\
where plan is a binary indicator for eligibility in the pension plan. However, we can only sample people with a net wealth less than $\$ 200,000$, so the sample is selected on the basis of wealth. As we will see, sampling based on a response variable is much more serious than sampling based on an exogenous explanatory variable.

In these two examples data were missing on all variables for a subset of the population as a result of survey design. In other cases, units are randomly drawn from the population, but data are missing on one or more variables for some units in the sample. Using a subset of a random sample because of missing data can lead to a\\
sample selection problem. As we will see, if the reason the observations are missing is appropriately exogenous, using the subsample has no serious consequences.

Our final example illustrates a more subtle form of a missing data problem.\\
Example 19.3 (Wage Offer Function): Consider estimating a wage offer equation for people of working age. By definition, this equation is supposed to represent all people of working age, whether or not a person is actually working at the time of the survey. Because we can only observe the wage offer for working people, we effectively select our sample on this basis.

This example is not as straightforward as the previous two. We treat it as a sample selection problem because data on a key variable-the wage offer, wage - are available only for a clearly defined subset of the population. This is sometimes called incidental truncation because wage ${ }^{\mathrm{o}}$ is missing as a result of the outcome of another variable, labor force participation.

The incidental truncation in this example has a strong self-selection component: people self-select into employment, so whether or not we observe wage ${ }^{\mathrm{o}}$ depends on an individual's labor supply decision. Whether we call examples like this sample selection or self-selection is largely irrelevant. The important point is that we must account for the nonrandom nature of the sample we have for estimating the wage offer equation.

In the next several sections we cover a variety of sample selection issues, including tests and corrections.

\subsection*{19.4 When Can Sample Selection Be Ignored?}
In Section 19.3 we briefly discussed how there can be no sample selection problem if the model we have specified applies to a subsample of a population for which we can obtain a random sample. Here, we discuss a more substantive question: under what circumstances does using a nonrandom sample from a specified population nevertheless consistently estimate the population parameters?

\subsection*{19.4.1 Linear Models: Estimation by OLS and 2SLS}
We begin by obtaining conditions under which estimation of the population model by two-stage least squares (2SLS) using a selected sample is consistent for the population parameters. These results are of interest in their own right, but we will also apply them to several situations later in the chapter.

We assume there is a population represented by the random vector $(\mathbf{x}, y, \mathbf{z})$, where $\mathbf{x}$ is a $1 \times K$ vector of explanatory variables, $y$ is the scalar response variable, and $\mathbf{z}$ is a $1 \times L$ vector of instrumental variables. The population model is the standard single-equation linear model with possibly endogenous explanatory variables:


\begin{equation*}
y=\beta_{1}+\beta_{2} x_{2}+\cdots+\beta_{K} x_{K}+u=\mathbf{x} \boldsymbol{\beta}+u \tag{19.30}
\end{equation*}


$\mathrm{E}\left(\mathbf{z}^{\prime} u\right)=\mathbf{0}$,\\
where we take $x_{1} \equiv 1$ for notational simplicity (an assumption that means $z_{1}$ is almost certainly equal to unity, too). From Chapter 5 we know that, if we could obtain a random sample from the population, then equation (19.31), along with the rank condition (particularly $\operatorname{rank}\left[\mathrm{E}\left(\mathbf{z}^{\prime} \mathbf{x}\right)\right]=K$ ), would be sufficient to consistently estimate $\boldsymbol{\beta}$. As we will see, in the context of sample selection, equation (19.31) is rarely sufficient for consistency of the 2SLS estimator on the selected sample.

A leading special case is $\mathbf{z}=\mathbf{x}$, in which case the explanatory variables are assumed to be uncorrelated with the error. Our general treatment allows elements of $\mathbf{x}$ to be correlated with $u$.

Rather than obtaining a random sample from the population, we only use data points that satisfy certain conditions. The idea is to think of drawing units randomly from the population, but now a random draw for unit $i,\left(\mathbf{x}_{i}, y_{i}, \mathbf{z}_{i}\right)$, is supplemented by drawing a selection indicator, $s_{i}$. By definition, $s_{i}=1$ if unit $i$ is used in the estimation, and $s_{i}=0$ if we do not use random draw $i$. Therefore, our "data" consist of $\left\{\left(\mathbf{x}_{i}, y_{i}, \mathbf{z}_{i}, s_{i}\right): i=1, \ldots, N\right\}$, where the value of $s_{i}$ determines whether we observe all of $\left(\mathbf{x}_{i}, y_{i}, \mathbf{z}_{i}\right)$.

Because parameter identification should always be studied in a population, we let $s$ denote a random variable with the distribution of $s_{i}$ for all $i$. In other words, $(\mathbf{x}, y, \mathbf{z}, s)$ now represents the population. Therefore, to determine the properties of any estimation procedure using the selected sample, we need to know about the distribution of $s$ and its dependence on $(\mathbf{x}, y, \mathbf{z})$.

To obtain conditions under which 2SLS on the selected sample consistently estimates $\boldsymbol{\beta}$, assume $\left\{\left(\mathbf{x}_{i}, y_{i}, \mathbf{z}_{i}, s_{i}\right): i=1, \ldots, N\right\}$ is a random sample from the population. The 2SLS estimator using the selected sample can be written as

$$
\begin{aligned}
\hat{\boldsymbol{\beta}}= & {\left[\left(N^{-1} \sum_{i=1}^{N} s_{i} \mathbf{z}_{i}^{\prime} \mathbf{x}_{i}\right)^{\prime}\left(N^{-1} \sum_{i=1}^{N} s_{i} \mathbf{z}_{i}^{\prime} \mathbf{z}_{i}\right)^{-1}\left(N^{-1} \sum_{i=1}^{N} s_{i} \mathbf{z}_{i}^{\prime} \mathbf{x}_{i}\right)\right]^{-1} } \\
& \times\left(N^{-1} \sum_{i=1}^{N} s_{i} \mathbf{z}_{i}^{\prime} \mathbf{x}_{i}\right)^{\prime}\left(N^{-1} \sum_{i=1}^{N} s_{i} \mathbf{z}_{i}^{\prime} \mathbf{z}_{i}\right)^{-1}\left(N^{-1} \sum_{i=1}^{N} s_{i} \mathbf{z}_{i}^{\prime} y_{i}\right)
\end{aligned}
$$

Substituting $y_{i}=\mathbf{x}_{i} \boldsymbol{\beta}+u_{i}$ gives


\begin{align*}
\hat{\boldsymbol{\beta}}= & \boldsymbol{\beta}+\left[\left(N^{-1} \sum_{i=1}^{N} s_{i} \mathbf{z}_{i}^{\prime} \mathbf{x}_{i}\right)^{\prime}\left(N^{-1} \sum_{i=1}^{N} s_{i} \mathbf{z}_{i}^{\prime} \mathbf{z}_{i}\right)^{-1}\left(N^{-1} \sum_{i=1}^{N} s_{i} \mathbf{z}_{i}^{\prime} \mathbf{x}_{i}\right)\right]^{-1} \\
& \times\left(N^{-1} \sum_{i=1}^{N} s_{i} \mathbf{z}_{i}^{\prime} \mathbf{x}_{i}\right)^{\prime}\left(N^{-1} \sum_{i=1}^{N} s_{i} \mathbf{z}_{i}^{\prime} \mathbf{z}_{i}\right)^{-1}\left(N^{-1} \sum_{i=1}^{N} s_{i} \mathbf{z}_{i}^{\prime} u_{i}\right) . \tag{19.32}
\end{align*}


It is easily seen from equation (19.32) that the key condition for consistency is $\mathrm{E}\left(s_{i} \mathbf{z}_{i}^{\prime} u_{i}\right)=\mathbf{0}$, along with the rank condition on the selected sample. Formally, we have the following result:\\
theorem 19.1 (Consistency of 2SLS under Sample Selection): In model (19.30), assume that $\mathrm{E}\left(u^{2}\right)<\infty, \mathrm{E}\left(x_{j}^{2}\right)<\infty, j=1, \ldots, K$, and $\mathrm{E}\left(z_{j}^{2}\right)<\infty, j=1, \ldots, L$. Further, assume that


\begin{align*}
& \mathrm{E}\left(s \mathbf{z}^{\prime} u\right)=\mathbf{0}  \tag{19.33}\\
& \operatorname{rank} \mathrm{E}\left(\mathbf{z}^{\prime} \mathbf{z} \mid s=1\right)=L  \tag{19.34}\\
& \operatorname{rank} \mathrm{E}\left(\mathbf{z}^{\prime} \mathbf{x} \mid s=1\right)=K \tag{19.35}
\end{align*}


Then the 2SLS estimator using the selected sample is consistent for $\boldsymbol{\beta}$ and $\sqrt{N}$ asymptotically normal.

Equation (19.32) essentially proves the consistency result under the assumptions of Theorem 19.1. Conditions (19.34) and (19.35) are fairly straightforward, and imply that the usual rank condition for 2SLS holds in the selected subpopulation. Naturally, it is possible that the rank condition holds on the entire population but not the $s=1$ subset. For example, if $s=1$ denotes females in the working population, then neither $\mathbf{x}$ nor $\mathbf{z}$ can include a female dummy variable. Generally, one needs sufficient variation in both the explanatory variables and instruments in the subpopulation in order for the rank condition to hold.

It is worth studying condition (19.33) in some detail. First, it is not generally enough to assume condition (19.31), that is, zero correlation between the instruments and errors in the population. One case where condition (19.31) is sufficient occurs when $s$ is independent of $(\mathbf{z}, u)$, so that $\mathrm{E}\left(s \mathbf{z}^{\prime} u\right)=\mathrm{E}(s) \mathrm{E}\left(\mathbf{z}^{\prime} u\right)=\mathbf{0}$ if $\mathrm{E}\left(\mathbf{z}^{\prime} u\right)=\mathbf{0}$. Of course, independence of selection and $(\mathbf{z}, u)$ is a very strong assumption. In the leading case where $\mathbf{z}=\mathbf{x}$, so that estimation is by OLS, the condition is the same as independence between $s$ and $(\mathbf{x}, y)$, which is tantamount to assuming that some\\
observations from an original random sample are dropped randomly, without regard to the values of $\left(\mathbf{x}_{i}, y_{i}\right)$. In the statistics literature, this has been called the missing completely at random (MCAR) assumption; see, for example, Little and Rubin (2002).

More interesting are situations where selection can depend on exogenous variables, but not on the unobserved error. An important sufficient condition for assumption (19.33), easily verified by applying iterated expectations, is\\
$\mathrm{E}(u \mid \mathbf{z}, s)=0$.

Assumption (19.36) allows selection to be correlated with $\mathbf{z}$ but not with $u$, and has been called exogenous sampling. It is easier to interpret this label by looking at a special case, which strengthens the sense in which the instruments are exogenous in the population. Rather than equation (19.31), make the population zero-conditionalmean assumption\\
$\mathrm{E}(u \mid \mathbf{z})=0$.

By basic properties of conditional expectations, if assumption (19.37) holds and $s$ is a deterministic function of $\mathbf{z}$, then assumption (19.36) holds (which means, of course, that assumption (19.33) holds). In other words, exogenous sampling occurs when $s=h(\mathbf{z})$ for some nonrandom function $h(\cdot)$; that is, $s$ is a nonrandom function of exogenous variables. But we must remember that the sense in which the instruments are exogenous is given by the stronger assumption (19.37).

In the case with exogenous explanatory variables, condition (19.36) is equivalent to\\
$\mathrm{E}(y \mid \mathbf{x}, s)=\mathrm{E}(y \mid \mathbf{x})=\mathbf{x} \boldsymbol{\beta}$,\\
where the first equality is how we would define exogenous sampling in the context of regression analysis. Notice that assumption (19.38) implies consistency of OLS on the selected sample regardless of whether data are missing on $y$ or elements of $\mathbf{x}$, or both. Assumption (19.38) rules out selection mechanisms that depend on the unobserved factors affecting $y$ in equation (19.30). This assumption is related to (a conditional mean version of) the missing at random (MAR) assumption (for example, Little and Rubin, 2002). MAR, which is often stated in terms of conditional distributions, further assumes that the variables determining selection ( $\mathbf{x}$ in this case) are always observed. We will have more to say on this kind of assumption in Section 19.8 on inverse probability weighting.

A sufficient condition for assumption (19.36) is that, conditional on $\mathbf{z}, u$ and $s$ are independent, which is informatively written as $\mathrm{D}(s \mid \mathbf{z}, u)=\mathrm{D}(s \mid \mathbf{z})$, or, because $s$ is binary,\\
$\mathrm{P}(s=1 \mid \mathbf{z}, u)=\mathrm{P}(s=1 \mid \mathbf{z})$.

Because $\mathbf{z}$ is observable-at least for part of the population-and $u$ is always unobservable, condition (19.39) is an example of what has been dubbed selection on observables, although this concept is usually employed in settings where $\mathbf{z}$ is always observed, something not required for the statement of Theorem 19.1. Again, we will have more to say on this kind of assumption in Section 19.8, where sample selection can also be based on variables that appear outside the model specification.

The asymptotic normality also follows in a fairly straightforward manner from equation (19.32); remember, the summands are i.i.d. random vectors, and the last term has a zero mean by assumption (19.33). Not surprisingly, the usual heteroskedasticity-robust variance matrix estimator, applied to the selected sample, is valid without further assumptions.

If we add the homoskedasticity assumption $\mathrm{E}\left(u^{2} \mid \mathbf{z}, s\right)=\mathrm{E}\left(u^{2}\right)=\sigma^{2}$ (which is the same as $\operatorname{Var}(u \mid \mathbf{z}, s)=\sigma^{2}$ if we maintain assumption (19.36)) to the assumptions of Theorem 19.1, then we can show that the "usual" variance matrix estimator for 2SLS is asymptotically valid. Doing so requires two steps. First, under $\mathrm{E}\left(u^{2} \mid \mathbf{z}, s\right)=\sigma^{2}$ the usual iterated expectations argument gives $\mathrm{E}\left(s u^{2} \mathbf{z}^{\prime} \mathbf{z}\right)=\sigma^{2} \mathrm{E}\left(s \mathbf{z}^{\prime} \mathbf{z}\right)$. This equation can be used to show that $\mathrm{Avar} \sqrt{N}(\hat{\boldsymbol{\beta}}-\boldsymbol{\beta})=\sigma^{2}\left\{\mathrm{E}\left(s \mathbf{x}^{\prime} \mathbf{z}\right)\left[\mathrm{E}\left(s \mathbf{z}^{\prime} \mathbf{z}\right)\right]^{-1} \mathrm{E}\left(s \mathbf{z}^{\prime} \mathbf{x}\right)\right\}^{-1}$. The second step is to show that the usual 2SLS estimator of $\sigma^{2}$ is consistent. This fact can be seen as follows. Under the homoskedasticity assumption, $\mathrm{E}\left(s u^{2}\right)=\mathrm{E}(s) \sigma^{2}$, where $\mathrm{E}(s)$ is just the fraction of the subpopulation in the overall population. The estimator of $\sigma^{2}$ (without degrees-of-freedom adjustment) is\\
$\left(\sum_{i=1}^{N} s_{i}\right)^{-1} \sum_{i=1}^{N} s_{i} \hat{u}_{i}^{2}$,\\
since $\sum_{i=1}^{N} s_{i}$ is simply the number of observations in the selected sample. Removing the " \^{} " from $u_{i}^{2}$ and applying the law of large numbers gives $N^{-1} \sum_{i=1}^{N} s_{i} \xrightarrow{p} \mathrm{E}(s)$ and $N^{-1} \sum_{i=1}^{N} s_{i} u_{i}^{2} \xrightarrow{p} \mathrm{E}\left(s u^{2}\right)=\mathrm{E}(s) \sigma^{2}$. Since the $N^{-1}$ terms cancel, expression (19.40) converges in probability to $\sigma^{2}$.

If $s$ is a function only of $\mathbf{z}$, or $s$ is independent of $(\mathbf{z}, u)$, and $\mathrm{E}\left(u^{2} \mid \mathbf{z}\right)=\sigma^{2}-$ that is, if the homoskedasticity assumption holds in the original population-then $\mathrm{E}\left(u^{2} \mid \mathbf{z}, s\right)=\sigma^{2}$. As mentioned previously, without the homoskedasticity assumption we would just use the heteroskedasticity-robust standard errors, just as if a random sample were available with heteroskedasticity present in the population model.

When $\mathbf{x}$ is exogenous and we apply OLS on the selected sample, Theorem 19.1 implies that we can select the sample on the basis of the explanatory variables.

Selection based on $y$ or on endogenous elements of $\mathbf{x}$ is not allowed because then $\mathrm{E}(u \mid \mathbf{x}, s) \neq \mathrm{E}(u)$.

Example 19.4 (Nonrandomly Missing IQ Scores): As an example of how Theorem 19.1 can be applied, consider the analysis in Griliches, Hall, and Hausman (1978) $(\mathrm{GHH})$. The structural equation of interest is\\
$\log ($ wage $)=\mathbf{z}_{1} \boldsymbol{\delta}_{1}+$ abil $+v, \quad \mathrm{E}\left(v \mid \mathbf{z}_{1}\right.$, abil,$\left.I Q\right)=0$,\\
and we assume that $I Q$ is a valid proxy for abil in the sense that abil $=\theta_{1} I Q+e$ and $\mathrm{E}\left(e \mid \mathbf{z}_{1}, I Q\right)=0$ (see Section 4.3.2). Write\\
$\log ($ wage $)=\mathbf{z}_{1} \boldsymbol{\delta}_{1}+\theta_{1} I Q+u$,\\
where $u=v+e$. Under the assumptions made, $\mathrm{E}\left(u \mid \mathbf{z}_{1}, I Q\right)=0$. It follows immediately from Theorem 19.1 that, if we choose the sample excluding all people with IQs below a fixed value, then OLS estimation of equation (19.41) will be consistent. This problem is not quite the one faced by GHH. Instead, GHH noticed that the probability of IQ missing was higher at lower IQs (because people were reluctant to give permission to obtain IQ scores). A simple way to model this situation is $s=1$ if $I Q+r \geq 0, s=0$ if $I Q+r<0$, where $r$ is an unobserved random variable. If $r$ is redundant in the structural equation and in the proxy variable equation for $I Q$, that is, if $\mathrm{E}\left(v \mid \mathbf{z}_{1}\right.$, abil, $\left.I Q, r\right)=0$ and $\mathrm{E}\left(e \mid \mathbf{z}_{1}, I Q, r\right)=0$, then $\mathrm{E}\left(u \mid \mathbf{z}_{1}, I Q, r\right)=0$. Since $s$ is a function of $I Q$ and $r$, it follows immediately that $\mathrm{E}\left(u \mid \mathbf{z}_{1}, I Q, s\right)=0$. Therefore, using OLS on the sample for which $I Q$ is observed yields consistent estimators.

If $r$ is correlated with either $v$ or $e, \mathrm{E}\left(u \mid \mathbf{z}_{1}, I Q, s\right) \neq \mathrm{E}(u)$ in general, and OLS estimation of equation (19.41) using the selected sample would not consistently estimate $\boldsymbol{\delta}_{1}$ and $\theta_{1}$. Therefore, even though $I Q$ is exogenous in the population equation (19.41), the sample selection is not exogenous. In Section 19.6.2 we cover a method that can be used to correct for sample selection bias.

Theorem 19.1 has other useful applications. Suppose that $\mathbf{x}$ is exogenous in equation (19.30) and that $s$ is a nonrandom function of $(\mathbf{x}, v)$, where $v$ is a variable not appearing in equation (19.30). If ( $u, v$ ) is independent of $\mathbf{x}$, then $\mathrm{E}(u \mid \mathbf{x}, v)=\mathrm{E}(u \mid v)$, and so\\
$\mathrm{E}(y \mid \mathbf{x})=\mathbf{x} \boldsymbol{\beta}+\mathrm{E}(u \mid \mathbf{x}, v)=\mathbf{x} \boldsymbol{\beta}+\mathrm{E}(u \mid v)$.\\
If we make an assumption about the functional form of $\mathrm{E}(u \mid v)$, for example, $\mathrm{E}(u \mid v)=\gamma v$, then we can write\\
$y=\mathbf{x} \boldsymbol{\beta}+y v+e, \quad \mathrm{E}(e \mid \mathbf{x}, v)=0$,\\
where $e=u-\mathrm{E}(u \mid v)$. Because $s$ is just a function of $(\mathbf{x}, v), \mathrm{E}(e \mid \mathbf{x}, v, s)=0$, and so $\boldsymbol{\beta}$ and $\gamma$ can be estimated consistently by the OLS regression $y$ on $\mathbf{x}, v$, using the selected sample. Effectively, including $v$ in the regression on the selected subsample eliminates the sample selection problem and allows us to consistently estimate $\boldsymbol{\beta}$. (Incidentally, because $v$ is independent of $\mathbf{x}$, we would not have to include it in equation (19.30) to consistently estimate $\boldsymbol{\beta}$ if we had a random sample from the population. However, including $v$ would result in an asymptotically more efficient estimator of $\boldsymbol{\beta}$ when $\operatorname{Var}(y \mid \mathbf{x}, v)$ is homoskedastic. See problem 4.5.) In Section 19.7 we will see how equation (19.42) can be implemented when $v$ depends on unknown parameters.

\subsection*{19.4.2 Nonlinear Models}
Results similar to those in the previous section hold for nonlinear models as well. We will cover explicitly the case of nonlinear regression and maximum likelihood. See problem 19.11 for the GMM case.

In the nonlinear regression case, if $\mathrm{E}(y \mid \mathbf{x}, s)=\mathrm{E}(y \mid \mathbf{x})$-so that selection is exogenous in the conditional mean sense-then NLS on the selected sample is consistent. Sufficient is that $s$ is a deterministic function of $\mathbf{x}$. The consistency argument is simple: NLS on the selected sample solves


\begin{equation*}
\min _{\boldsymbol{\beta}} N^{-1} \sum_{i=1}^{N} s_{i}\left[y_{i}-m\left(\mathbf{x}_{i}, \boldsymbol{\beta}\right)\right]^{2} \tag{19.43}
\end{equation*}


so it suffices to show that $\boldsymbol{\beta}_{\mathrm{o}}$ in $\mathrm{E}(y \mid \mathbf{x})=m\left(\mathbf{x}, \boldsymbol{\beta}_{\mathrm{o}}\right)$ minimizes $\mathrm{E}\left\{s[y-m(\mathbf{x}, \boldsymbol{\beta})]^{2}\right\}$ over $\boldsymbol{\beta}$. By iterated expectations,\\
$\mathrm{E}\left\{s[y-m(\mathbf{x}, \boldsymbol{\beta})]^{2}\right\}=\mathrm{E}\left(s \mathrm{E}\left\{[y-m(\mathbf{x}, \boldsymbol{\beta})]^{2} \mid \mathbf{x}, s\right\}\right)$\\
Next, write $[y-m(\mathbf{x}, \boldsymbol{\beta})]^{2}=u^{2}+2\left[m\left(\mathbf{x}, \boldsymbol{\beta}_{\mathrm{o}}\right)-m(\mathbf{x}, \boldsymbol{\beta})\right] u+\left[m\left(\mathbf{x}, \boldsymbol{\beta}_{\mathrm{o}}\right)-m(\mathbf{x}, \boldsymbol{\beta})\right]^{2}$, where $u=y-m\left(\mathbf{x}, \boldsymbol{\beta}_{\mathrm{o}}\right)$. By assumption, $\mathrm{E}(u \mid \mathbf{x}, s)=0$. Therefore,\\
$\mathrm{E}\left\{[y-m(\mathbf{x}, \boldsymbol{\beta})]^{2} \mid \mathbf{x}, s\right\}=\mathrm{E}\left(u^{2} \mid \mathbf{x}, s\right)+\left[m\left(\mathbf{x}, \boldsymbol{\beta}_{\mathrm{o}}\right)-m(\mathbf{x}, \boldsymbol{\beta})\right]^{2}$,\\
and the second term is clearly minimized at $\boldsymbol{\beta}=\boldsymbol{\beta}_{\mathrm{o}}$. We do have to assume that $\boldsymbol{\beta}_{\mathrm{o}}$ is the unique value of $\boldsymbol{\beta}$ that makes $\mathrm{E}\left\{s\left[m(\mathbf{x}, \boldsymbol{\beta})-m\left(\mathbf{x}, \boldsymbol{\beta}_{\mathrm{o}}\right)\right]^{2}\right\}$ zero. This is the identification condition on the subpopulation.

It can also be shown that, if $\operatorname{Var}(y \mid \mathbf{x}, s)=\operatorname{Var}(y \mid \mathbf{x})$ and $\operatorname{Var}(y \mid \mathbf{x})=\sigma_{\mathrm{o}}^{2}$, then the usual, nonrobust NLS statistics are valid. If heteroskedasticity exists either in the population or the subpopulation, standard heteroskedasticity-robust inference can be used. The arguments are very similar to those for 2SLS in the previous subsection.

Another important case is the general conditional maximum likelihood setup. Assume that the distribution of $\mathbf{y}$ given $\mathbf{x}$ and $s$ is the same as the distribution of $\mathbf{y}$ given $\mathbf{x}: \mathrm{D}(\mathbf{y} \mid \mathbf{x}, s)=\mathrm{D}(\mathbf{y} \mid \mathbf{x})$. This is a stronger form of ignorability of selection, but it always holds if $s$ is a nonrandom function of $\mathbf{x}$, or if $s$ is independent of $(\mathbf{x}, \mathbf{y})$. In any case, $\mathrm{D}(\mathbf{y} \mid \mathbf{x}, s)=\mathrm{D}(\mathbf{y} \mid \mathbf{x})$ ensures that the MLE on the selected sample is consistent and that the usual MLE statistics are valid. The analogy argument should be familiar by now. Conditional MLE on the selected sample solves\\
$\max _{\theta} N^{-1} \sum_{i=1}^{N} s_{i} \ell\left(\mathbf{y}_{i}, \mathbf{x}_{i} ; \boldsymbol{\theta}\right)$,\\
where $\ell\left(\mathbf{y}_{i}, \mathbf{x}_{i} ; \boldsymbol{\theta}\right)$ is the $\log$ likelihood for observation $i$. Now for each $\mathbf{x}, \boldsymbol{\theta}_{\mathrm{o}}$ maximizes $\mathrm{E}[\ell(\mathbf{y}, \mathbf{x} ; \boldsymbol{\theta}) \mid \mathbf{x}]$ over $\boldsymbol{\theta}$. But $\mathrm{E}[s \ell(\mathbf{y}, \mathbf{x} ; \boldsymbol{\theta})]=\mathrm{E}\{s \mathrm{E}[\ell(\mathbf{y}, \mathbf{x} ; \boldsymbol{\theta}) \mid \mathbf{x}, s]\}=\mathrm{E}\{s \mathrm{E}[\ell(\mathbf{y}, \mathbf{x} ; \boldsymbol{\theta}) \mid \mathbf{x}]\}$, since, by assumption, the conditional distribution of $\mathbf{y}$ given $(\mathbf{x}, s)$ does not depend on $s$. Since $\mathrm{E}[\ell(\mathbf{y}, \mathbf{x} ; \boldsymbol{\theta}) \mid \mathbf{x}]$ is maximized at $\boldsymbol{\theta}_{\mathrm{o}}$, so is $\mathrm{E}\{s \mathrm{E}[\ell(\mathbf{y}, \mathbf{x} ; \boldsymbol{\theta}) \mid \mathbf{x}]\}$. We must make the stronger assumption that $\boldsymbol{\theta}_{\mathrm{o}}$ is the unique maximum, just as in the previous cases: if the selected subset of the population is too small, we may not be able to identify $\boldsymbol{\theta}_{\mathrm{o}}$. Inference can be carried out using the usual MLE statistics obtained from the selected subsample because the information equality now holds conditional on $\mathbf{x}$ and $s$ under the assumption that $\mathrm{D}(\mathbf{y} \mid \mathbf{x}, s)=\mathrm{D}(\mathbf{y} \mid \mathbf{x})$. We omit the details.

Problem 19.11 asks you to work through the case of GMM estimation of general nonlinear models based on conditional moment restrictions.

\subsection*{19.5 Selection on the Basis of the Response Variable: Truncated Regression}
Let $\left(\mathbf{x}_{i}, y_{i}\right)$ denote a random draw from a population. In this section we explicitly treat the case where the sample is selected on the basis of $y_{i}$.

In applying the following methods it is important to remember that there is an underlying population of interest, often described by a linear conditional expectation: $\mathrm{E}\left(y_{i} \mid \mathbf{x}_{i}\right)=\mathbf{x}_{i} \boldsymbol{\beta}$. If we could observe a random sample from the population, then we would just use standard regression analysis. The problem comes about because the sample we can observe is chosen at least partly based on the value of $y_{i}$. Unlike in the case where selection is based only on $\mathbf{x}_{i}$, selection based on $y_{i}$ causes problems for standard OLS analysis on the selected sample.

A classic example of selection based on $y_{i}$ is Hausman and Wise's (1977) study of the determinants of earnings. Hausman and Wise recognized that their sample from a negative income tax experiment was truncated because only families with income below 1.5 times the poverty level were allowed to participate in the program; no data\\
were available on families with incomes above the threshold value. The truncation rule was known, and so the effects of truncation could be accounted for.

A similar example is example 19.2. We do not observe data on families with wealth above $\$ 200,000$. This case is different from the top coding example we discussed in Section 19.2.3. Here, we observe nothing about families with high wealth: they are entirely excluded from the sample. In the top coding case, we have a random sample of families, and we always observe $\mathbf{x}_{i}$; the information on $\mathbf{x}_{i}$ is useful even if wealth is top coded.

We assume that $y_{i}$ is a continuous random variable and that the selection rule takes the form\\
$s_{i}=1\left[a_{1}<y_{i}<a_{2}\right]$,\\
where $a_{1}$ and $a_{2}$ are known constants such that $a_{1}<a_{2}$. A good way to think of the sample selection in this case is that we draw $\left(\mathbf{x}_{i}, y_{i}\right)$ randomly from the population. If $y_{i}$ falls in the interval $\left(a_{1}, a_{2}\right)$, then we observe both $y_{i}$ and $\mathbf{x}_{i}$. If $y_{i}$ is outside this interval, then we do not observe $y_{i}$ or $\mathbf{x}_{i}$. Thus, all we know is that there is some subset of the population that does not enter our data set because of the selection rule. We know how to characterize the part of the population not being sampled because we know the constants $a_{1}$ and $a_{2}$.

In most applications we are still interested in estimating $\mathrm{E}\left(y_{i} \mid \mathbf{x}_{i}\right)=\mathbf{x}_{i} \boldsymbol{\beta}$. However, because of sample selection based on $y_{i}$, we must-at least in a parametric contextspecify a full conditional distribution of $y_{i}$ given $\mathbf{x}_{i}$. Parameterize the conditional density of $y_{i}$ given $\mathbf{x}_{i}$ by $f\left(\cdot \mid \mathbf{x}_{i} ; \boldsymbol{\beta}, \boldsymbol{\gamma}\right)$, where $\boldsymbol{\beta}$ are the conditional mean parameters and $\gamma$ is a $G \times 1$ vector of additional parameters. The cdf of $y_{i}$ given $\mathbf{x}_{i}$ is $F\left(\cdot \mid \mathbf{x}_{i} ; \boldsymbol{\beta}, \boldsymbol{\gamma}\right)$.

What we can use in estimation is the density of $y_{i}$ conditional on $\mathbf{x}_{i}$ and the fact that we observe $\left(y_{i}, \mathbf{x}_{i}\right)$. In other words, we must condition on $a_{1}<y_{i}<a_{2}$ or, equivalently, $s_{i}=1$. The cdf of $y_{i}$ conditional on ( $\mathbf{x}_{i}, s_{i}=1$ ) is simply\\
$\mathrm{P}\left(y_{i} \leq y \mid \mathbf{x}_{i}, s_{i}=1\right)=\frac{\mathrm{P}\left(y_{i} \leq y, s_{i}=1 \mid \mathbf{x}_{i}\right)}{\mathrm{P}\left(s_{i}=1 \mid \mathbf{x}_{i}\right)}$.\\
Because $y_{i}$ is continuously distributed, $\mathrm{P}\left(s_{i}=1 \mid \mathbf{x}_{i}\right)=\mathrm{P}\left(a_{1}<y_{i}<a_{2} \mid \mathbf{x}_{i}\right)=F\left(a_{2} \mid \mathbf{x}_{i} ; \boldsymbol{\beta}, \boldsymbol{\gamma}\right) -F\left(a_{1} \mid \mathbf{x}_{i} ; \boldsymbol{\beta}, \boldsymbol{\gamma}\right)>0$ for all possible values of $\mathbf{x}_{i}$. The case $a_{2}=\infty$ corresponds to truncation only from below, in which case $F\left(a_{2} \mid \mathbf{x}_{i} ; \boldsymbol{\beta}, \boldsymbol{\gamma}\right) \equiv 1$. If $a_{1}=-\infty$ (truncation only from above), then $F\left(a_{1} \mid \mathbf{x}_{i} ; \boldsymbol{\beta}, \boldsymbol{\gamma}\right)=0$. To obtain the numerator when $a_{1}<y< a_{2}$, we have\\
$\mathrm{P}\left(y_{i} \leq y, s_{i}=1 \mid \mathbf{x}_{i}\right)=\mathrm{P}\left(a_{1}<y_{i} \leq y \mid \mathbf{x}_{i}\right)=F\left(y \mid \mathbf{x}_{i} ; \boldsymbol{\beta}, \boldsymbol{\gamma}\right)-F\left(a_{1} \mid \mathbf{x}_{i} ; \boldsymbol{\beta}, \boldsymbol{\gamma}\right)$.

When we put this equation over $\mathrm{P}\left(s_{i}=1 \mid \mathbf{x}_{i}\right)$ and take the derivative with respect to the dummy argument $y$, we obtain the density of $y_{i}$ given $\left(\mathbf{x}_{i}, s_{i}=1\right)$ :


\begin{equation*}
p\left(y \mid \mathbf{x}_{i}, s_{i}=1\right)=\frac{f\left(y \mid \mathbf{x}_{i} ; \boldsymbol{\beta}, \gamma\right)}{F\left(a_{2} \mid \mathbf{x}_{i} ; \boldsymbol{\beta}, \gamma\right)-F\left(a_{1} \mid \mathbf{x}_{i} ; \boldsymbol{\beta}, \gamma\right)} \tag{19.45}
\end{equation*}


for $a_{1}<y<a_{2}$.\\
Given a model for $f(y \mid \mathbf{x} ; \boldsymbol{\beta}, \boldsymbol{\gamma})$, the log-likelihood function for any $\left(\mathbf{x}_{i}, y_{i}\right)$ in the sample can be obtained by plugging $y_{i}$ into equation (19.45) and taking the log. The CMLEs of $\boldsymbol{\beta}$ and $\boldsymbol{\gamma}$ using the selected sample are efficient in the class of estimators that do not use information about the distribution of $\mathbf{x}_{i}$. Standard errors and test statistics can be computed using the general theory of conditional MLE.

In most applications of truncated samples, the population conditional distribution is assumed to be $\operatorname{Normal}\left(\mathbf{x} \boldsymbol{\beta}, \sigma^{2}\right)$, in which case we have the truncated Tobit model or truncated normal regression model. The truncated Tobit model is related to the censored Tobit model for data-censoring applications (see Section 19.2.3), but there is a key difference: in censored regression, we have a random sample of units and we observe the covariates $\mathbf{x}$ for all people, even those for whom the response in not known. If we drop observations entirely when the response is not observed, we obtain the truncated regression model. If in a top coding example we use the information in the top coded observations, we are in the censored regression case. If we drop all top coded observations, we are in the truncated regression case. (Given a choice, we should use a censored regression analysis, as it uses all of the information in the sample.)

From our analysis of the censored regression model in Section 19.2.3, it is not surprising that heteroskedasticity or nonnormality in truncated regression results in inconsistent estimators of $\boldsymbol{\beta}$. This outcome is unfortunate because, if not for the sample selection problem, we could consistently estimate $\boldsymbol{\beta}$ under $\mathrm{E}(y \mid \mathbf{x})=\mathbf{x} \boldsymbol{\beta}$, without specifying $\operatorname{Var}(y \mid \mathbf{x})$ or the conditional distribution. Distribution-free methods for the truncated regression model have been suggested by Powell (1986) under the assumption of a symmetric error distribution; see Powell (1994) for a recent survey.

Truncating a sample on the basis of $y$ is related to choice-based sampling. Traditional choice-based sampling applies when $y$ is a discrete response taking on a finite number of values, where sampling frequencies differ depending on the outcome of $y$. (In the truncation case, the sampling frequency is one when $y$ falls in the interval ( $a_{1}, a_{2}$ ) and zero when $y$ falls outside of the interval.) We do not cover choice-based sampling here; see Manksi and McFadden (1981), Imbens (1992), and Cosslett\\
(1993). In Section 20.2 we cover some estimation methods for stratified sampling, which can be applied to some choice-based samples.

\subsection*{19.6 Incidental Truncation: A Probit Selection Equation}
We now turn to sample selection corrections when selection is determined by a probit model. This setup applies to problems different from those in Section 19.5, where the problem was that a survey or program was designed to intentionally exclude part of the population. We are now interested in selection problems that are due to incidental truncation, attrition in the context of program evalution, and general nonresponse that leads to missing data on the response variable or the explanatory variables.

\subsection*{19.6.1 Exogenous Explanatory Variables}
The incidental truncation problem is motivated by Gronau's (1974) model of the wage offer and labor force participation.

Example 19.5 (Labor Force Participation and the Wage Offer): Interest lies in estimating $\mathrm{E}\left(w_{i}^{\mathrm{o}} \mid \mathbf{x}_{i}\right)$, where $w_{i}^{\mathrm{o}}$ is the hourly wage offer for a randomly drawn individual $i$. If $w_{i}^{\mathrm{o}}$ were observed for everyone in the (working age) population, we would proceed in a standard regression framework. However, a potential sample selection problem arises because $w_{i}^{\mathrm{o}}$ is observed only for people who work.

We can cast this problem as a weekly labor supply model:


\begin{equation*}
\max _{h} \operatorname{utili}_{i}\left(w_{i}^{\mathrm{o}} h+a_{i}, h\right) \quad \text { subject to } 0 \leq h \leq 168 \tag{19.46}
\end{equation*}


where $h$ is hours worked per week and $a_{i}$ is nonwage income of person $i$. Let $s_{i}(h) \equiv \operatorname{util}_{i}\left(w_{i}^{\circ} h+a_{i}, h\right)$, and assume that we can rule out the solution $h_{i}=168$. Then the solution can be $h_{i}=0$ or $0<h_{i}<168$. If $\mathrm{d} s_{i} / \mathrm{d} h \leq 0$ at $h=0$, then the optimum is $h_{i}=0$. Using this condition, straightforward algebra shows that $h_{i}=0$ if and only if


\begin{equation*}
w_{i}^{\mathrm{o}} \leq-m u_{i}^{h}\left(a_{i}, 0\right) / m u_{i}^{q}\left(a_{i}, 0\right), \tag{19.47}
\end{equation*}


where $m u_{i}^{h}(\cdot, \cdot)$ is the marginal disutility of working and $m u_{i}^{q}(\cdot, \cdot)$ is the marginal utility of income. Gronau (1974) called the right-hand side of equation (19.47) the reservation wage, $w_{i}^{r}$, which is assumed to be strictly positive.

We now make the parametric assumptions


\begin{equation*}
w_{i}^{o}=\exp \left(\mathbf{x}_{i 1} \boldsymbol{\beta}_{1}+u_{i 1}\right), \quad w_{i}^{r}=\exp \left(\mathbf{x}_{i 2} \boldsymbol{\beta}_{2}+\gamma_{2} a_{i}+u_{i 2}\right) \tag{19.48}
\end{equation*}


where ( $u_{i 1}, u_{i 2}$ ) is independent of ( $\mathbf{x}_{i 1}, \mathbf{x}_{i 2}, a_{i}$ ). Here, $\mathbf{x}_{i 1}$ contains productivity characteristics, and possibly demographic characteristics, of individual $i$, and $\mathbf{x}_{i 2}$ contains\\
variables that determine the marginal utility of leisure and income; these may overlap with $\mathbf{x}_{i 1}$. From equation (19.48) we have the log wage equation\\
$\log w_{i}^{o}=\mathbf{x}_{i 1} \boldsymbol{\beta}_{1}+u_{i 1}$.

But the wage offer $w_{i}^{\mathrm{o}}$ is observed only if the person works, that is, only if $w_{i}^{\mathrm{o}}>w_{i}^{r}$, or $\log w_{i}^{\mathrm{o}}-\log w_{i}^{r}=\mathbf{x}_{i 1} \boldsymbol{\beta}_{1}-\mathbf{x}_{i 2} \boldsymbol{\beta}_{2}-\gamma_{2} a_{i}+u_{i 1}-u_{i 2} \equiv \mathbf{x}_{i} \boldsymbol{\delta}_{2}+v_{i 2}>0$.

This behavior introduces a potential sample selection problem if we use data only on working people to estimate equation (19.49).

This example differs in an important respect from top coding examples. With top coding, the censoring rule is known for each unit in the population. In Gronau's example, we do not know $w_{i}^{r}$, so we cannot use $w_{i}^{\mathrm{o}}$ in a censored regression analysis. If $w_{i}^{r}$ were observed and exogenous and $\mathbf{x}_{i 1}$ were always observed, then we would be in the censored regression framework with censoring from below. If $w_{i}^{r}$ were observed and exogenous but $\mathbf{x}_{i 1}$ were observed only when $w_{i}^{\mathrm{o}}$ is, we would be in the truncated Tobit framework. But $w_{i}^{r}$ is allowed to depend on unobservables, and so we need a new framework.

If we drop the $i$ subscript, let $y_{1} \equiv \log w^{\mathrm{o}}$, and let $y_{2}$ be the binary labor force participation indicator, Gronau's model can be written for a random draw from the population as


\begin{align*}
& y_{1}=\mathbf{x}_{1} \boldsymbol{\beta}_{1}+u_{1}  \tag{19.50}\\
& y_{2}=1\left[\mathbf{x} \boldsymbol{\delta}_{2}+v_{2}>0\right] \tag{19.51}
\end{align*}


We discuss estimation of this model under the following set of assumptions:\\
ASSUMPTION 19.1: (a) $\left(\mathbf{x}, y_{2}\right)$ are always observed, $y_{1}$ is observed only when $y_{2}=1$; (b) ( $u_{1}, v_{2}$ ) is independent of $\mathbf{x}$ with zero mean; (c) $v_{2} \sim \operatorname{Normal}(0,1)$; and (d) $\mathrm{E}\left(u_{1} \mid v_{2}\right)=\gamma_{1} v_{2}$.

Assumption 19.1a emphasizes the sample selection nature of the problem. Part $b$ is a strong, but standard, form of exogeneity of $\mathbf{x}$. We will see that assumption 19.1c is needed to derive a conditional expectation given the selected sample. It is probably the most restrictive assumption because it is an explicit distributional assumption. Assuming $\operatorname{Var}\left(v_{2}\right)=1$ is without loss of generality because $y_{2}$ is a binary variable.

Assumption 19.1d requires linearity in the population regression of $u_{1}$ on $v_{2}$. It always holds if ( $u_{1}, v_{2}$ ) is bivariate normal-a standard assumption in these contexts-but assumption 19.1d holds under weaker assumptions. In particular, we do not need to assume that $u_{1}$ itself is normally distributed.

Amemiya (1985) calls equations (19.50) and (19.51) the Type II Tobit model. This name is fine as a label, but we must understand that it is a model of sample selection, and it has nothing to do with $y_{1}$ being a corner solution outcome. Unfortunately, in almost all treatments of this model, $y_{1}$ is set to zero when $y_{2}=0$. Setting $y_{1}$ to zero (or any value) when $y_{2}=0$ is misleading and can lead to inappropriate use of the model. For example, it makes no sense to set the wage offer to zero just because we do not observe it. As another example, it makes no sense to set the price per dollar of life insurance $\left(y_{1}\right)$ to zero for someone who did not buy life insurance (so $y_{2}=1$ if and only if a person owns a life insurance policy).

We also have some interest in the parameters of the selection equation (19.51); for example, in Gronau's model it is a reduced-form labor force participation equation. In program evaluation with attrition, the selection equation explains the probability of dropping out of the program.

Because of the similarities in the statistical structures for the Type II Tobit model for sample selection and the exponential Type II Tobit model for corner solutions, now is a good time to review the differences in how one applies these models. Recall that when using the ET2T for a corner solution response, the goal is to model the distribution of the fully observable outcome in a flexible way. The possible outcomes include zero-which is a legitimate value and means we need to model the probability of a response at zero, as well as the density for positive values. One can debate about which model is best, but the important point is that there is no missing data problem. (It is for that reason that using the label "selection model" for the hurdle model we discussed in Section 17.6.3 is somewhat misleading.)

The present situation is different from having a corner at zero. In the sample selection setting, the variable $y_{1}$ in equation (19.50), which often is in logarithmic form (in which case the setup appears quite similar to the ET2T model), is not always observed. We are interested in the mean response in the population, $\mathrm{E}\left(y_{1} \mid \mathbf{x}_{1}\right)$, which is assumed to follow a standard linear model. Consequently, once the parameters $\boldsymbol{\beta}_{1}$ have been consistently estimated, it is trivial to interpret the estimated equation because the elements of $\hat{\boldsymbol{\beta}}_{1}$ directly measure the partial effects of interest. In other words, in the sample selection context the parameter estimates are directly interpretable as partial effects on $\mathrm{E}\left(y_{1} \mid \mathbf{x}_{1}\right)$, unlike in the ET2T case, where the various conditional means of the observed response have no simple forms.

We can allow a little more generality in the model by replacing $\mathbf{x}$ in equation (19.51) with $\mathbf{x}_{2}$; then, as will become clear, $\mathbf{x}_{1}$ would only need to be observed whenever $y_{1}$ is, whereas $\mathbf{x}_{2}$ must always be observed. This extension is not especially useful for something like Gronau's model because it implies that $\mathbf{x}_{1}$ contains elements that cannot also appear in $\mathbf{x}_{2}$. Because the selection equation is not typically a\\
structural equation, it is undesirable to impose exclusion restrictions in equation (19.51). If a variable affecting $y_{1}$ is observed only along with $y_{1}$, the instrumental variables method that we cover in Section 19.6.2 is more attractive.

To derive an estimating equation, let $\left(y_{1}, y_{2}, \mathbf{x}, u_{1}, v_{2}\right)$ denote a random draw from the population. Since $y_{1}$ is observed only when $y_{2}=1$, what we can hope to estimate is $\mathrm{E}\left(y_{1} \mid \mathbf{x}, y_{2}=1\right)$ [along with $\left.\mathrm{P}\left(y_{2}=1 \mid \mathbf{x}\right)\right]$. How does $\mathrm{E}\left(y_{1} \mid \mathbf{x}, y_{2}=1\right)$ depend on the vector of interest, $\boldsymbol{\beta}_{1}$ ? First, under assumption 19.1 and equation (19.50),\\
$\mathrm{E}\left(y_{1} \mid \mathbf{x}, v_{2}\right)=\mathbf{x}_{1} \boldsymbol{\beta}_{1}+\mathrm{E}\left(u_{1} \mid \mathbf{x}, v_{2}\right)=\mathbf{x}_{1} \boldsymbol{\beta}_{1}+\mathrm{E}\left(u_{1} \mid v_{2}\right)=\mathbf{x}_{1} \boldsymbol{\beta}_{1}+\gamma_{1} v_{2}$,\\
where the second equality follows because ( $u_{1}, v_{2}$ ) is independent of $\mathbf{x}$. Equation (19.52) is very useful. The first thing to note is that, if $\gamma_{1}=0$-which implies that $u_{1}$ and $v_{2}$ are uncorrelated-then $\mathrm{E}\left(y_{1} \mid \mathbf{x}, v_{2}\right)=\mathrm{E}\left(y_{1} \mid \mathbf{x}\right)=\mathrm{E}\left(y_{1} \mid \mathbf{x}_{1}\right)=\mathbf{x}_{1} \boldsymbol{\beta}_{1}$. Because $y_{2}$ is a function of $\left(\mathbf{x}, v_{2}\right)$, it follows immediately that $\mathrm{E}\left(y_{1} \mid \mathbf{x}, y_{2}\right)=\mathrm{E}\left(y_{1} \mid \mathbf{x}_{1}\right)$. In other words, if $\gamma_{1}=0$, then there is no sample selection problem, and $\boldsymbol{\beta}_{1}$ can be consistently estimated by OLS using the selected sample.

What if $\gamma_{1} \neq 0$ ? Using iterated expectations on equation (19.52),\\
$\mathrm{E}\left(y_{1} \mid \mathbf{x}, y_{2}\right)=\mathbf{x}_{1} \boldsymbol{\beta}_{1}+\gamma_{1} \mathrm{E}\left(v_{2} \mid \mathbf{x}, y_{2}\right)=\mathbf{x}_{1} \boldsymbol{\beta}_{1}+\gamma_{1} h\left(\mathbf{x}, y_{2}\right)$,\\
where $h\left(\mathbf{x}, y_{2}\right)=\mathrm{E}\left(v_{2} \mid \mathbf{x}, y_{2}\right)$. If we knew $h\left(\mathbf{x}, y_{2}\right)$, then, from Theorem 19.1, we could estimate $\boldsymbol{\beta}_{1}$ and $\gamma_{1}$ from the regression $y_{1}$ on $\mathbf{x}_{1}$ and $h\left(\mathbf{x}, y_{2}\right)$, using only the selected sample. Because the selected sample has $y_{2}=1$, we need only find $h(\mathbf{x}, 1)$. But $h(\mathbf{x}, 1)=\mathrm{E}\left(v_{2} \mid v_{2}>-\mathbf{x} \boldsymbol{\delta}_{2}\right)=\lambda\left(\mathbf{x} \boldsymbol{\delta}_{2}\right)$, where $\lambda(\cdot) \equiv \phi(\cdot) / \Phi(\cdot)$ is the inverse Mills ratio, and so we can write\\
$\mathrm{E}\left(y_{1} \mid \mathbf{x}, y_{2}=1\right)=\mathbf{x}_{1} \boldsymbol{\beta}_{1}+\gamma_{1} \lambda\left(\mathbf{x} \boldsymbol{\delta}_{2}\right)$.

Equation (19.53), which can be found in numerous places (see, for example, Heckman, 1979, and Amemiya, 1985) makes it clear that an OLS regression of $y_{1}$ on $\mathbf{x}_{1}$ using the selected sample omits the term $\lambda\left(\mathbf{x} \boldsymbol{\delta}_{2}\right)$ and generally leads to inconsistent estimation of $\boldsymbol{\beta}_{1}$. As pointed out by Heckman (1979), the presence of selection bias can be viewed as an omitted variable problem in the selected sample. An interesting point is that, even though only $\mathbf{x}_{1}$ appears in the population expectation, $\mathrm{E}\left(y_{1} \mid \mathbf{x}\right)$, other elements of $\mathbf{x}$ appear in the expectation on the subpopulation, $\mathrm{E}\left(y_{1} \mid \mathbf{x}, y_{2}=1\right)$.

Equation (19.53) also suggests a way to consistently estimate $\boldsymbol{\beta}_{1}$. Following Heckman (1979), we can consistently estimate $\boldsymbol{\beta}_{1}$ and $\gamma_{1}$ using the selected sample by regressing $y_{i 1}$ on $\mathbf{x}_{i 1}, \lambda\left(\mathbf{x}_{i} \boldsymbol{\delta}_{2}\right)$. The problem is that $\boldsymbol{\delta}_{2}$ is unknown, so we cannot compute the additional regressor $\lambda\left(\mathbf{x}_{i} \boldsymbol{\delta}_{2}\right)$. Nevertheless, a consistent estimator of $\boldsymbol{\delta}_{2}$ is available from the first-stage probit estimation of the selection equation.

Procedure 19.1: (a) Obtain the probit estimate $\hat{\boldsymbol{\delta}}_{2}$ from the model


\begin{equation*}
\mathrm{P}\left(y_{i 2}=1 \mid \mathbf{x}_{i}\right)=\Phi\left(\mathbf{x}_{i} \boldsymbol{\delta}_{2}\right) \tag{19.54}
\end{equation*}


using all $N$ observations. Then obtain the estimated inverse Mills ratios $\hat{\lambda}_{i 2} \equiv \lambda\left(\mathbf{x}_{i} \hat{\boldsymbol{\delta}}_{2}\right)$ (at least for $i=1, \ldots, N_{1}$ ).\\
(b) Obtain $\hat{\boldsymbol{\beta}}_{1}$ and $\hat{\gamma}_{1}$ from the OLS regression on the selected sample,\\
$y_{i 1}$ on $\mathbf{x}_{i 1}, \hat{\lambda}_{i 2}, \quad i=1,2, \ldots, N_{1}$.

These estimators are consistent and $\sqrt{N}$-asymptotically normal.\\
The procedure is sometimes called Heckit after Heckman (1976) and the tradition of putting "it" on the end of procedures related to probit (such as Tobit).

A very simple test for selection bias is available from regression (19.55). Under the null of no selection bias, $\mathrm{H}_{0}: \gamma_{1}=0$, we have $\operatorname{Var}\left(y_{1} \mid \mathbf{x}, y_{2}=1\right)=\operatorname{Var}\left(y_{1} \mid \mathbf{x}\right)= \operatorname{Var}\left(u_{1}\right)$, and so homoskedasticity holds under $\mathrm{H}_{0}$. Further, from the results on generated regressors in Chapter 6, the asymptotic variance of $\hat{\gamma}_{1}$ (and $\hat{\boldsymbol{\beta}}_{1}$ ) is not affected by $\hat{\boldsymbol{\delta}}_{2}$ when $\gamma_{1}=0$. Thus, a standard $t$ test on $\hat{\gamma}_{1}$ is a valid test of the null hypothsesis of no selection bias.

When $\gamma_{1} \neq 0$, obtaining a consistent estimate for the asymptotic variance of $\hat{\boldsymbol{\beta}}_{1}$ is complicated for two reasons. The first is that, if $\gamma_{1} \neq 0$, then $\operatorname{Var}\left(y_{1} \mid \mathbf{x}, y_{2}=1\right)$ is not constant. As we know, heteroskedasticity itself is easy to correct for using the robust standard errors. However, we should also account for the fact that $\hat{\boldsymbol{\delta}}_{2}$ is an estimator of $\boldsymbol{\delta}_{2}$. The adjustment to the variance of $\left(\hat{\boldsymbol{\beta}}_{1}, \hat{\gamma}_{1}\right)$ because of the two-step estimation is cumbersome-it is not enough to simply make the standard errors heteroskedasticityrobust. Some statistical packages now have this feature built in.

As a technical point, we do not need $\mathbf{x}_{1}$ to be a strict subset of $\mathbf{x}$ for $\boldsymbol{\beta}_{1}$ to be identified, and procedure 19.1 does carry through when $\mathbf{x}_{1}=\mathbf{x}$. However, if $\mathbf{x}_{i} \hat{\boldsymbol{\delta}}_{2}$ does not have much variation in the sample, then $\hat{\lambda}_{i 2}$ can be approximated well by a linear function of $\mathbf{x}$. If $\mathbf{x}=\mathbf{x}_{1}$, this correlation can introduce severe collinearity among the regressors in regression (19.55), which can lead to large standard errors of the elements of $\hat{\boldsymbol{\beta}}_{1}$. When $\mathbf{x}_{1}=\mathbf{x}, \boldsymbol{\beta}_{1}$ is identified only due to the nonlinearity of the inverse Mills ratio.

The situation is not quite as bad as in Section 9.5.1. There, identification failed for certain values of the structural parameters. Here, we still have identification for any value of $\boldsymbol{\beta}_{1}$ in equation (19.50), but it is unlikely we can estimate $\boldsymbol{\beta}_{1}$ with much precision. Even if we can, we would have to wonder whether a statistically significant inverse Mills ratio term is due to sample selection or functional form misspecification in the population model (19.50).

\begin{table}[h]
\begin{center}
\captionsetup{labelformat=empty}
\caption{Table 19.1\\
Wage Offer Equation for Married Women}
\begin{tabular}{|l|l|l|}
\hline
\multicolumn{3}{|l|}{Dependent Variable: $\log$ (wage)} \\
\hline
Independent Variable & OLS & Heckit \\
\hline
\multirow[t]{2}{*}{educ} & . 108 & . 109 \\
\hline
 & (.014) & (.016) \\
\hline
\multirow[t]{2}{*}{exper} & . 042 & . 044 \\
\hline
 & (.012) & (.016) \\
\hline
\multirow[t]{2}{*}{exper ${ }^{2}$} & -. 00081 & -. 00086 \\
\hline
 & (.00039) & (.00044) \\
\hline
\multirow[t]{2}{*}{constant} & -. 522 & -. 578 \\
\hline
 & (.199) & (.307) \\
\hline
\multirow[t]{2}{*}{$\hat{\lambda}_{2}$} & - & . 032 \\
\hline
 &  & (.134) \\
\hline
Sample size & 428 & 428 \\
\hline
$R$-squared & . 157 & . 157 \\
\hline
\end{tabular}
\end{center}
\end{table}

Example 19.6 (Wage Offer Equation for Married Women): We use the data in MROZ.RAW to estimate a wage offer function for married women, accounting for potential selectivity bias into the workforce. Of the 753 women, we observe the wage offer for 428 working women. The labor force participation equation contains the variables in Table 15.1, including other income, age, number of young children, and number of older children—in addition to educ, exper, and exper ${ }^{2}$. The results of OLS on the selected sample and the Heckit method are given in Table 19.1.

The differences between the OLS and Heckit estimates are practically small, and the inverse Mills ratio term is statistically insignificant. The fact that the intercept estimates differ somewhat is usually unimportant. (The standard errors reported for Heckit are the unadjusted ones from regression (19.55). If $\hat{\lambda}_{2}$ were statistically significant, we should obtain the corrected standard errors.)

The Heckit results in Table 19.1 use four exclusion restrictions in the structural equation, because nwifeinc, age, kidsltb, and kidsge6 are all excluded from the wage offer equation. If we allow all variables in the selection equation to also appear in the wage offer equation, the Heckit estimates become very imprecise. The coefficient on educ becomes $.119(\mathrm{se}=.034)$, compared with the OLS estimate $.100(\mathrm{se}=.015)$. The coefficient on kidsltb-which now appears in the wage offer equation-is - .188 $(\mathrm{se}=.232)$ in the Heckit estimation, and $-.056(\mathrm{se}=.009)$ in the OLS estimation. The imprecision of the Heckit estimates is due to the severe collinearity that comes from adding $\hat{\lambda}_{2}$ to the equation, because $\hat{\lambda}_{2}$ is now a function only of the explanatory variables in the wage offer equation. In fact, using the selected sample, regressing $\hat{\lambda}_{2}$ on the seven explanatory variables gives $R$-squared $=.962$. Unfortunately, comparing\\
the OLS and Heckit results does not allow us to resolve some important issues. For example, the OLS results suggest that another young child reduces the wage offer by about 5.6 percent ( $t$ statistic $\approx-6.2$ ), other things being equal. Is this effect real, or is it simply due to our inability to adequately correct for sample selection bias? Unless we have a variable that affects labor force participation without affecting the wage offer, we cannot answer this question.

If we replace parts c and d in assumption 19.1 with the stronger assumption that $\left(u_{1}, v_{2}\right)$ is bivariate normal with mean zero, $\operatorname{Var}\left(u_{1}\right)=\sigma_{1}^{2}, \operatorname{Cov}\left(u_{1}, v_{2}\right)=\sigma_{12}$, and $\operatorname{Var}\left(v_{2}\right)=1$, then partial maximum likelihood estimation can be used, as described generally in problem 13.7. Partial MLE will be more efficient than the two-step procedure under joint normality of $u_{1}$ and $v_{2}$, and it will produce standard errors and likelihood ratio statistics that can be used directly (this conclusion follows from problem 13.7). The drawbacks are that it is less robust than the two-step procedure and that it is sometimes difficult to get the problem to converge.

The reason we cannot perform full conditional MLE-contrast the situation in Section 17.6.3 for the ET2T model-is that $y_{1}$ is only observed when $y_{2}=1$. Thus, while we can use the full density of $y_{2}$ given $\mathbf{x}$, which is $f\left(y_{2} \mid \mathbf{x}\right)= \left[\Phi\left(\mathbf{x} \boldsymbol{\delta}_{2}\right)\right]^{y_{2}}\left[1-\Phi\left(\mathbf{x} \boldsymbol{\delta}_{2}\right)\right]^{1-y_{2}}, y_{2}=0,1$, we can only use the density $f\left(y_{1} \mid y_{2}, \mathbf{x}\right)$ when $y_{2}=1$. To find $f\left(y_{1} \mid y_{2}, \mathbf{x}\right)$ at $y_{2}=1$, we can use Bayes' rule to write $f\left(y_{1} \mid y_{2}, \mathbf{x}\right)= f\left(y_{2} \mid y_{1}, \mathbf{x}\right) f\left(y_{1} \mid \mathbf{x}\right) / f\left(y_{2} \mid \mathbf{x}\right)$. Therefore, $f\left(y_{1} \mid y_{2}=1, \mathbf{x}\right)=\mathrm{P}\left(y_{2}=1 \mid y_{1}, \mathbf{x}\right) f\left(y_{1} \mid \mathbf{x}\right) / \mathrm{P}\left(y_{2}=1 \mid \mathbf{x}\right)$. But $y_{1} \mid \mathbf{x} \sim \operatorname{Normal}\left(\mathbf{x}_{1} \boldsymbol{\beta}_{1}, \sigma_{1}^{2}\right)$. Further, $y_{2}=1\left[\mathbf{x} \boldsymbol{\delta}_{2}+\sigma_{12} \sigma_{1}^{-2}\left(y_{1}-\mathbf{x}_{1} \boldsymbol{\beta}_{1}\right)\right. \left.+e_{2}>0\right]$, where $e_{2}$ is independent of $\left(\mathbf{x}, y_{1}\right)$ and $e_{2} \sim \operatorname{Normal}\left(0,1-\sigma_{12}^{2} \sigma_{1}^{-2}\right)$ (this conclusion follows from standard conditional distribution results for joint normal random variables). Therefore,\\
$\mathrm{P}\left(y_{2}=1 \mid y_{1}, \mathbf{x}\right)=\Phi\left\{\left[\mathbf{x} \boldsymbol{\delta}_{2}+\sigma_{12} \sigma_{1}^{-2}\left(y_{1}-\mathbf{x}_{1} \boldsymbol{\beta}_{1}\right)\right]\left(1-\sigma_{12}^{2} \sigma_{1}^{-2}\right)^{-1 / 2}\right\}$.\\
Combining all of these pieces [and noting the cancellation of $\left.\mathrm{P}\left(y_{2}=1 \mid \mathbf{x}\right)\right]$ we get

$$
\begin{aligned}
\ell_{i}(\boldsymbol{\theta})= & \left(1-y_{i 2}\right) \log \left[1-\Phi\left(\mathbf{x}_{i} \boldsymbol{\delta}_{2}\right)\right]+y_{i 2}\left(\operatorname { l o g } \Phi \left\{\left[\mathbf{x}_{i} \boldsymbol{\delta}_{2}+\sigma_{12} \sigma_{1}^{-2}\left(y_{i 1}-\mathbf{x}_{i 1} \boldsymbol{\beta}_{1}\right)\right]\right.\right. \\
& \left.\left.\times\left(1-\sigma_{12}^{2} \sigma_{1}^{-2}\right)^{-1 / 2}\right\}+\log \phi\left[\left(y_{i 1}-\mathbf{x}_{i 1} \boldsymbol{\beta}_{1}\right) / \sigma_{1}\right]-\log \left(\sigma_{1}\right)\right)
\end{aligned}
$$

The partial $\log$ likelihood is obtained by summing $\ell_{i}(\boldsymbol{\theta})$ across all observations; $y_{i 2}=1$ picks out when $y_{i 1}$ is observed and therefore contains information for estimating $\boldsymbol{\beta}_{1}$.

Ahn and Powell (1993) show how to consistently estimate $\boldsymbol{\beta}_{1}$ without making any distributional assumptions; in particular, the selection equation need not have the probit form. Vella (1998) contains a useful survey.

\subsection*{19.6.2 Endogenous Explanatory Variables}
We now study the sample selection model when one of the elements of $\mathbf{x}_{1}$ is thought to be correlated with $u_{1}$. Or, all the elements of $\mathbf{x}_{1}$ are exogenous in the population model but data are missing on an element of $\mathbf{x}_{1}$, and the reason data are missing might be systematically related to $u_{1}$. For simplicity, we focus on the case of a single endogenous explanatory variable. Having multiple endogenous explanatory variables adds no complications beyond the usual identification conditions.

The model in the population is


\begin{equation*}
y_{1}=\mathbf{z}_{1} \boldsymbol{\delta}_{1}+\alpha_{1} y_{2}+u_{1} \tag{19.56}
\end{equation*}



\begin{equation*}
y_{2}=\mathbf{z}_{2} \boldsymbol{\delta}_{2}+v_{2} \tag{19.57}
\end{equation*}


$y_{3}=1\left[\mathbf{z} \boldsymbol{\delta}_{3}+v_{3}>0\right]$.

The first equation is the structural equation of interest, the second equation is a linear projection for the potentially endogenous or missing variable $y_{2}$, and the third equation is the selection equation. We allow arbitrary correlation among $u_{1}, v_{2}$, and $v_{3}$.

The setup in equations (19.56)-(19.58) encompasses at least three cases of interest. The first occurs when $y_{2}$ is always observed but is endogenous in equation (19.56). An example is seen when $y_{1}$ is $\log \left(\right.$ wage $\left.^{\mathrm{o}}\right)$ and $y_{2}$ is years of schooling: years of schooling is generally available whether or not someone is in the workforce. The model also applies when $y_{2}$ is observed only along with $y_{1}$, as would happen if $y_{1}= \log \left(\right.$ wage $\left.^{\mathrm{o}}\right)$ and $y_{2}$ is the ratio of the benefits offer to wage offer. As a second example, let $y_{1}$ be the percentage of voters supporting the incumbent in a congressional district, and let $y_{2}$ be intended campaign expenditures. Then $y_{3}=1$ if the incumbent runs for reelection, and we only observe $\left(y_{1}, y_{2}\right)$ when $y_{3}=1$. A third application is to missing data only on $y_{2}$, as in example 19.4 where $y_{2}$ is IQ score. In the last two cases, $y_{2}$ might in fact be exogenous in equation (19.56), but endogenous sample selection effectively makes $y_{2}$ endogenous in the selected sample.

If $y_{1}$ and $y_{2}$ were always observed along with $\mathbf{z}$, we would just estimate equation (19.56) by 2SLS if $y_{2}$ is endogenous. We can use the results from Section 19.4.1 to show that 2SLS with the inverse Mills ratio added to the regressors is consistent. Regardless of the data availability on $y_{1}$ and $y_{2}$, in the second step we use only observations for which both $y_{1}$ and $y_{2}$ are observed.

ASSUMPTION 19.2: (a) $\left(\mathbf{z}, y_{3}\right)$ is always observed, $\left(y_{1}, y_{2}\right)$ is observed when $y_{3}=1$; (b) $\left(u_{1}, v_{3}\right)$ is independent of $\mathbf{z}$; (c) $v_{3} \sim \operatorname{Normal}(0,1)$; (d) $\mathrm{E}\left(u_{1} \mid v_{3}\right)=\gamma_{1} v_{3}$; and (e) $\mathrm{E}\left(\mathbf{z}_{2}^{\prime} v_{2}\right)=\mathbf{0}$ and, writing $\mathbf{z}_{2} \boldsymbol{\delta}_{2}=\mathbf{z}_{1} \boldsymbol{\delta}_{22}+\mathbf{z}_{21} \boldsymbol{\delta}_{22}, \boldsymbol{\delta}_{22} \neq \mathbf{0}$.

Parts $\mathrm{b}, \mathrm{c}$, and d are identical to the corresponding assumptions in assumption 19.1 when all explanatory variables are observed and exogenous. Assumption e is new, resulting from the endogeneity of $y_{2}$ in equation (19.56). It is important to see that assumption 19.2e is identical to the rank condition needed for identifying equation (19.56) in the absence of sample selection.

The relationship between $\mathbf{z}_{2}$ and $\mathbf{z}$ is somewhat subtle. The assumptions allow the possibility that $\mathbf{z}_{2}=\mathbf{z}$, but it is helpful to think of $\mathbf{z}$ including at least one factor that "primarily" affects selection that is not also in $\mathbf{z}_{2}$. This way, we can think of needing at least one instrument for $y_{2}$-as in the case without selection - and then at least one more exogenous variable that affects selection. This thinking forces discipline on us when we (potentially) have the problem of an endogenous explanatory variable and sample selection.

One can choose $\mathbf{z}_{2}=\mathbf{z}$, but then we should have at least two elements in $\mathbf{z}$ that are not in $\mathbf{z}_{1}$. As we will see, choosing $\mathbf{z}_{2}=\mathbf{z}$ means that the (implicit) reduced form for $y_{2}$ will tend to suffer from collinearity because then the inverse Mills ratio will be a function of the same variables, $\mathbf{z}$, appearing linearly as regressors. Unless we are interested in the reduced form parameters, the collinearity introduced by having the same regressors appearing linearly as those appearing in the inverse Mills ratio is not of much concern-remember, 2SLS uses only the fitted values from the reduced form-but it is possible the small-sample performance of the procedure is affected.

Importantly, if we choose $\mathbf{z}_{2}$ to be a strict subset of $\mathbf{z}$, we are not making exclusion restrictions in the reduced form; we are simply choosing between variables that are viewed as instruments for $y_{2}$ versus variables that affect selection. By contrast, we usually do not want to make exclusion restrictions in the selection equation, as the subsequent procedure is inconsistent if those restrictions are violated. Therefore, it is a good idea to choose $\mathbf{z}$ to be the vector of all exogenous variables.

To derive an estimating equation, write (in the population)


\begin{equation*}
y_{1}=\mathbf{z}_{1} \boldsymbol{\delta}_{1}+\alpha_{1} y_{2}+g\left(\mathbf{z}, y_{3}\right)+e_{1}, \tag{19.59}
\end{equation*}


where $g\left(\mathbf{z}, y_{3}\right) \equiv \mathrm{E}\left(u_{1} \mid \mathbf{z}, y_{3}\right)$ and $e_{1} \equiv u_{1}-\mathrm{E}\left(u_{1} \mid \mathbf{z}, y_{3}\right)$. By definition, $\mathrm{E}\left(e_{1} \mid \mathbf{z}, y_{3}\right)=$ 0 . If we knew $g\left(\mathbf{z}, y_{3}\right)$ then, from Theorem 19.1, we could just estimate equation (19.59) by 2SLS on the selected sample ( $y_{3}=1$ ) using instruments $[\mathbf{z}, g(\mathbf{z}, 1)]$. It turns out that we do know $g(\mathbf{z}, 1)$ up to some estimable parameters: $\mathrm{E}\left(u_{1} \mid \mathbf{z}, y_{3}=1\right)= \gamma_{1} \lambda\left(\mathbf{z} \boldsymbol{\delta}_{3}\right)$. Since $\boldsymbol{\delta}_{3}$ can be consistently estimated by probit of $y_{3}$ on $\mathbf{z}$ (using the entire sample), we have the following:

Procedure 19.2: (a) Obtain $\hat{\boldsymbol{\delta}}_{3}$ from probit of $y_{3}$ on $\mathbf{z}$ using all observations. Obtain the estimated inverse Mills ratios, $\hat{\lambda}_{i 3}=\lambda\left(\mathbf{z}_{i} \hat{\boldsymbol{\delta}}_{3}\right)$.\\
(b) Using the selected subsample (for which we observe both $y_{1}$ and $y_{2}$ ), estimate the equation\\
$y_{i 1}=\mathbf{z}_{i 1} \boldsymbol{\delta}_{1}+\alpha_{1} y_{i 2}+\gamma_{1} \hat{\lambda}_{i 3}+$ error $_{i}$\\
by 2SLS, using instruments $\left(\mathbf{z}_{i 2}, \hat{\lambda}_{i 3}\right)$.\\
The steps in this procedure show that identification actually requires that $\mathbf{z}_{21}$ appear in the linear projection of $y_{2}$ onto $\mathbf{z}_{1}, \mathbf{z}_{22}$, and $\lambda\left(\mathbf{z} \boldsymbol{\delta}_{3}\right)$ in the selected subpopulation. It would be unusual if this condition were not true when the rank condition 19.2e holds in the population.

The hypothesis-of-no-selection problem (allowing $y_{2}$ to be endogenous or not), $\mathrm{H}_{0}: \gamma_{1}=0$, is tested using the usual 2SLS $t$ statistic for $\hat{\gamma}_{1}$. When $\gamma_{1} \neq 0$, standard errors and test statistics should be corrected for the generated regressors problem, as in Chapter 6, or using the bootstrap (with subsamples drawn from the complete list of observations).

Example 19.7 (Education Endogenous and Sample Selection): In example 19.6 we now allow educ to be endogenous in the wage offer equation, and we test for sample selection bias. Just as if we did not have a sample selection problem, we need IVs for educ that do not appear in the wage offer equation. As in example 5.3, we use parents' education (motheduc, fatheduc) and husband's education as IVs. In addition, we need some variables that affect labor force participation but not the wage offer; we use the same four variables as in example 19.6. Therefore, all variables except educ (and, of course, the wage offer) are treated as exogenous.

We include all of the exogenous variables in the labor force participation equation: exper, exper ${ }^{2}$, nwifeinc, age, kidslt6, kidsge6, motheduc, fatheduc, and huseduc (not educ). (As it turns out, the three education variables are marginally jointly significant, with $p$-value $=.046$.) If we add the inverse Mills ratio, $\hat{\lambda}_{3}$, to the equation and estimate it by 2SLS, using all of the exogenous variables as instruments-in addition to $\hat{\lambda}_{3}$-the coefficient on $\hat{\lambda}_{3}$ is .040 (se $=.133$ ). Therefore, there is little evidence of sample selection bias. Further, the estimated coefficient on education is .088 $(s e=.021)$, which is very close to the estimate when we drop the IMR: the 2SLS coefficient without the IMR included is .087 ( $\mathrm{se}=.021$ ). Thus, there is no practical consequence of correcting for selection bias.

If we use only the parents' and husband's education variables as instruments for educ, the coefficient on $\hat{\lambda}_{3}$ becomes .036 (se $=.134$ ), and so the evidence for selection is still nonexistent. The education coefficient changes somewhat, to .081 ( $s e=.022$ ), but this is due to the different list of instruments, not the selection correction. Without $\hat{\lambda}_{3}$ in the equation, the 2SLS estimate using the smaller list of instruments is .080\\
(se $=.022$ ). Therefore, in this data set, the estimated return to education changes depending on whether education is allowed to be endogenous, and then somewhat depending on the instruments used for it. But whether a sample selection correction is used has essentially no effect.

Importantly, procedure 19.2 applies to any kind of endogenous variable $y_{2}$, including binary and other discrete variables, without any additional assumptions. Why? Because the reduced form for $y_{2}$ is just a linear projection; we do not have to assume, for example, that $v_{2}$ is normally distributed or even independent of $\mathbf{z}$. As an example, we might wish to look at the effects of participation in a job training program on the subsequent wage offer, accounting for the fact that not all who participated in the program will be employed in the following period ( $y_{2}$ is always observed in this case). If participation is voluntary, an instrument for it might be whether the person was randomly chosen as a potential participant.

In order for the selection correction to work for $y_{2}$ with a variety of distributions, it is important to carry out the estimation exactly as described in procedure 19.2. It is tempting to first obtain fitted values, say $\hat{y}_{2}$, insert these for $y_{2}$, and then apply the Heckman regression. (If $y_{2}$ is always observed, one would probably just estimate a linear reduced form using all the observations; if $y_{2}$ is missing along with $y_{1}$, one would probably use a standard Heckman correction on the reduced form.) The problem with using fitted values in place of $y_{2}$, and then, say, running the regression $y_{i 1}$ on $\mathbf{z}_{i 1}, \hat{y}_{i 2}, \hat{\lambda}_{i 3}$ on the selected sample, is that this procedure places strong assumptions on the reduced form error, $v_{2}$ in equation (19.57). In effect, it adds $-\alpha_{1} v_{2}$ to the structural error, $u_{1}$, with the result that we would have to assume, at a minimum, that $u_{1}-\alpha_{1} v_{2}$ is independent of $\mathbf{z}$ and that $\mathrm{E}\left(u_{1}-\alpha_{1} v_{2} \mid v_{3}\right)$ is linear. Such assumptions are virtually impossible unless $y_{2}$ is continuously distributed. Of course, joint normality of $\left(u_{1}, v_{2}, v_{3}\right)$ is sufficient, but this approach is much more restrictive than assumption 19.2. Thus, applying 2SLS to equation (19.60) is much preferred to first replacing $y_{2}$ with fitted values.

Even if $y_{2}$ is exogenous in the population equation (19.56), when $y_{2}$ is sometimes missing we generally need an instrument for $y_{2}$ when selection is not ignorable (that is, $\mathrm{E}\left(u_{1} \mid \mathbf{z}_{1}, y_{2}, y_{3}\right) \neq \mathrm{E}\left(u_{1}\right)$ ). In example 19.4 we could use family background variables and another test score, such as $K W W$, as IVs for $I Q$, assuming these are always observed. We would generally include all such variables in the reduced-form selection equation. Procedure 19.2 works whether we assume $I Q$ is a proxy variable for ability or an indicator of ability (see Chapters 4 and 5).

To illustrate the likely consequences of using the same variable as an instrument for $y_{2}$ and to predict selection, reconsider example 19.7. As we already know, the number of young children predicts a lower probability of labor force participation. In\\
this data set, education has a positive and statistically significant partial correlation with kidslt6. Therefore, we might try applying procedure 19.2 with kidslto the only element in $\mathbf{z}_{2}=\mathbf{z}$ that is not in $\mathbf{z}_{1}$. That is, the probit for inlf includes exper, exper ${ }^{2}$, and kidslt6, and then the the instrumental variables are exper, exper ${ }^{2}, \hat{\lambda}_{3}$ and kidslt6. Using this approach, the estimated education coefficient is .329 ( $s e=.188$ ), which is much too large to be plausible. Further, the $95 \%$ confidence interval (which ignores the two-step estimation) includes zero, showing the effects of the severe collinearity in the IV estimation.

If we make stronger assumptions, it is possible to estimate model (19.56)-(19.58) by partial maximum likelihood of the kind discussed in problem 13.7. One possibility is to assume that $\left(u_{1}, v_{2}, v_{3}\right)$ is trivariate normal and independent of $\mathbf{z}$. In addition to ruling out discrete $y_{2}$, such a procedure would be computationally difficult. If $y_{2}$ is binary, we can model it as $y_{2}=1\left[\mathbf{z}_{2} \boldsymbol{\delta}_{2}+v_{2}>0\right]$, where $v_{2} \mid \mathbf{z} \sim \operatorname{Normal}(0,1)$. But maximum likelihood estimation that allows any correlation matrix for $\left(u_{1}, v_{2}, v_{3}\right)$ is complicated and less robust than procedure 19.2.

\subsection*{19.6.3 Binary Response Model with Sample Selection}
We can estimate binary response models with sample selection if we assume that the latent errors are bivariate normal and independent of the explanatory variables. Write the model as


\begin{equation*}
y_{1}=1\left[\mathbf{x}_{1} \boldsymbol{\beta}_{1}+u_{1}>0\right] \tag{19.61}
\end{equation*}


$y_{2}=1\left[\mathbf{x} \boldsymbol{\delta}_{2}+v_{2}>0\right]$,\\
where the second equation is the sample selection equation and $y_{1}$ is observed only when $y_{2}=1$; we assume that $\mathbf{x}$ is always observed. For example, suppose $y_{1}$ is an employment indicator and $\mathbf{x}_{1}$ contains a job training binary indicator (which we assume is exogenous), as well as other human capital and family background variables. We might lose track of some people who are eligible to participate in the program; this is an example of sample attrition. If attrition is systematically related to $u_{1}$, estimating equation (19.61) on the sample at hand can result in an inconsistent estimator of $\boldsymbol{\beta}_{1}$.

If we assume that $\left(u_{1}, v_{2}\right)$ is independent of $\mathbf{x}$ with a zero-mean normal distribution (and unit variances), we can apply partial maximum likelihood. What we need is the density of $y_{1}$ conditional on $\mathbf{x}$ and $y_{2}=1$. We have essentially found this density in Chapter 15: in equation (15.55) set $\alpha_{1}=0$, replace $\mathbf{z}$ with $\mathbf{x}$, and replace $\boldsymbol{\delta}_{1}$ with $\boldsymbol{\beta}_{1}$. The parameter $\rho_{1}$ is still the correlation between $u_{1}$ and $v_{2}$. A two-step procedure can be applied: first, estimate $\boldsymbol{\delta}_{2}$ by probit of $y_{2}$ on $\mathbf{x}$. Then, estimate $\boldsymbol{\beta}_{1}$ and $\rho_{1}$ in the second stage using equation (15.55) along with $\mathrm{P}\left(y_{1}=0 \mid \mathbf{x}, y_{2}=1\right)$.

A convincing analysis requires at least one variable in $\mathbf{x}$-that is, something that determines selection-that is not also in $\mathbf{x}_{1}$. Otherwise, identification is off of the nonlinearities in the probit models.

Importantly, some simple strategies for "correcting" for sample selection are not valid. For example, on one hand, it is tempting to estimate the selection equation by probit and then plug the estimated inverse Mills ratio into the second-stage probit, using only the observations with $y_{i 2}=1$. There is no way to justify this as a sample selection correction. On the other hand, inserting the IMR into the second-stage probit is a legitimate test of the null hypothesis of no selection bias. One can start from equation (15.55) with $\alpha_{1}=0, \mathbf{z}=\mathbf{x}, \mathbf{z}_{1}=\mathbf{x}_{1}$, and $\boldsymbol{\delta}_{1}=\boldsymbol{\beta}_{1}$, which is $\mathrm{E}\left(y_{1} \mid \mathbf{x}\right.$, $y_{2}=1$ ). Fixing $\boldsymbol{\delta}_{2}$-because we subsequently insert its probit estimate-we can compute the gradient of the mean function with respect to $\left(\boldsymbol{\beta}_{1}^{\prime}, \rho_{1}\right)^{\prime}$, and then insert $\rho_{1}=0$ (which holds under the null). If $m\left(\mathbf{x}, \boldsymbol{\beta}_{1}, \rho_{1} ; \boldsymbol{\delta}_{2}\right)$ denotes the mean function, then it can be shown that $\nabla_{\boldsymbol{\beta}_{1}} m\left(\mathbf{x}, \boldsymbol{\beta}_{1}, 0 ; \boldsymbol{\delta}_{2}\right)=\phi\left(\mathbf{x}_{1} \boldsymbol{\beta}_{1}\right) \mathbf{x}_{1}$ and $\nabla_{\rho_{1}} m\left(\mathbf{x}, \boldsymbol{\beta}_{1}, 0 ; \boldsymbol{\delta}_{2}\right)= \phi\left(\mathbf{x}_{1} \boldsymbol{\beta}_{1}\right) \lambda\left(\mathbf{x} \boldsymbol{\delta}_{2}\right)$, where $\lambda(\cdot)$ is the IMR. If we use these gradients in the score tests developed in Section 15.5.3, we obtain a simple score statistic for sample selection. Actually, the variable addition version of the test is simpler. In the first step, estimate $\boldsymbol{\delta}_{2}$ by probit of $y_{i 2}$ on $\mathbf{x}_{i}$, using all of the observations. (Under the null, $\hat{\boldsymbol{\delta}}_{2}$ is the MLE because there is no sample selection.) Construct the IMRs, $\hat{\lambda}_{i 2}=\lambda\left(\mathbf{x}_{i} \hat{\boldsymbol{\delta}}_{2}\right)$. Next, using the observations for which $y_{i 2}=1$ (that is, for which $y_{i 1}$ is observed), run probit of $y_{i 1}$ on $\mathbf{x}_{i 1}, \hat{\lambda}_{i 2}$ and use the usual $t$ statistic on $\hat{\lambda}_{i 2}$ to test the null hypothesis $H_{0}: \rho_{1}=0$. Because the coefficient on $\hat{\lambda}_{i 2}$ is zero under the null, there is no need to adjust the $t$ statistic for the first-stage estimation; see Section 12.4.2.

Allowing for endogenous explanatory variables in equation (19.61) along with sample selection is difficult, and it is a useful area of future research.

\subsection*{19.6.4 An Exponential Response Function}
Another nonlinear model for which it is easy to obtain a simple test for selection bias, and relatively easy to correct for sample selection bias, is the exponential response model. The mechanics are very similar to those in Section 18.5 for an exponential response function and a binary endogenous explanatory variable. In particular, we still use equation (18.48) but only for the subsample with $y_{i 2}=1$ (the selected sample). The two-step method proposed by Terza (1998), which consists of estimating a probit in the first stage followed by nonlinear regression (or a quasi-MLE, such as the Poisson or exponential), consistently estimates the parameters. A simple test of sample selection bias is obtained by adding the log of the inverse Mills ratio, $\log \left[\lambda\left(\mathbf{z}_{i} \hat{\boldsymbol{\delta}}_{2}\right)\right]$, to the exponential function, and estimating the resulting "model" by, say, the Poisson QMLE using the selected sample. The robust $t$ statistic for $\log \left[\lambda\left(\mathbf{z}_{i} \hat{\boldsymbol{\delta}}_{2}\right)\right]$ that allows the likelihood to be misspecified is a valid test of the null hypothesis of no selection bias.

\subsection*{19.7 Incidental Truncation: A Tobit Selection Equation}
We now study the case where more information is available on sample selection, primarily in the context of incidental truncation. In particular, we assume that selection is based on the outcome of a Tobit, rather than a probit, equation. The model in Section 19.7.1 is a special case of the model studied by Vella (1992) in the context of testing for selectivity bias.

\subsection*{19.7.1 Exogenous Explanatory Variables}
We now consider the case where the selection equation is of the censored Tobit form. The population model is


\begin{equation*}
y_{1}=\mathbf{x}_{1} \boldsymbol{\beta}_{1}+u_{1} \tag{19.63}
\end{equation*}


$y_{2}=\max \left(0, \mathbf{x} \boldsymbol{\delta}_{2}+v_{2}\right)$,\\
where $\left(\mathbf{x}, y_{2}\right)$ is always observed in the population but $y_{1}$ is observed only when $y_{2}>0$. A standard example occurs when $y_{1}$ is the log of the hourly wage offer and $y_{2}$ is weekly or annual hours of labor supply.

ASSUMPTION 19.3: (a) $\left(\mathbf{x}, y_{2}\right)$ is always observed in the population, but $y_{1}$ is observed only when $y_{2}>0$; (b) $\left(u_{1}, v_{2}\right)$ is independent of $\mathbf{x}$; (c) $v_{2} \sim \operatorname{Normal}\left(0, \tau_{2}^{2}\right)$; and (d) $\mathrm{E}\left(u_{1} \mid v_{2}\right)=\gamma_{1} v_{2}$.

These assumptions are very similar to the assumptions for a probit selection equation. The only difference is that $v_{2}$ now has an unknown variance, since $y_{2}$ is a censored as opposed to binary variable.

Amemiya (1985) calls equations (19.63) and (19.64) the Type III Tobit model, but we emphasize that equation (19.63) is the structural population equation of interest and that equation (19.64) simply determines when $y_{1}$ is observed. In the labor economics example, we are interested in the wage offer equation, and equation (19.64) is a reduced-form hours equation. As with a probit selection equation, it makes no sense to define $y_{1}$ to be, say, zero, just because we do not observe $y_{1}$.

The starting point is equation (19.52), just as in the probit selection case. Now define the selection indicator as $s_{2}=1$ if $y_{2}>0$, and $s_{2}=0$ otherwise. Since $s_{2}$ is a function of $\mathbf{x}$ and $v_{2}$, it follows immediately that\\
$\mathrm{E}\left(y_{1} \mid \mathbf{x}, v_{2}, s_{2}\right)=\mathbf{x}_{1} \boldsymbol{\beta}_{1}+\gamma_{1} v_{2}$.

This equation means that, if we could observe $v_{2}$, then an OLS regression of $y_{1}$ on $\mathbf{x}_{1}$, $v_{2}$ using the selected subsample would consistently estimate $\left(\boldsymbol{\beta} \boldsymbol{\beta}_{1}, \boldsymbol{\gamma}_{1}\right)$, as we discussed in Section 19.4.1. While $v_{2}$ cannot be observed when $y_{2}=0$ (because when $y_{2}=0$,\\
we only know that $v_{2} \leq-\mathbf{x} \boldsymbol{\delta}_{2}$ ), for $y_{2}>0, v_{2}=y_{2}-\mathbf{x} \boldsymbol{\delta}_{2}$. Thus, if we knew $\boldsymbol{\delta}_{2}$, we would know $v_{2}$ whenever $y_{2}>0$. It seems reasonable that, because $\boldsymbol{\delta}_{2}$ can be consistently estimated by Tobit on the whole sample, we can replace $v_{2}$ with consistent estimates.

Procedure 19.3: (a) Estimate equation (19.64) by standard Tobit using all $N$ observations. For $y_{i 2}>0$ (say $i=1,2, \ldots, N_{1}$ ), define\\
$\hat{v}_{i 2}=y_{i 2}-\mathbf{x}_{i} \hat{\boldsymbol{\delta}}_{2}$.\\
(b) Using observations for which $y_{i 2}>0$, estimate $\boldsymbol{\beta}_{1}, \gamma_{1}$ by the OLS regression\\
$y_{i 1}$ on $\mathbf{x}_{i 1}, \hat{v}_{i 2} \quad i=1,2, \ldots, N_{1}$.

This regression produces consistent, $\sqrt{N}$-asymptotically normal estimators of $\boldsymbol{\beta}_{1}$ and $\gamma_{1}$ under assumption 19.3.

The statistic to test for selectivity bias is just the usual $t$ statistic on $\hat{v}_{i 2}$ in regression (19.67), which was suggested by Vella (1992).

It seems likely that there is an efficiency gain over procedure 19.1. If $v_{2}$ were known and we could use regression (19.67) for the entire population, there would definitely be an efficiency gain: the error variance is reduced by conditioning on $v_{2}$ along with $\mathbf{x}$, and there would be no heteroskedasticity in the population. See problem 4.5.

Unlike in the probit selection case, $\mathbf{x}_{1}=\mathbf{x}$ causes no problems here: $v_{2}$ always has separate variation from $\mathbf{x}_{1}$ because of variation in $y_{2}$. We do not need to rely on the nonlinearity of the inverse Mills ratio.

Example 19.8 (Wage Offer Equation for Married Women): We now apply procedure 19.3 to the wage offer equation for married women in example 19.6. (We assume education is exogenous.) The only difference is that the first-step estimation is Tobit, rather than probit, and we include the Tobit residuals as the additional explanatory variables, not the inverse Mills ratio. In regression (19.67), the coefficient on $\hat{v}_{2}$ is -.000053 ( $\mathrm{se}=.000041$ ), which is somewhat more evidence of a sample selection problem, but we still do not reject the null hypothesis $\mathrm{H}_{0}: \gamma_{1}=0$ at even the 15 percent level against a two-sided alternative. Further, the coefficient on educ is .103 (se $=.015$ ), which is not much different from the OLS and Heckit estimates. (Again, we use the usual OLS standard error.) When we include all exogenous variables in the wage offer equation, the estimates from procedure 19.3 are much more stable than the Heckit estimates. For example, the coefficient on educ becomes .093 $(s e=.016)$, which is comparable to the OLS estimates discussed in example 19.6.

For partial maximum likelihood estimation, we assume that ( $u_{1}, v_{2}$ ) is jointly normal, and we use the density for $f\left(y_{2} \mid \mathbf{x}\right)$ for the entire sample and the conditional density $f\left(y_{1} \mid \mathbf{x}, y_{2}, s_{2}=1\right)=f\left(y_{1} \mid \mathbf{x}, y_{2}\right)$ for the selected sample. This approach is fairly straightforward because, when $y_{2}>0, y_{1} \mid \mathbf{x}, y_{2} \sim \operatorname{Normal}\left[\mathbf{x}_{1} \boldsymbol{\beta}_{1}+\right. \left.\gamma_{1}\left(y_{2}-\mathbf{x} \boldsymbol{\delta}_{2}\right), \eta_{1}^{2}\right]$, where $\eta_{1}^{2}=\sigma_{1}^{2}-\sigma_{12}^{2} / \tau_{2}^{2}, \sigma_{1}^{2}=\operatorname{Var}\left(u_{1}\right)$, and $\sigma_{12}=\operatorname{Cov}\left(u_{1}, v_{2}\right)$. The log likelihood for observation $i$ is\\
$\ell_{i}(\boldsymbol{\theta})=s_{i 2} \log f\left(y_{i 1} \mid \mathbf{x}_{i}, y_{i 2} ; \boldsymbol{\theta}\right)+\log f\left(y_{i 2} \mid \mathbf{x}_{i} ; \boldsymbol{\delta}_{2}, \tau_{2}^{2}\right)$,\\
where $f\left(y_{i 1} \mid \mathbf{x}_{i}, y_{i 2} ; \boldsymbol{\theta}\right)$ is the $\operatorname{Normal}\left[\mathbf{x}_{i 1} \boldsymbol{\beta}_{1}+\gamma_{1}\left(y_{i 2}-\mathbf{x}_{i} \boldsymbol{\delta}_{2}\right), \eta_{1}^{2}\right]$ distribution, evaluated at $y_{i 1}$, and $f\left(y_{i 2} \mid \mathbf{x}_{i} ; \boldsymbol{\delta}_{2}, \tau_{2}^{2}\right)$ is the standard censored Tobit density (see equation (17.19)). As shown in problem 13.7, the usual MLE theory can be used even though the log-likelihood function is not based on a full conditional density.

It is possible to obtain sample selection corrections and tests for various other nonlinear models when the selection rule is of the Tobit form. For example, suppose that the binary variable $y_{1}$ given $\mathbf{z}$ follows a probit model, but it is observed only when $y_{2}>0$. A valid test for selection bias is to include the Tobit residuals, $\hat{v}_{2}$, in a probit of $y_{1}$ on $\mathbf{z}, \hat{v}_{2}$ using the selected sample; see Vella (1992). This procedure also produces consistent estimates (up to scale), as can be seen by applying the maximum likelihood results in Section 19.4.2 along with two-step estimation results.

Honoré, Kyriazidou, and Udry (1997) show how to estimate the parameters of the Type III Tobit model without making distributional assumptions.

\subsection*{19.7.2 Endogenous Explanatory Variables}
We explicitly consider the case of a single endogenous explanatory variable, as in Section 19.6.2. We use equations (19.56) and (19.57), and, in place of equation (19.58), we have a Tobit selection equation:\\
$y_{3}=\max \left(0, \mathbf{z} \boldsymbol{\delta}_{3}+v_{3}\right)$.

ASSUMPTION 19.4: (a) $\left(\mathbf{z}, y_{3}\right)$ is always observed, $\left(y_{1}, y_{2}\right)$ is observed when $y_{3}>0$; (b) $\left(u_{1}, v_{3}\right)$ is independent of $\mathbf{z}$; (c) $v_{3} \sim \operatorname{Normal}\left(0, \tau_{3}^{2}\right)$; (d) $\mathrm{E}\left(u_{1} \mid v_{3}\right)=\gamma_{1} v_{3}$; and (e) $\mathrm{E}\left(\mathbf{z}^{\prime} v_{2}\right)=\mathbf{0}$ and, writing $\mathbf{z} \boldsymbol{\delta}_{2}=\mathbf{z}_{1} \boldsymbol{\delta}_{21}+\mathbf{z}_{2} \boldsymbol{\delta}_{22}, \boldsymbol{\delta}_{22} \neq \mathbf{0}$.

Again, these assumptions are very similar to those used with a probit selection mechanism.

To derive an estimating equation, write\\
$y_{1}=\mathbf{z}_{1} \delta_{1}+\alpha_{1} y_{2}+\gamma_{1} v_{3}+e_{1}$,\\
where $e_{1} \equiv u_{1}-\mathrm{E}\left(u_{1} \mid v_{3}\right)$. Since ( $e_{1}, v_{3}$ ) is independent of $\mathbf{z}$ by assumption 19.4b, $\mathrm{E}\left(e_{1} \mid \mathbf{z}, v_{3}\right)=0$. From Theorem 19.1, if $v_{3}$ were observed, we could estimate equation (19.70) by 2SLS on the selected sample using instruments ( $\mathbf{z}, v_{3}$ ). As before, we can estimate $v_{3}$ when $y_{3}>0$, since $\boldsymbol{\delta}_{3}$ can be consistently estimated by Tobit of $y_{3}$ on $\mathbf{z}$ (using the entire sample).

Procedure 19.4: (a) Obtain $\hat{\boldsymbol{\delta}}_{3}$ from Tobit of $y_{3}$ on $\mathbf{z}$ using all observations. Obtain the Tobit residuals $\hat{v}_{i 3}=y_{i 3}-\mathbf{z}_{i} \hat{\boldsymbol{\delta}}_{3}$ for $y_{i 3}>0$.\\
(b) Using the selected subsample, estimate the equation\\
$y_{i 1}=\mathbf{z}_{i 1} \boldsymbol{\delta}_{1}+\alpha_{1} y_{i 2}+\gamma_{1} \hat{v}_{i 3}+$ error $_{i}$\\
by 2SLS, using instruments $\left(\mathbf{z}_{i}, \hat{v}_{i 3}\right)$. The estimators are $\sqrt{N}$-consistent and asymptotically normal under assumption 19.4.

Comments similar to those after procedure 19.2 hold here as well. Strictly speaking, identification really requires that $\mathbf{z}_{2}$ appear in the linear projection of $y_{2}$ onto $\mathbf{z}_{1}$, $\mathbf{z}_{2}$, and $v_{3}$ in the selected subpopulation. The null of no selection bias is tested using the 2SLS $t$ statistic (or maybe its heteroskedasticity-robust version) on $\hat{v}_{i 3}$. When $\gamma_{1} \neq 0$, standard errors should be corrected using two-step methods.

As in the case with a probit selection equation, the endogenous variable $y_{2}$ can be continuous, discrete, a corner solution, and so on. Extending the method to multiple endogenous explanatory variables is straightforward. The only restriction is the usual one for linear models: we need enough instruments to identify the structural equation. See problem 19.9 for an application to the Mroz data.

An interesting special case of model (19.56), (19.57), and (19.69) occurs when $y_{2}=y_{3}$. Actually, because we only use observations for which $y_{3}>0, y_{2}=y_{3}^{*}$ is also allowed, where $y_{3}^{*}=\mathbf{z} \boldsymbol{\delta}_{3}+v_{3}$. Either way, the variable that determines selection also appears in the structural equation. This special case could be useful when sample selection is caused by a corner solution outcome on $y_{3}$ (in which case $y_{2}=y_{3}$ is natural) or because $y_{3}^{*}$ is subject to data censoring (in which case $y_{2}=y_{3}^{*}$ is more realistic). An example of the former occurs when $y_{3}$ is hours worked and we assume hours appears in the wage offer function. As a data-censoring example, suppose that $y_{1}$ is a measure of growth in an infant's weight starting from birth and that we observe $y_{1}$ only if the infant is brought into a clinic within three months. Naturally, birth weight depends on age, and so $y_{3}^{*}$-length of time between the first and second measurements, which has quantitiative meaning-appears as an explanatory variable in the equation for $y_{1}$. We have a data-censoring problem for $y_{3}^{*}$, which causes a sample selection problem for $y_{1}$. In this case, we would estimate a censored normal regres-\\
sion model for $y_{3}$ [or, possibly, $\left.\log \left(y_{3}\right)\right]$ to account for the data censoring. We would include the residuals $\hat{v}_{i 3}=y_{i 3}-\mathbf{z}_{i} \hat{\boldsymbol{\delta}}_{3}$ in equation (19.71) for the noncensored observations. As our extra instrument we might use distance from the child's home to the clinic.

\subsection*{19.7.3 Estimating Structural Tobit Equations with Sample Selection}
We briefly show how a structural Tobit model can be estimated using the methods of the previous section. As an example, consider the structural labor supply model


\begin{equation*}
\log \left(w^{\mathrm{o}}\right)=\mathbf{z}_{1} \boldsymbol{\beta}_{1}+u_{1} \tag{19.72}
\end{equation*}


$h=\max \left[0, \mathbf{z}_{2} \boldsymbol{\beta}_{2}+\alpha_{2} \log \left(w^{\mathrm{o}}\right)+u_{2}\right]$.

This system involves simultaneity and sample selection because we observe $w^{\mathrm{o}}$ only if $h>0$.

The general form of the model is


\begin{equation*}
y_{1}=\mathbf{z}_{1} \boldsymbol{\beta}_{1}+u_{1} \tag{19.74}
\end{equation*}


$y_{2}=\max \left(0, \mathbf{z}_{2} \boldsymbol{\beta}_{2}+\alpha_{2} y_{1}+u_{2}\right)$.

ASSUMPTION 19.5: (a) $\left(\mathbf{z}, y_{2}\right)$ is always observed; $y_{1}$ is observed when $y_{2}>0$; (b) ( $u_{1}, u_{2}$ ) is independent of $\mathbf{z}$ with a zero-mean bivariate normal distribution; and (c) $\mathbf{z}_{1}$ contains at least one element whose coefficient is different from zero that is not in $\mathbf{z}_{2}$.

As always, it is important to see that equations (19.74) and (19.75) constitute a model describing a population. If $y_{1}$ were always observed, then equation (19.74) could be estimated by OLS. If, in addition, $u_{1}$ and $u_{2}$ were uncorrelated, equation (19.75) could be estimated by type I Tobit. Correlation between $u_{1}$ and $u_{2}$ could be handled by the methods of Section 17.5.2. Now, we require new methods, whether or not $u_{1}$ and $u_{2}$ are uncorrelated, because $y_{1}$ is not observed when $y_{2}=0$.

The restriction in assumption 19.5c is needed to identify the structural parameters $\left(\boldsymbol{\beta}_{2}, \alpha_{2}\right)\left(\boldsymbol{\beta}_{1}\right.$ is always identified). To see that this condition is needed, and for finding the reduced form for $y_{2}$, it is useful to introduce the latent variable


\begin{equation*}
y_{2}^{*} \equiv \mathbf{z}_{2} \boldsymbol{\beta}_{2}+\alpha_{2} y_{1}+u_{2} \tag{19.76}
\end{equation*}


so that $y_{2}=\max \left(0, y_{2}^{*}\right)$. If equations (19.74) and (19.76) make up the system of interest-that is, if $y_{1}$ and $y_{2}^{*}$ are always observed-then $\boldsymbol{\beta}_{1}$ is identified without further restrictions, but identification of $\alpha_{2}$ and $\boldsymbol{\beta}_{2}$ requires exactly assumption 19.5c. This turns out to be sufficient even when $y_{2}$ follows a Tobit model and we have nonrandom sample selection.

The reduced form for $y_{2}^{*}$ is $y_{2}^{*}=\mathbf{z} \boldsymbol{\delta}_{2}+v_{2}$. Therefore, we can write the reduced form of equation (19.75) as\\
$y_{2}=\max \left(0, \mathbf{z} \boldsymbol{\delta}_{2}+v_{2}\right)$.

But then equations (19.74) and (19.77) constitute the model we studied in Section 17.6.1. The vector $\boldsymbol{\delta}_{2}$ is consistently estimated by Tobit, and $\boldsymbol{\beta}_{1}$ is estimated as in procedure 19.3. The only remaining issue is how to estimate the structural parameters of equation (19.75), $\alpha_{2}$ and $\boldsymbol{\beta}_{2}$. In the labor supply case, these are the labor supply parameters.

Assuming identification, estimation of ( $\alpha_{2}, \boldsymbol{\beta}_{2}$ ) is fairly straightforward after having estimated $\boldsymbol{\beta}_{1}$. To see this point, write the reduced form of $y_{2}$ in terms of the structural parameters as\\
$y_{2}=\max \left[0, \mathbf{z}_{2} \boldsymbol{\beta}_{2}+\alpha_{2}\left(\mathbf{z}_{1} \boldsymbol{\beta}_{1}\right)+v_{2}\right]$.

Under joint normality of $u_{1}$ and $u_{2}, v_{2}$ is normally distributed. Therefore, if $\boldsymbol{\beta}_{1}$ were known, $\boldsymbol{\beta}_{2}$ and $\alpha_{2}$ could be estimated by standard Tobit using $\mathbf{z}_{2}$ and $\mathbf{z}_{1} \boldsymbol{\beta}_{1}$ as regressors. Operationalizing this procedure requires replacing $\boldsymbol{\beta}_{1}$ with its consistent estimator. Thus, using all observations, $\boldsymbol{\beta}_{2}$ and $\alpha_{2}$ are estimated from the Tobit equation\\
$y_{i 2}=\max \left[0, \mathbf{z}_{i 2} \boldsymbol{\beta}_{2}+\alpha_{2}\left(\mathbf{z}_{i 1} \hat{\boldsymbol{\beta}}_{1}\right)+\right.$ error $\left._{i}\right]$.

To summarize, we have the following:\\
Procedure 19.5: (a) Use procedure 19.3 to obtain $\hat{\boldsymbol{\beta}}_{1}$.\\
(b) Obtain $\hat{\boldsymbol{\beta}}_{2}$ and $\hat{\alpha}_{2}$ from the Tobit in equation (19.79).

In applying this procedure, it is important to note that the explanatory variable in equation (19.79) is $\mathbf{z}_{i 1} \hat{\boldsymbol{\beta}}_{1}$ for all $i$. These are not the fitted values from regression (19.67), which depend on $\hat{v}_{i 2}$. Also, it may be tempting to use $y_{i 1}$ in place of $\mathbf{z}_{i 1} \hat{\boldsymbol{\beta}}_{1}$ for that part of the sample for which $y_{i 1}$ is observed. This approach is not a good idea: the estimators are inconsistent in this case.

The estimation in equation (19.79) makes it clear that the procedure fails if $\mathbf{z}_{1}$ does not contain at least one variable not in $\mathbf{z}_{2}$. If $\mathbf{z}_{1}$ is a subset of $\mathbf{z}_{2}$, then $\mathbf{z}_{i 1} \hat{\boldsymbol{\beta}}_{1}$ is a linear combination of $\mathbf{z}_{i 2}$, and so perfect multicollinearity will exist in equation (19.79).

Estimating $\operatorname{Avar}\left(\hat{\alpha}_{2}, \hat{\boldsymbol{\beta}}_{2}\right)$ is even messier than estimating $\operatorname{Avar}\left(\hat{\boldsymbol{\beta}}_{1}\right)$, since $\left(\hat{\alpha}_{2}, \hat{\boldsymbol{\beta}}_{2}\right)$ comes from a three-step procedure. Often just the usual Tobit standard errors and test statistics reported from equation (19.79) are used, even though these are not strictly valid. By setting the problem up as a large GMM problem, as illustrated in Chapter 14, correct standard errors and test statistics can be obtained. Bootstrapping can be used, too.

Under assumption 19.5, a full maximum likelihood approach is possible. In fact, the log-likelihood function can be constructed from equations (19.74) and (19.78), and it has a form very similar to equation (19.68). The only difference is that nonlinear restrictions are imposed automatically on the structural parameters. In addition to making it easy to obtain valid standard errors, MLE is desirable because it allows us to estimate $\sigma_{2}^{2}=\operatorname{Var}\left(u_{2}\right)$, which is needed to estimate average partial effects in equation (19.75).

In examples such as labor supply, it is not clear where the elements of $\mathbf{z}_{1}$ that are not in $\mathbf{z}_{2}$ might come from. One possibility is a union binary variable, if we believe that union membership increases wages (other factors accounted for) but has no effect on labor supply once wage and other factors have been controlled for. This approach would require knowing union status for people whether or not they are working in the period covered by the survey. In some studies past experience is assumed to affect wage-which it certainly does-and is assumed not to appear in the labor supply function, a tenuous assumption.

\subsection*{19.8 Inverse Probability Weighting for Missing Data}
We now turn to a different method for correcting for general missing data problems, inverse probability weighting (IPW). Compared with Heckman-type approaches, IPW applies very generally to any estimation problem that involves minimization or maximization. However, the assumptions under which IPW produces consistent estimators of the population parameters are quite different from those used in Heckman-type methods. We highlight the differences in this section, which follows the M-estimation setup in Chapter 12 and in Wooldridge (2007).

Again, we characterize missing data using a binary selection indicator, $s_{i}$. Therefore, a random draw from the population consists of ( $\mathbf{w}_{i}, s_{i}$ ), and all or part of $\mathbf{w}_{i}$ is not observed if $s_{i}=0$. We are interested in estimating $\boldsymbol{\theta}_{o}$, the solution to the population problem


\begin{equation*}
\min _{\boldsymbol{\theta} \in \Theta} \mathrm{E}\left[q\left(\mathbf{w}_{i}, \boldsymbol{\theta}\right)\right], \tag{19.80}
\end{equation*}


where $q(\mathbf{w}, \cdot)$ is the objective function for given $\mathbf{w}$. If we use the selected sample to estimate $\boldsymbol{\theta}_{o}$ we solve


\begin{equation*}
\min _{\boldsymbol{\theta} \in \Theta} N^{-1} \sum_{i=1}^{N} s_{i} q\left(\mathbf{w}_{i}, \boldsymbol{\theta}\right) . \tag{19.81}
\end{equation*}


We call the solution to this problem the unweighted M-estimator, $\hat{\boldsymbol{\theta}}_{u}$, to distinguish it from the weighted estimator that will be introduced later. We have already seen examples, particularly for the linear model, where $\hat{\boldsymbol{\theta}}_{u}$ is not consistent for $\boldsymbol{\theta}_{o}$.

Inverse probability weighting is a general approach to nonrandom sampling that dates back to Horvitz and Thompson (1952). IPW has been used more recently for regression models with missing data (for example, Robins, Rotnitzky, and Zhou 1995) and in the treatment effects literature (Hirano, Imbens, and Ridder (2003) and Chapter 21). The key is that we have some variables that are "good" predictors of selection, something we make precise in the following assumption.

ASSUMPTION 19.6: (a) The vector $\mathbf{w}_{i}$ is observed whenever $s_{i}=1$.\\
(b) There is a random vector $\mathbf{z}_{i}$ such that $\mathrm{P}\left(s_{i}=1 \mid \mathbf{w}_{i}, \mathbf{z}_{i}\right)=\mathrm{P}\left(s_{i}=1 \mid \mathbf{z}_{i}\right) \equiv p\left(\mathbf{z}_{i}\right)$.\\
(c) For all $\mathbf{z} \in \mathscr{Z} \subset \mathbb{R}^{J}, p(\mathbf{z})>0$.\\
(d) $\mathbf{z}_{i}$ is observed when $s_{i}=1$.

When $\mathbf{z}_{i}$ is always observed, assumption 19.6 is essentially the missing at random assumption mentioned in Section 19.4. Wooldridge (2007) shows that allowing $\mathbf{z}_{i}$ to be unobserved when $s_{i}=0$ allows coverage of certain stratified sampling schemes, which we treat in Chapter 20, and IPW solutions to censored duration data, which we cover in Chapter 22. One also sees assumption 19.6b described as ignorable selection (conditional on $\mathbf{z}_{i}$ ). In the treatment effects literature-see Chapter 21assumption 19.6b is essentially the unconfoundedness assumption.

Assumption 19.6 encompasses the selection on observables assumption mentioned in Section 19.4. In the general M-estimation framework, selection on observables typically applies when $\mathbf{w}_{i}$ partitions as $\left(\mathbf{x}_{i}, \mathbf{y}_{i}\right), \mathbf{x}_{i}$ is always observed but $\mathbf{y}_{i}$ is not, and $\mathbf{z}_{i}$ is a vector that is always observed and includes $\mathbf{x}_{i}$. Then, $s_{i}$ is allowed to be a function of observables $\mathbf{z}_{i}$, but $s_{i}$ cannot be related to unobserved factors affecting $\mathbf{y}_{i}$. Assumption 19.6 does not apply to the selection on unobservables case, at least as that terminology has been used in econometrics. Selection on unobservables is essentially what we covered in Sections 19.5 and 19.6, where we explicitly allowed selection to be correlated with unobservables after conditioning on exogenous variables. Heckman's (1976) solution to the incidental truncation problem, which applies to linear models and a few others, requires at least one exogenous variable that affects selection but does not have a partial effect on the outcome. In assumption 19.6, the $\mathbf{z}_{i}$ should have properties different from exogenous variables that are used for identification in the Heckman approach. In effect, the $\mathbf{z}_{i}$ should be good proxies for the unobservables that affect $\mathbf{y}_{i}$ and also determine selection. Sometimes, $\mathbf{z}_{i}$ includes outcomes on $\mathbf{y}_{i}$ and even $\mathbf{x}_{i}$ from previous time periods.

As an example, consider the linear model for a random draw, $y_{i}=\mathbf{x}_{i} \boldsymbol{\beta}_{o}+u_{i}$, $\mathrm{E}\left(\mathbf{x}_{i}^{\prime} u_{i}\right)=\mathbf{0}$, and suppose data are missing on $y_{i}$ when $s_{i}=0$. Let $\mathbf{z}_{i}$ be a vector of variables that includes $\mathbf{x}_{i}$, and, for concreteness, assume that $s_{i}=1\left[\mathbf{z}_{i} \boldsymbol{\delta}_{o}+v_{i} \geq 0\right]$, where $v_{i}$ is independent of $\mathbf{z}_{i}$. To use a Heckman selection correction we would assume $\mathrm{E}\left(\mathbf{z}_{i}^{\prime} u_{i}\right)=\mathbf{0}$ but allow arbitrary correlation between $v_{i}$ and $u_{i}$, so that selection is correlated with $u_{i}$ even after netting out $\mathbf{z}_{i}$. By contrast, assumption 19.6b is the same as $\mathrm{P}\left(s_{i}=1 \mid \mathbf{z}_{i}, u_{i}\right)=\mathrm{P}\left(s_{i}=1 \mid \mathbf{z}_{i}\right)$, which essentially means that $u_{i}$ and $v_{i}$ should be independent. Thus, under assumption 19.6b, we assume that $\mathbf{z}_{i}$ is such a good predictor of selection that, conditional on $\mathbf{z}_{i}, s_{i}$ is independent of the unobservables in the regression equation, $u_{i}$. However, $\mathbf{z}_{i}$ and $u_{i}$ can be arbitrarily correlated.

If we could observe the selection probabilities then, under assumption 19.6, solving the missing data problem would be easy. To see why, let $g(\mathbf{w})$ be any scalar function such that the mean, $\mu=\mathrm{E}\left[g\left(\mathbf{w}_{i}\right)\right]$, exists. Then, using iterated expectations,


\begin{align*}
\mathrm{E}\left[s_{i} g\left(\mathbf{w}_{i}\right) / p\left(\mathbf{z}_{i}\right)\right] & =\mathrm{E}\left\{\mathrm{E}\left[s_{i} g\left(\mathbf{w}_{i}\right) / p\left(\mathbf{z}_{i}\right) \mid \mathbf{w}_{i}, \mathbf{z}_{i}\right]\right\} \\
& =\mathrm{E}\left\{\mathrm{E}\left(s_{i} \mid \mathbf{w}_{i}, \mathbf{z}_{i}\right) g\left(\mathbf{w}_{i}\right) / p\left(\mathbf{z}_{i}\right)\right\} \\
& =\mathrm{E}\left\{\mathrm{P}\left(s_{i}=1 \mid \mathbf{w}_{i}, \mathbf{z}_{i}\right) g\left(\mathbf{w}_{i}\right) / p\left(\mathbf{z}_{i}\right)\right\} \\
& =\mathrm{E}\left\{p\left(\mathbf{z}_{i}\right) g\left(\mathbf{w}_{i}\right) / p\left(\mathbf{z}_{i}\right)\right\}=\mathrm{E}\left[g\left(\mathbf{w}_{i}\right)\right], \tag{19.82}
\end{align*}


where the last equality follows from $\mathrm{P}\left(s_{i}=1 \mid \mathbf{w}_{i}, \mathbf{z}_{i}\right)=p\left(\mathbf{z}_{i}\right)$. This result shows that the population mean of any function of $\mathbf{w}_{i}$ can be recovered by weighting the selected observation by the inverse of the probability of selection. It follows immediately that a consistent (actually, unbiased) estimator of $\mu$ is $N^{-1} \sum_{i=1}^{n}\left[s_{i} g\left(\mathbf{w}_{i}\right) / p\left(\mathbf{z}_{i}\right)\right]$. Actually, a somewhat more common estimator, based on the fact that $\mathrm{E}\left[s_{i} / p\left(\mathbf{z}_{i}\right)\right]=1$, is


\begin{equation*}
\hat{\mu}_{I P W}=\left(\sum_{i=1}^{N}\left[s_{i} / p\left(\mathbf{z}_{i}\right)\right]\right)^{-1}\left(\sum_{i=1}^{N}\left[s_{i} g\left(\mathbf{w}_{i}\right) / p\left(\mathbf{z}_{i}\right)\right]\right) \tag{19.83}
\end{equation*}


which is a weighted average of the sampled data where the weights add to one. While equation (19.83) appears to depend on $N$ (the number of times the population was sampled), $N$ is not needed to compute $\hat{\mu}_{I P W}$. The sampling weights implicit in equation (19.83) are often reported in survey data to obtain means in the presence of missing data.

We can now see how to use IPW estimation in the context of M-estimation. The IPW estimator, $\tilde{\boldsymbol{\theta}}_{w}$, solves


\begin{equation*}
\min _{\theta \in \Theta} N^{-1} \sum_{i=1}^{N}\left[s_{i} / p\left(\mathbf{z}_{i}\right)\right] q\left(\mathbf{w}_{i}, \boldsymbol{\theta}\right) \tag{19.84}
\end{equation*}


From the previous argument, the mean of each summand is $\mathrm{E}\left[q\left(\mathbf{w}_{i}, \boldsymbol{\theta}\right)\right]$, which is minimized at $\boldsymbol{\theta}_{o}$, and so, under the mild conditions of the uniform weak law of large numbers, $\tilde{\boldsymbol{\theta}}_{w}$ is consistent for $\boldsymbol{\theta}_{o}$; see Theorem 12.2.

In most cases where the selection probabilities $p\left(\mathbf{z}_{i}\right)$ are not known, $\mathbf{z}_{i}$ is assumed to be always observed, so that a model for $\mathrm{P}\left(s_{i}=1 \mid \mathbf{z}_{i}\right)$ can be estimated by binary response maximum likelihood. Here we consider the case where we use a binary response model for $\mathrm{P}\left(s_{i}=1 \mid \mathbf{z}_{i}\right)$, which requires that $\mathbf{z}_{i}$ is always observed.

ASSUMPTION 19.7: (a) $G(\mathbf{z}, \gamma)$ is a parametric model for $p(\mathbf{z})$, where $\gamma \in \boldsymbol{\Gamma} \subset \mathbb{R}^{M}$ and $G(\mathbf{z}, \gamma)>0$, all $\mathbf{z}$ and $\gamma$.\\
(b) There there exists $\gamma_{o} \in \boldsymbol{\Gamma}$ such that $p(\mathbf{z})=G\left(\mathbf{z}, \gamma_{o}\right)$.\\
(c) $\hat{\gamma}$ is the binary response maximum likelihood estimator, and the regularity conditions for MLE hold such that


\begin{equation*}
\sqrt{N}\left(\hat{\gamma}-\gamma_{o}\right)=\left\{\mathrm{E}\left[\mathbf{d}_{i}\left(\gamma_{o}\right) \mathbf{d}_{i}\left(\gamma_{o}\right)^{\prime}\right]\right\}^{-1}\left(N^{-1 / 2} \sum_{i=1}^{N} \mathbf{d}_{i}\left(\gamma_{o}\right)\right)+o_{p}(1), \tag{19.85}
\end{equation*}


where $\mathbf{d}_{i}(\gamma) \equiv \nabla_{\gamma} G(\mathbf{z}, \gamma)^{\prime}\left[s_{i}-G(\mathbf{z}, \gamma)\right] /\{G(\mathbf{z}, \gamma)[1-G(\mathbf{z}, \gamma)]\}$ is the $M \times 1$ score vector for the MLE.

Given $\hat{\gamma}$, we can form $G\left(\mathbf{z}_{i}, \hat{\gamma}\right)$ for all $i$ with $s_{i}=1$, and then obtain the weighted $M$-estimator, $\hat{\boldsymbol{\theta}}_{w}$, by solving


\begin{equation*}
\min _{\theta \in \Theta} N^{-1} \sum_{i=1}^{N}\left[s_{i} / G\left(\mathbf{z}_{i}, \hat{\gamma}\right)\right] q\left(\mathbf{w}_{i}, \boldsymbol{\theta}\right) . \tag{19.86}
\end{equation*}


Replacing the unknown probability $p(\mathbf{z})=G\left(\mathbf{z}, \gamma_{o}\right)$ with $G\left(\mathbf{z}_{i}, \hat{\gamma}\right)$ does not affect consistency of $\hat{\boldsymbol{\theta}}_{w}$ under the general conditions for two-step estimators; see Section 12.4.1. More interesting is finding the asymptotic distribution of $\sqrt{N}\left(\hat{\boldsymbol{\theta}}_{w}-\boldsymbol{\theta}_{o}\right)$.

The following result assumes that the objective function $q(\mathbf{w}, \cdot)$ is twice continuously differentiable on the interior of $\boldsymbol{\Theta}$, as in Section 13.10.2. Write $\mathbf{r}\left(\mathbf{w}_{i}, \boldsymbol{\theta}\right) \equiv \nabla_{\boldsymbol{\theta}} q\left(\mathbf{w}_{i}, \boldsymbol{\theta}\right)^{\prime}$ as the $\mathrm{P} \times 1$ score of the unweighted objective function, $H(\mathbf{w}, \boldsymbol{\theta}) \equiv \nabla_{\boldsymbol{\theta}}^{2} q(\mathbf{w}, \boldsymbol{\theta})$ as the $\mathrm{P} \times \mathrm{P}$ Hessian of $q\left(\mathbf{w}_{i}, \boldsymbol{\theta}\right)$, and $k\left(s_{i}, \mathbf{z}_{i}, \mathbf{w}_{i}, \boldsymbol{\gamma}, \boldsymbol{\theta}\right) \equiv\left[s_{i} / G\left(\mathbf{z}_{i}, \boldsymbol{\gamma}\right)\right] r\left(\mathbf{w}_{i}, \boldsymbol{\theta}\right)$ as the selected, weighted score function; in particular, $k\left(s_{i}, \mathbf{z}_{i}, \mathbf{w}_{i}, \boldsymbol{\gamma}, \boldsymbol{\theta}\right)$ is zero whenever $s_{i}=0$.

It is easily shown that the conditions of the "surprising" efficiency result in Section 13.10.2 hold. Therefore,


\begin{equation*}
\sqrt{N}\left(\hat{\boldsymbol{\theta}}_{w}-\boldsymbol{\theta}_{o}\right) \stackrel{a}{\sim} \operatorname{Normal}\left(0, \mathbf{A}_{o}^{-1} \mathbf{D}_{o} \mathbf{A}_{o}^{-1}\right) \tag{19.87}
\end{equation*}


where $\mathbf{A}_{o} \equiv \mathrm{E}\left[H\left(\mathbf{w}_{i}, \boldsymbol{\theta}_{o}\right)\right], \mathbf{D}_{o} \equiv \mathrm{E}\left(\mathbf{e}_{i} \mathbf{e}_{i}^{\prime}\right), \mathbf{e}_{i} \equiv \mathbf{k}_{i}-\mathrm{E}\left(\mathbf{k}_{i} \mathbf{d}_{i}^{\prime}\right)\left[\mathrm{E}\left(\mathbf{d}_{i} \mathbf{d}_{i}^{\prime}\right)\right]^{-1} \mathbf{d}_{i}$, and $\mathbf{k}_{i}$ and $\mathbf{d}_{i}$ are evaluated at $\left(\boldsymbol{\gamma}_{o}, \boldsymbol{\theta}_{o}\right)$ and $\boldsymbol{\gamma}_{o}$, respectively. Further, consistent estimators of $\mathbf{A}_{o}$ and $\mathbf{D}_{o}$, respectively, are


\begin{equation*}
\hat{\mathbf{A}} \equiv N^{-1} \sum_{i=1}^{N}\left[s_{i} / G\left(\mathbf{z}_{i}, \hat{\gamma}\right)\right] H\left(\mathbf{w}_{i}, \hat{\boldsymbol{\theta}}_{w}\right) \tag{19.88}
\end{equation*}


and\\
$\hat{\mathbf{D}} \equiv N^{-1} \sum_{i=1}^{N} \hat{\mathbf{e}}_{i} \hat{\mathbf{e}}_{i}^{\prime}$,\\
where the $\hat{\mathbf{e}}_{i} \equiv \hat{\mathbf{k}}_{i}-\left(N^{-1} \sum_{i=1}^{N} \hat{\mathbf{k}}_{i} \hat{\mathbf{d}}_{i}^{\prime}\right)\left(N^{-1} \sum_{i=1}^{N} \hat{\mathbf{d}}_{i} \hat{\mathbf{d}}_{i}^{\prime}\right)^{-1} \hat{\mathbf{d}}_{i}$ are the $\mathrm{P} \times 1$ residuals from the multivariate regression of $\hat{\mathbf{k}}_{i}$ on $\hat{\mathbf{d}}, i=1, \ldots, N$, and all hatted quantities are evaluated at $\hat{\gamma}$ or $\hat{\boldsymbol{\theta}}_{w}$. The asymptotic variance of $\hat{\boldsymbol{\theta}}_{w}$ is consistently estimated as $\hat{\mathbf{A}}^{-1} \hat{\mathbf{D}} \hat{\mathbf{A}}^{-1} / N$.

We can compare expression (19.87) with the asymptotic variance that we would obtain by using a known value of $\gamma_{o}$ in place of the conditional MLE, $\hat{\gamma}$. As before, let $\tilde{\boldsymbol{\theta}}_{w}$ denote the estimator that uses $1 / G\left(z_{i}, \gamma_{o}\right)$ as the weights. Then\\
$\sqrt{N}\left(\tilde{\boldsymbol{\theta}}_{w}-\boldsymbol{\theta}_{o}\right) \stackrel{a}{\sim} \operatorname{Normal}\left(0, \mathbf{A}_{o}^{-1} \mathbf{B}_{o} \mathbf{A}_{o}^{-1}\right)$,\\
where $\mathbf{B}_{o} \equiv \mathrm{E}\left(\mathbf{k}_{i} \mathbf{k}_{i}^{\prime}\right)$. Because $\mathbf{B}_{o}-\mathbf{D}_{o}$ is positive semidefinite, Avar $\sqrt{N}\left(\tilde{\boldsymbol{\theta}}_{w}-\boldsymbol{\theta}_{o}\right)-$ Avar $\sqrt{N}\left(\hat{\boldsymbol{\theta}}_{w}-\boldsymbol{\theta}_{o}\right)$ is positive semidefinite.

As an example, consider the linear regression model $y=\mathbf{x} \boldsymbol{\beta}_{o}+u, \mathrm{E}\left(\mathbf{x}^{\prime} u\right)=\mathbf{0}$, and suppose the estimated probabilities $\hat{p}_{i}=G\left(\mathbf{z}_{i}, \hat{\gamma}\right)$ are from a logit estimation. The gradient for the logit estimation is $\hat{\mathbf{d}}_{i}^{\prime}=\mathbf{z}_{i}\left[s_{i}-\Lambda\left(\mathbf{z}_{i}, \hat{\gamma}\right)\right]$, a $1 \times M$ vector (the dimension of $\mathbf{z}_{i}$, which includes a constant). The selected, weighted gradient for the linear regression problem is $\hat{\mathbf{k}}_{i}^{\prime}=s_{i} \mathbf{x}_{i} \hat{u}_{i} / \hat{p}_{i}$, where $\hat{u}_{i}=y_{i}-\mathbf{x}_{i} \hat{\boldsymbol{\beta}}_{w}$ are the residuals after the IPW estimation. The adjustment to the asymptotic variance is obtained by getting the (row) residuals $\hat{\mathbf{e}}_{i}^{\prime}$ from the regression $s_{i} \mathbf{x}_{i} \hat{u}_{i} / \hat{p}_{i}$ on $\mathbf{z}_{i}\left[s_{i}-\Lambda\left(\mathbf{z}_{i}, \hat{\gamma}\right)\right]$ using all observations. Then, $\hat{\mathbf{D}}$ is constructed as in expression (19.90), and, in this case, $\hat{\mathbf{A}} \equiv N^{-1} \sum_{i=1}^{N}\left(s_{i} / \hat{p}_{i}\right) \mathbf{x}_{i}^{\prime} \mathbf{x}_{i}$. The asymptotic variance of $\hat{\boldsymbol{\beta}}_{w}$ is estimated as\\
$\left[\sum_{i=1}^{N}\left(s_{i} / \hat{p}_{i}\right) \mathbf{x}_{i}^{\prime} \mathbf{x}_{i} \cdot\right]^{-1}\left(\sum_{i=1}^{N} \hat{\mathbf{e}}_{i} \hat{\mathbf{e}}_{i}^{\prime}\right)\left[\sum_{i=1}^{N}\left(s_{i} / \hat{p}_{i}\right) \mathbf{x}_{i}^{\prime} \mathbf{x}_{i} \cdot\right]^{-1}$.

The conservative estimate would replace $\hat{\mathbf{e}}_{i}$ with $s_{i} \mathbf{x}_{i}^{\prime} \hat{u}_{i} / \hat{p}_{i}$, in which case the estimator looks just like a "heteroskedasticity"-robust sandwich estimator in the context of weighted least squares.

As with all general efficiency results, the conclusion that a difference in asymptotic variances is positive semidefinite does not guarantee an efficiency gain in particular cases: the statement includes the possibility that there is no difference in the asymptotic variances.

One case where there is no efficiency gain from using the estimated probabilities occurs when the missing data mechanism is, in a particular sense, exogenous. Consider the linear regression model. If we impose the assumption\\
$\mathrm{E}(u \mid \mathbf{x}, \mathbf{z})=0$,\\
which now means that $\mathrm{E}(y \mid \mathbf{x}, \mathbf{z})=\mathrm{E}(y \mid \mathbf{x})=\mathbf{x} \boldsymbol{\beta}$, and $\mathrm{P}(s=1 \mid \mathbf{x}, y, \mathbf{z})=\mathrm{P}(s=1 \mid \mathbf{z})$, then the asymptotic variance is given by expression (19.90) whether or not we estimate the weights. Further, the probability weights can come from a misspecified estimation problem without affecting consistency. For maximum likelihood problems, the exogeneity condition is $\mathrm{D}(y \mid \mathbf{x}, \mathbf{z})=\mathrm{D}(y \mid \mathbf{x})$. An important situation where this condition and assumption (19.92) hold occurs when $\mathbf{z}=\mathbf{x}$ and the model for $\mathrm{D}(y \mid \mathbf{x})$ or $\mathrm{E}(y \mid \mathbf{x})$ is correctly specified; the exogeneity condition on selection is $\mathrm{P}(s=1 \mid \mathbf{x}, y)=\mathrm{P}(s=1 \mid \mathbf{x})$.

Not surprisingly, in the general case of exogenous sampling, the unweighted estimator is also consistent. Along with assumption 19.6b, we generally define exogenous selection in the context of $M$-estimation as follows. Assume $\boldsymbol{\theta}_{o}$ satisfies


\begin{equation*}
\boldsymbol{\theta}_{o}=\underset{\boldsymbol{\theta} \in \boldsymbol{\Theta}}{\operatorname{argmin}} \mathrm{E}[q(\mathbf{w}, \boldsymbol{\theta}) \mid \mathbf{z}] \tag{19.93}
\end{equation*}


for all outcomes $\mathbf{z}$. This requirement may seem abstract, but it holds for the usual estimation methods when the appropriate underlying feature of the conditional distribution is correctly specified and $\mathbf{z}$ is exogenous in the model. For example, for nonlinear least squares and quasi-MLE in the linear exponential family, assumption (19.93) holds whenever $\mathrm{E}(y \mid \mathbf{x}, \mathbf{z})=\mathrm{E}(y \mid \mathbf{x})$ and the latter is correctly specified. This condition holds for conditional MLE, too.

Wooldridge (2007, Theorem 4.3) shows that, under exogenous sampling, the unweighted estimator is more efficient than the weighted estimator under a version of the conditional information matrix equality. Generally, the condition is stated as


\begin{equation*}
\mathrm{E}\left[\nabla_{\boldsymbol{\theta}} q\left(\mathbf{w}, \boldsymbol{\theta}_{o}\right)^{\prime} \nabla_{\boldsymbol{\theta}} q\left(\mathbf{w}, \boldsymbol{\theta}_{o}\right) \mid \mathbf{z}\right]=\sigma_{o}^{2} \mathrm{E}\left[\nabla_{\theta}^{2} q\left(\mathbf{w}, \boldsymbol{\theta}_{o}\right) \mid \mathbf{z}\right] \tag{19.94}
\end{equation*}


which holds under assumption (19.93) under suitable regularity conditions for correctly specified CMLE, nonlinear least squares under homoskedasticity, and quasiMLE under the assumption that the GLM variance assumption holds. The usual\\
nonrobust variance matrix estimators on the selected sample are valid. See Wooldridge (2007) for further discussion.

The previous results help to inform the decision of when weighting should and should not be used. If features of an unconditional distribution, say $\mathbf{D}(\mathbf{w})$, are of interest, unweighted estimators consistently estimate the parameters only if $\mathrm{P}(s=1 \mid \mathbf{w}) =\mathrm{P}(s=1)$-that is, the data are missing completely at random (Rubin, 1976). Of course, consistency of the weighted estimator relies on the presence of $\mathbf{z}$ such that $\mathrm{P}(s=1 \mid \mathbf{w}, \mathbf{z})=\mathrm{P}(s=1 \mid \mathbf{z})$, but maintaining such an assumption is often one's only recourse.

The decision to weight is more subtle when we begin with the premise that some feature of a conditional distribution, $\mathrm{D}(\mathbf{y} \mid \mathbf{x})$, is of interest. Wooldridge (2007) describes several scenarios concerning both consistent estimation and efficient estimation. Here, we briefly discuss a problem that can arise with weighting if data are missing on some of the conditioning variables, $\mathbf{x}$. The issue arises because, in general, if data are missing on elements of $\mathbf{x}$-due to attrition, say, or nonresponse-these elements generally cannot be included in the selection predictors, $\mathbf{z}$. But then suppose that selection is entirely a function of $\mathbf{x}: \mathrm{P}(s=1 \mid \mathbf{x}, \mathbf{y})=\mathrm{P}(s=1 \mid \mathbf{x})$ (even though elements of $\mathbf{x}$ are not observed). Then the unweighted estimator is consistent if (the feature of) $\mathrm{D}(\mathbf{y} \mid \mathbf{x})$ is correctly specified. Generally, if $\mathbf{z}$ omits any element of $\mathbf{x}$, the weighted estimator would actually be inconsistent. In effect, the wrong weights are used because the condition $\mathrm{P}(s=1 \mid \mathbf{x}, \mathbf{y}, \mathbf{z})=\mathrm{P}(s=1 \mid \mathbf{z})$ cannot hold. See Wooldridge (2007) for additional discussion.

When $\mathbf{z}$ can be chosen to include $\mathbf{x}$, the case for weighting is much stronger: if selection does depend only on $\mathbf{x}$, that fact will be picked up in large enough samples if the model for $\mathrm{P}(s=1 \mid \mathbf{z})$ is sufficiently flexible. Although this approach covers the case of treatment effects estimation-see Chapter 21-it does not cover general missing data problems. Therefore, one has to be cautious when using IPW when some conditioning variables are missing: important differences between the weighted and unweighted estimates cannot necessarily be attributed to a problem with the unweighted estimator, and we can never know why the two estimates are different (unless we have access to a random sample).

When the selection probabilities reflect stratified sampling and are determined from the sample design, the case for weighting is also stronger; see Chapter 20.

\subsection*{19.9 Sample Selection and Attrition in Linear Panel Data Models}
In our treatment of panel data models we have assumed that a balanced panel is available-each cross section unit has the same time periods available. Often, some\\
time periods are missing for some units in the population of interest, and we are left with an unbalanced panel. Unbalanced panels can arise for several reasons. First, the survey design may simply rotate people or firms out of the sample based on prespecified rules. For example, if a survey of individuals begins at time $t=1$, at time $t=2$ some of the original people may be dropped and new people added. At $t=3$ some additional people might be dropped and others added; and so on. This is an example of a rotating panel.

Provided the decision to rotate units out of a panel is made randomly, unbalanced panels are fairly easy to deal with, as we will see shortly. A more complicated problem arises when attrition from a panel is due to units electing to drop out. If this decision is based on factors that are systematically related to the response variable, even after we condition on explanatory variables, a sample selection problem can resultjust as in the cross section case. Nevertheless, a panel data set provides us with the means to handle, in a simple fashion, attrition that is based on a time-constant, unobserved effect, provided we use first-differencing methods; we show this in Section 19.9.3.

A different kind of sample selection problem occurs when people do not disappear from the panel but certain variables are unobserved for at least some time periods. This is the incidental truncation problem discussed in Section 19.6. A leading case is estimating a wage offer equation using a panel of individuals. Even if the population of interest is people who are employed in the initial year, some people will become unemployed in subsequent years. For those people we cannot observe a wage offer, just as in the cross-sectional case. This situation is different from the attrition problem where people leave the sample entirely and, usually, do not reappear in later years. In the incidental truncation case we observe some variables on everyone in each time period.

\subsection*{19.9.1 Fixed and Random Effects Estimation with Unbalanced Panels}
We begin by studying assumptions under which the usual fixed effects estimator on the unbalanced panel is consistent. The model is the usual linear, unobserved effects model under random sampling in the cross section: for any $i$,


\begin{equation*}
y_{i t}=\mathbf{x}_{i t} \boldsymbol{\beta}+c_{i}+u_{i t}, \quad t=1, \ldots, T \tag{19.95}
\end{equation*}


where $\mathbf{x}_{i t}$ is $1 \times K$ and $\boldsymbol{\beta}$ is the $K \times 1$ vector of interest. As before, we assume that $N$ cross section observations are available and the asymptotic analysis is as $N \rightarrow \infty$. We first cover the case where $c_{i}$ is allowed to be correlated with $\mathbf{x}_{i t}$, so that all elements of $\mathbf{x}_{i t}$ are time varying.

We treated the case where all $T$ time periods are available in Chapters 10 and 11. Now we consider the case where some time periods might be missing for some of the cross section draws. Think of $t=1$ as the first time period for which data on anyone in the population are available, and $t=T$ as the last possible time period. For a random draw $i$ from the population, let $\mathbf{s}_{i} \equiv\left(s_{i 1}, \ldots, s_{i T}\right)^{\prime}$ denote the $T \times 1$ vector of selection indicators: $s_{i t}=1$ if $\left(\mathbf{x}_{i t}, y_{i t}\right)$ is observed, and zero otherwise. Generally, we have an unbalanced panel. We can treat $\left\{\left(\mathbf{x}_{i}, \mathbf{y}_{i}, \mathbf{s}_{i}\right): i=1,2, \ldots, N\right\}$ as a random sample from the population; the selection indicators tell us which time periods are missing for each $i$.

We can easily find assumptions under which the fixed effects estimator on the unbalanced panel is consistent by writing it as


\begin{align*}
\hat{\boldsymbol{\beta}} & =\left(N^{-1} \sum_{i=1}^{N} \sum_{t=1}^{T} s_{i t} \ddot{\mathbf{x}}_{i t}^{\prime} \ddot{\mathbf{x}}_{i t}\right)^{-1}\left(N^{-1} \sum_{i=1}^{N} \sum_{t=1}^{T} s_{i t} \ddot{\mathbf{x}}_{i t}^{\prime} \ddot{y}_{i t}\right) \\
& =\boldsymbol{\beta}+\left(N^{-1} \sum_{i=1}^{N} \sum_{t=1}^{T} s_{i t} \ddot{\mathbf{x}}_{i t}^{\prime} \ddot{\mathbf{x}}_{i t}\right)^{-1}\left(N^{-1} \sum_{i=1}^{N} \sum_{t=1}^{T} s_{i t} \ddot{\mathbf{x}}_{i t}^{\prime} u_{i t}\right), \tag{19.96}
\end{align*}


where we define

$$
\ddot{\mathbf{x}}_{i t} \equiv \mathbf{x}_{i t}-T_{i}^{-1} \sum_{r=1}^{T} s_{i r} \mathbf{x}_{i r}, \quad \ddot{y}_{i t} \equiv y_{i t}-T_{i}^{-1} \sum_{r=1}^{T} s_{i r} y_{i r}, \quad \text { and } \quad T_{i} \equiv \sum_{t=1}^{T} s_{i t}
$$

That is, $T_{i}$ is the number of time periods observed for cross section $i$, and we apply the within transformation on the available time periods.

If fixed effects on the unbalanced panel is to be consistent, we should have $\mathrm{E}\left(s_{i t} \ddot{\mathbf{x}}_{i t}^{\prime} u_{i t}\right)=0$ for all $t$. Now, since $\ddot{\mathbf{x}}_{i t}$ depends on all of $\mathbf{x}_{i}$ and $\mathbf{s}_{i}$, a form of strict exogeneity is needed.

ASSUMPTION 19.8: (a) $\mathrm{E}\left(u_{i t} \mid \mathbf{x}_{i}, \mathbf{s}_{i}, c_{i}\right)=0, t=1,2, \ldots, T$; (b) $\sum_{t=1}^{T} \mathrm{E}\left(s_{i t} \ddot{\mathbf{x}}_{i t}^{\prime} \ddot{\mathbf{x}}_{i t}\right)$ is nonsingular; and (c) $\mathrm{E}\left(\mathbf{u}_{i} \mathbf{u}_{i}^{\prime} \mid \mathbf{x}_{i}, \mathbf{s}_{i}, c_{i}\right)=\sigma_{u}^{2} \mathbf{I}_{T}$.

Under assumption 19.8a, $\mathrm{E}\left(s_{i t} \ddot{\mathbf{x}}_{i t}^{\prime} u_{i t}\right)=\mathbf{0}$ from the law of iterated expectations (because $s_{i t} \ddot{\mathbf{x}}_{i t}$ is a function of $\left(\mathbf{x}_{i}, \mathbf{s}_{i}\right)$ ). The second assumption is the rank condition on the expected outer product matrix, after accounting for sample selection; naturally, it rules out time-constant elements in $\mathbf{x}_{i t}$. These first two assumptions ensure consistency of FE on the unbalanced panel.

In the case of a randomly rotating panel, and in other cases where selection is entirely random, $\mathbf{s}_{i}$ is independent of $\left(\mathbf{u}_{i}, \mathbf{x}_{i}, c_{i}\right)$, in which case assumption 19.8a\\
follows under the standard fixed effects assumption $\mathrm{E}\left(u_{i t} \mid \mathbf{x}_{i}, c_{i}\right)=0$ for all $t$. In this case, the natural assumptions on the population model imply consistency and asympotic normality on the unbalanced panel. Assumption 19.8a also holds under much weaker conditions. In particular, it does not assume anything about the relationship between $\mathbf{s}_{i}$ and $\left(\mathbf{x}_{i}, c_{i}\right)$. Therefore, if we think selection in all time periods is correlated with $c_{i}$ or $\mathbf{x}_{i}$, but that $u_{i t}$ is mean independent of $\mathbf{s}_{i}$ given ( $\mathbf{x}_{i}, c_{i}$ ) for all $t$, then FE on the unbalanced panel is consistent and asymptotically normal. This assumption may be a reasonable approximation, especially for short panels. What assumption 19.8a rules out is selection that is partially correlated with the idiosyncratic errors, $u_{i t}$.

When we add assumption 19.8c, standard inference procedures based on FE are valid. In particular, under assumptions 19.8a and 19.8c,\\
$\operatorname{Var}\left(\sum_{t=1}^{T} s_{i t} \ddot{\mathbf{x}}_{i t}^{\prime} u_{i t}\right)=\sigma_{u}^{2}\left[\sum_{t=1}^{T} \mathrm{E}\left(s_{i t} \ddot{\mathbf{x}}_{i t}^{\prime} \ddot{\mathbf{x}}_{i t}\right)\right]$.\\
Therefore, the asymptotic variance of the fixed effects estimator is estimated as


\begin{equation*}
\hat{\sigma}_{u}^{2}\left(\sum_{i=1}^{N} \sum_{t=1}^{T} s_{i t} \ddot{\mathbf{x}}_{i t}^{\prime} \ddot{\mathbf{x}}_{i t}\right)^{-1} . \tag{19.97}
\end{equation*}


The estimator $\hat{\sigma}_{u}^{2}$ can be derived from

$$
\mathrm{E}\left(\sum_{t=1}^{T} s_{i t} \ddot{u}_{i t}^{2}\right)=\mathrm{E}\left[\sum_{t=1}^{T} s_{i t} \mathrm{E}\left(\ddot{u}_{i t}^{2} \mid \mathbf{s}_{i}\right)\right]=\mathrm{E}\left\{T_{i}\left[\sigma_{u}^{2}\left(1-1 / T_{i}\right)\right]\right\}=\sigma_{u}^{2} \mathrm{E}\left[\left(T_{i}-1\right)\right]
$$

Now, define the FE residuals as $\hat{u}_{i t}=\ddot{y}_{i t}-\ddot{\mathbf{x}}_{i t} \hat{\boldsymbol{\beta}}$ when $s_{i t}=1$. Then, because $N^{-1} \sum_{i=1}^{N}\left(T_{i}-1\right) \xrightarrow{p} \mathrm{E}\left(T_{i}-1\right)$,

$$
\hat{\sigma}_{u}^{2}=\left[N^{-1} \sum_{i=1}^{N}\left(T_{i}-1\right)\right]^{-1} N^{-1} \sum_{i=1}^{N} \sum_{t=1}^{T} s_{i t} \hat{u}_{i t}^{2}=\left[\sum_{i=1}^{N}\left(T_{i}-1\right)\right]^{-1} \sum_{i=1}^{N} \sum_{t=1}^{T} s_{i t} \hat{u}_{i t}^{2}
$$

is consistent for $\sigma_{u}^{2}$ as $N \rightarrow \infty$. Standard software packages also make a degrees-offreedom adjustment by subtracting $K$ from $\sum_{i=1}^{N}\left(T_{i}-1\right)$. It follows that all of the usual test statistics based on an unbalanced fixed effects analysis are valid. In particular, the dummy variable regression discussed in Chapter 10 produces asymptotically valid statistics.

Because the FE estimator uses time demeaning, any unit $i$ for which $T_{i}=1$ drops out of the fixed effects estimator. To use these observations we would need to add more assumptions, such as the random effects assumption $\mathrm{E}\left(c_{i} \mid \mathbf{x}_{i}, \mathbf{s}_{i}\right)=0$.

Relaxing assumption 19.8c is easy: just apply the robust variance matrix estimator in equation (10.59) to the unbalanced panel. The only changes are that the rows of $\ddot{\mathbf{X}}_{i}$ are $s_{i t} \ddot{\mathbf{x}}_{i t}$ and the elements of $\ddot{\mathbf{u}}_{i}$ are $s_{i t} \ddot{u}_{i t}, t=1, \ldots, T$.

Under assumption 19.8, it is also valid to used a standard fixed effects analysis on any balanced subset of the unbalanced panel; in fact, we can condition on any outcomes of the $s_{i t}$. For example, if we use unit $i$ only when observations are available in all time periods, we are conditioning on $s_{i t}=1$ for all $t$. Comparing FE on the unbalanced panel with FE on a balanced panel is a sensible specification check.

Using similar arguments, it can be shown that any kind of differencing method on any subset of the observed panel is consistent. For example, with $T=3$, we observe cross section units with data for one, two, or three time periods. Those units with $T_{i}=1$ drop out, but any other combinations of differences can be used in a pooled OLS analysis. The analogues of assumption 19.8 for first differencing-for example, assumption 19.8c is replaced with $\mathrm{E}\left(\Delta \mathbf{u}_{i} \Delta \mathbf{u}_{i}^{\prime} \mid \mathbf{x}_{i}, \mathbf{s}_{i}, c_{i}\right)=\sigma_{e}^{2} \mathbf{I}_{T-1}$-ensure that the usual statistics from pooled OLS on the unbalanced first differences are asymptotically valid.

Random effects estimation on an unbalanced panel is somewhat more complicated because of the GLS transformation, and its consistency hinges on stronger assumptions concerning the selection mechanism. In addition to a rank condition, the key extra requirement (compared with FE) is\\
$\mathrm{E}\left(c_{i} \mid \mathbf{x}_{i}, \mathbf{s}_{i}\right)=\mathrm{E}\left(c_{i}\right)$,\\
so that the heterogeneity is mean independent of the covariates and selection in all time periods. Assumption (19.98) is generally very restrictive. It does allow for units being randomly rotated in and out of a panel or for rotation to depend on the observed factors in $\mathbf{x}_{i t}$.

A careful analysis of GLS estimation with unbalanced panels is notationally cumbersome when we explicitly introduce the selection indicators. Nevertheless, the description of the random effects estimator is easy to describe because of the "exchangeable" nature of the random effects covariance structure. In particular, under assumption RE. 3 from Section 10.4 (for the balanced case), the covariance between any two time periods is the same. Along with the assumption of constant variance over time, this leads to a straightforward GLS transformation. In the balanced case, the data are quasi-time-demeaned using (an estimate of) the parameter $\lambda= 1-\left[\sigma_{u}^{2} /\left(\sigma_{u}^{2}+T \sigma_{c}^{2}\right)\right]^{1 / 2}$; see Section 10.7.2. For the unbalanced case, where selection is exogenous in the sense of assumption (19.98), quasi-time-demeaned applies, but the parameter depends on $i$, namely,\\
$\lambda_{i}=1-\left[\sigma_{u}^{2} /\left(\sigma_{u}^{2}+T_{i} \sigma_{c}^{2}\right)\right]^{1 / 2}$,\\
where $T_{i}=\sum_{r=1}^{T} s_{i r}$ is the number of time periods available for unit $i$. Unlike in the balanced case, $\lambda_{i}$ is properly viewed as a random variable because it is a function of $T_{i}$. But, under assumption 19.8a and assumption (19.98), $T_{i}$ is appropriately exogenous because $\mathrm{E}\left(v_{i t} \mid \mathbf{x}_{i}, T_{i}\right)=0$, where $v_{i t}=c_{i}+u_{i t}$. Therefore, transforming the equation using any function of $T_{i}$ still leads to consistent estimation under assumption (19.98). See Baltagi (2001, Section 9.2) for the case where $T_{i}$ is treated as nonrandom.

If we knew $\lambda_{i}$, the RE estimator on the unbalanced panel would just be pooled OLS on the quasi-time-demeaned data $\breve{y}_{i t}=y_{i t}-\lambda_{i} \bar{y}_{i}$ and $\breve{\mathbf{x}}_{i t}=\mathbf{x}_{i t}-\lambda_{i} \overline{\mathbf{x}}_{i}$, where, naturally, the time averages are of the observed data. To operationalize the RE estimator, we need initial estimators of $\sigma_{u}^{2}$ and $\sigma_{c}^{2}$, but these are obtained by modifying the estimators in the unbalanced case.

With modern software, implementing random effects on an unbalanced panel is straightforward. But one should know that the RE estimator imposes strong assumptions on the nature of the missing data mechanism. Further, as in the balanced case, there are good reasons to believe the RE variance structure should not be taken literally. That is, even if we assume, in addition to assumption 19.8a, assumption 19.8c, and assumption (19.98),


\begin{equation*}
\operatorname{Var}\left(c_{i} \mid \mathbf{x}_{i}, \mathbf{s}_{i}\right)=\operatorname{Var}\left(c_{i}\right)=\sigma_{c}^{2}, \tag{19.100}
\end{equation*}


it is still likely that the composite variance covariance matrix in the population, $\operatorname{Var}\left(\mathbf{v}_{i} \mid \mathbf{x}_{i}\right)$, does not have the RE structure. We already know in the balanced case that using the incorrect variance matrix in GLS estimation does not cause inconsistency in estimating $\boldsymbol{\beta}$, provided the explanatory variables are strictly exogenous with respect to $v_{i t}=c_{i}+u_{i t}$. Not surprisingly, in the unbalanced case we effectively must add strict exogeneity of the selection process. Under this assumption, RE on the unbalanced panel is generally consistent and $\sqrt{N}$-asymptotically normal for any variance structure for $\operatorname{Var}\left(\mathbf{v}_{i} \mid \mathbf{x}_{i}, \mathbf{c}_{i}\right)$. But we need to compute a fully robust variance matrix estimator for $\hat{\boldsymbol{\beta}}_{R E}$, which is typically available in modern econometrics packages for the balanced and unbalanced cases.

\subsection*{19.9.2 Testing and Correcting for Sample Selection Bias}
The results in the previous subsection imply that sample selection in a fixed effects context is only a problem when selection is related to the idiosyncratic errors, $u_{i t}$. Therefore, any test for selection bias should test only this assumption. A simple test was suggested by Nijman and Verbeek (1992) in the context of random effects esti-\\
mation, but it works for fixed effects as well: add, say, the lagged selection indicator, $s_{i, t-1}$, to the equation, estimate the model by fixed effects (on the unbalanced panel), and do a $t$ test (perhaps making it fully robust) for the significance of $s_{i, t-1}$. (This method loses the first time period for all observations.) Under the null hypothesis, $u_{i t}$ is uncorrelated with $s_{i r}$ for all $r$, and so selection in the previous time period should not be significant in the equation at time $t$. (Incidentally, it never makes sense to put $s_{i t}$ in the equation at time $t$ because $s_{i t}=1$ for all $i$ and $t$ in the selected subsample.)

Putting $s_{i, t-1}$ does not work if $s_{i, t-1}$ is unity whenever $s_{i t}$ is unity because then there is no variation in $s_{i, t-1}$ in the selected sample. This is the case in attrition problems if (say) a person can only appear in period $t$ if he or she appeared in $t-1$. An alternative is to include a lead of the selection indicator, $s_{i, t+1}$. For observations $i$ that are in the sample every time period, $s_{i, t+1}$ is always zero. But for attriters, $s_{i, t+1}$ switches from zero to one in the period just before attrition. Alternatively, we can define a variable at $t$, say $r_{i, t+1}$, which is the number of periods after period $t$ that unit $i$ is in the sample. Either way, if we use fixed effects or first differencing, we need $T>2$ time periods to carry out the test.

For RE we have other possibilities, such as adding $T_{i}$ as an additional regressor and using a $t$ test.

For incidental truncation problems it makes sense to extend Heckman's (1976) test to the unobserved effects panel data context. This is done in Wooldridge (1995a). Write the equation of interest as


\begin{equation*}
y_{i t 1}=\mathbf{x}_{i t 1} \boldsymbol{\beta}_{1}+c_{i 1}+u_{i t 1}, \quad t=1, \ldots, T \tag{19.101}
\end{equation*}


Initially, suppose that $y_{i t 1}$ is observed only if the binary selection indicator, $s_{i t 2}$, is unity. Let $\mathbf{x}_{i t}$ denote the set of all exogenous variables at time $t$; we assume that these are observed in every time period, and $\mathbf{x}_{i t 1}$ is a subset of $\mathbf{x}_{i t}$. Suppose that, for each $t$, $s_{i t 2}$ is determined by the probit equation


\begin{equation*}
s_{i t 2}=1\left[\mathbf{x}_{i} \boldsymbol{\psi}_{t 2}+v_{i t 2}>0\right], \quad v_{i t 2} \mid \mathbf{x}_{i} \sim \operatorname{Normal}(0,1), \tag{19.102}
\end{equation*}


where $\mathbf{x}_{i}$ contains unity. This is best viewed as a reduced-form selection equation: we let the explanatory variables in all time periods appear in the selection equation at time $t$ to allow for general selection models, including those with unobserved effect and the Chamberlain (1980) device discussed in Section 15.8.2, as well as certain dynamic models of selection. A Mundlak (1978) approach would replace $\mathbf{x}_{i}$ with $\left(\mathbf{x}_{i t}, \overline{\mathbf{x}}_{i}\right)$ at time $t$ and assume that coefficients are constant across time. (See equation (15.68).) Then the parameters can be estimated by pooled probit, greatly conserving on degrees of freedom. Such conservation may be important for small $N$. For testing purposes, under the null hypothesis it does not matter whether equation (19.102) is\\
the proper model of sample selection, but we will need to assume equation (19.102), or a Mundlak version of it, when correcting for sample selection.

Under the null hypothesis in assumption 19.8a (with the obvious notational changes), the inverse Mills ratio obtained from the sample selection probit should not be significant in the equation estimated by fixed effects. Thus, let $\hat{\lambda}_{i t 2}$ be the estimated Mills ratios from estimating equation (19.102) by pooled probit across $i$ and $t$. Then a valid test of the null hypothesis is a $t$ statistic on $\hat{\lambda}_{i t 2}$ in the FE estimation on the unbalanced panel. Under assumption 19.8c the usual $t$ statistic is valid, but the approach works whether or not the $u_{i t 1}$ are homoskedastic and serially uncorrelated: just compute the robust standard error. Wooldridge (1995a) shows formally that the first-stage estimation of $\boldsymbol{\psi}_{2}$ does not affect the limiting distribution of the $t$ statistic under $\mathrm{H}_{0}$. This conclusion also follows from the results in Chapter 12 on M-estimation.

Correcting for sample selection requires much more care. Unfortunately, under any assumptions that actually allow for an unobserved effect in the underlying selection equation, adding $\hat{\lambda}_{i t 2}$ to equation (19.101) and using FE does not produce consistent estimators. To see why, suppose\\
$s_{i t 2}=1\left[\mathbf{x}_{i t} \boldsymbol{\delta}_{2}+c_{i 2}+a_{i t 2}>0\right], \quad a_{i t 2} \mid\left(\mathbf{x}_{i}, c_{i 1}, c_{i 2}\right) \sim \operatorname{Normal}(0,1)$.

Then, to get equation (19.102), $v_{i t 2}$ depends on $a_{i t 2}$ and, at least partially, on $c_{i 2}$. Now, suppose we make the strong assumption $\mathrm{E}\left(u_{i t 1} \mid \mathbf{x}_{i}, c_{i 1}, c_{i 2}, \mathbf{v}_{i 2}\right)=g_{i 1}+\rho_{1} v_{i t 2}$, which would hold under the assumption that the $\left(u_{i t 1}, a_{i t 2}\right)$ are independent across $t$ conditional on $\left(\mathbf{x}_{i}, c_{i 1}, c_{i 2}\right)$. Then we have\\
$y_{i t 1}=\mathbf{x}_{i t 1} \boldsymbol{\beta}_{1}+\rho_{1} \mathrm{E}\left(v_{i t 2} \mid \mathbf{x}_{i}, \mathbf{s}_{i 2}\right)+\left(c_{i 1}+g_{i 1}\right)+e_{i t 1}+\rho_{1}\left[v_{i t 2}-\mathrm{E}\left(v_{i t 2} \mid \mathbf{x}_{i}, \mathbf{s}_{i 2}\right)\right]$.\\
The composite error, $e_{i t 1}+\rho_{1}\left[v_{i t 2}-\mathrm{E}\left(v_{i t 2} \mid \mathbf{x}_{i}, \mathbf{s}_{i 2}\right)\right]$, is uncorrelated with any function of $\left(\mathbf{x}_{i}, \mathbf{s}_{i 2}\right)$. The problem is that $\mathrm{E}\left(v_{i t 2} \mid \mathbf{x}_{i}, \mathbf{s}_{i 2}\right)$ depends on all elements in $\mathbf{s}_{i 2}$, and this expectation is complicated for even small $T$.

A method that does work is available using Chamberlain's approach to panel data models, but we need some linearity assumptions on the expected values of $u_{i t 1}$ and $c_{i 1}$ given $\mathbf{x}_{i}$ and $v_{i t 2}$.

ASSUMPTION 19.9: (a) The selection equation is given by equation (19.102); (b) $\mathrm{E}\left(u_{i t 1} \mid \mathbf{x}_{i}, v_{i t 2}\right)=\mathrm{E}\left(u_{i t 1} \mid v_{i t 2}\right)=\rho_{t 1} v_{i t 2}, \quad t=1, \ldots, T$; and (c) $\mathrm{E}\left(c_{i 1} \mid \mathbf{x}_{i}, v_{i t 2}\right)= \mathrm{L}\left(c_{i 1} \mid 1, \mathbf{x}_{i}, v_{i t 2}\right)$.

The second assumption is standard and follows under joint normality of $\left(u_{i t 1}, v_{i t 2}\right)$ when this vector is independent of $\mathbf{x}_{i}$. Assumption 19.9c implies that\\
$\mathrm{E}\left(c_{i 1} \mid \mathbf{x}_{i}, v_{i t 2}\right)=\mathbf{x}_{i} \boldsymbol{\pi}_{1}+\phi_{t 1} v_{i t 2}$,\\
where, by equation (19.102) and iterated expectations, $\mathrm{E}\left(c_{i 1} \mid \mathbf{x}_{i}\right)=\mathbf{x}_{i} \boldsymbol{\pi}_{1}+\mathrm{E}\left(v_{i t 2} \mid \mathbf{x}_{i t}\right) =\mathbf{x}_{i} \boldsymbol{\pi}_{1}$. These assumptions place no restrictions on the serial dependence in ( $u_{i t 1}, v_{i t 2}$ ). They do imply that\\
$\mathrm{E}\left(y_{i t 1} \mid \mathbf{x}_{i}, v_{i t 2}\right)=\mathbf{x}_{i t 1} \boldsymbol{\beta}_{1}+\mathbf{x}_{i} \boldsymbol{\pi}_{1}+y_{t 1} v_{i t 2}$.

Conditioning on $s_{i t 2}=1$ gives\\
$\mathrm{E}\left(y_{i t 1} \mid \mathbf{x}_{i}, s_{i t 2}=1\right)=\mathbf{x}_{i t 1} \boldsymbol{\beta}_{1}+\mathbf{x}_{i} \boldsymbol{\pi}_{1}+\gamma_{t 1} \lambda\left(\mathbf{x}_{i} \boldsymbol{\psi}_{t 2}\right)$.\\
Therefore, we can consistently estimate $\boldsymbol{\beta}_{1}$ by first estimating a probit of $s_{i t 2}$ on $\mathbf{x}_{i}$ for each $t$ and then saving the inverse Mills ratio, $\hat{\lambda}_{i t 2}$, all $i$ and $t$. Next, run the pooled OLS regression using the selected sample:\\
$y_{i t 1}$ on $\mathbf{x}_{i t 1}, \mathbf{x}_{i}, \hat{\lambda}_{i t 2}, d 2_{t} \hat{\lambda}_{i t 2}, \ldots, d T_{t} \hat{\lambda}_{i t 2} \quad$ for all $s_{i t 2}=1$,\\
where $d 2_{t}$ through $d T_{t}$ are time dummies. If $\gamma_{t 1}$ in equation (19.104) is constant across $t$, simply include $\hat{\lambda}_{i t 2}$ by itself in equation (19.105). Again, a simplification that can be practically useful is to replace $\mathbf{x}_{i}$ with the time average, $\overline{\mathbf{x}}_{i}$.

The asymptotic variance of $\hat{\boldsymbol{\beta}}_{1}$ needs to be corrected for general heteroskedasticity and serial correlation, as well as first-stage estimation of the $\boldsymbol{\psi}_{t 2}$. These corrections can be made using the formulas for two-step M-estimation from Chapter 12; Wooldridge (1995a) contains the formulas. Alternatively, the panel bootstrap can be used, where the resampling is done using the cross section units.

If the selection equation is of the Tobit form, we have somewhat more flexibility. Write the selection equation now as\\
$y_{i t 2}=\max \left(0, \mathbf{x}_{i} \boldsymbol{\psi}_{t 2}+v_{i t 2}\right), \quad v_{i t 2} \mid \mathbf{x}_{i} \sim \operatorname{Normal}\left(0, \sigma_{t 2}^{2}\right)$,\\
where $y_{i t 1}$ is observed if $y_{i t 2}>0$. Then, under assumption 19.8, with the Tobit selection equation in place of equation (19.102), consistent estimation follows from the pooled regression (19.105) where $\hat{\lambda}_{i t 2}$ is replaced by the Tobit residuals, $\hat{v}_{i t 2}$ when $y_{i t 2}>0\left(s_{i t 2}=1\right)$. The Tobit residuals are obtained from the $T$ cross section Tobits in equation (19.106); alternatively, especially with small $N$, we can use a Mundlak-type approach and use pooled Tobit with $\mathbf{x}_{i} \boldsymbol{\psi}_{t 2}$ replaced with $\mathbf{x}_{i t} \boldsymbol{\delta}_{2}+\overline{\mathbf{x}}_{i} \boldsymbol{\pi}_{2}$; see equation (17.82).

It is easy to see that we can add $\alpha_{1} y_{i t 2}$ to the structural equation (19.101), provided we make an explicit exclusion restriction in assumption 19.9. In particular, we must assume that $\mathrm{E}\left(c_{i 1} \mid \mathbf{x}_{i}, v_{i t 2}\right)=\mathbf{x}_{i 1} \boldsymbol{\pi}_{1}+\phi_{t 1} v_{i t 2}$, and that $\mathbf{x}_{i t 1}$ is a strict subset of $\mathbf{x}_{i t}$. Then, because $y_{i t 2}$ is a function of $\left(\mathbf{x}_{i}, v_{i t 2}\right)$, we can write $\mathrm{E}\left(y_{i t 1} \mid \mathbf{x}_{i}, v_{i t 2}\right)=\mathbf{x}_{i t 1} \boldsymbol{\beta}_{1}+ \alpha_{1} y_{i t 2}+\mathbf{x}_{i 1} \boldsymbol{\pi}_{1}+\gamma_{t 1} v_{i t 2}$. We obtain the Tobit residuals, $\hat{v}_{i t 2}$ for each $t$, and then run the\\
regression $y_{i t 1}$ on $\mathbf{x}_{i t 1}, y_{i t 2}, \mathbf{x}_{i 1}$, and $\hat{v}_{i t 2}$ (possibly interacted with time dummies) for the selected sample. If we do not have an exclusion restriction, this regression suffers from perfect multicollinearity. As an example, we can easily include hours worked in a wage offer function for panel data, provided we have a variable affecting labor supply (such as the number of young children) but not the wage offer.

A pure fixed effects approach is more fruitful when the selection equation is of the Tobit form. The following assumption comes from Wooldridge (1995a):

ASSUMPTION 19.10: (a) The selection equation is equation (19.106). (b) For some unobserved effect $g_{i 1}, \mathrm{E}\left(u_{i t 1} \mid \mathbf{x}_{i}, c_{i 1}, g_{i 1}, \mathbf{v}_{i 2}\right)=\mathrm{E}\left(u_{i t 1} \mid g_{i 1}, v_{i t 2}\right)=g_{i 1}+\rho_{1} v_{i t 2}$.

Under part b of this assumption,\\
$\mathrm{E}\left(y_{i t 1} \mid \mathbf{x}_{i}, \mathbf{v}_{i 2}, c_{i 1}, g_{i 1}\right)=\mathbf{x}_{i t 1} \boldsymbol{\beta}_{1}+\rho_{1} v_{i t 2}+f_{i 1}$,\\
where $f_{i 1}=c_{i 1}+g_{i 1}$. The same expectation holds when we also condition on $\mathbf{s}_{i 2}$ (since $\mathbf{s}_{i 2}$ is a function of $\mathbf{x}_{i}, \mathbf{v}_{i 2}$ ). Therefore, estimating equation (19.107) by fixed effects on the unbalanced panel would consistently estimate $\boldsymbol{\beta}_{1}$ and $\rho_{1}$. As usual, we replace $v_{i t 2}$ with the Tobit residuals $\hat{v}_{i t 2}$ whenever $y_{i t 2}>0$. A $t$ test of $\mathrm{H}_{0}: \rho_{1}=0$ is valid very generally as a test of the null hypothesis of no sample selection. If the $\left\{u_{i t 1}\right\}$ satisfy the standard homoskedasticity and serial uncorrelatedness assumptions, then the usual $t$ statistic is valid. A fully robust test may be warranted. (Again, with an exclusion restriction, we can add $y_{i t 2}$ as an additional explanatory variable.)

Wooldridge (1995a) discusses an important case where assumption 19.10b holds: in the Tobit version of equation (19.103) with ( $\mathbf{u}_{i 1}, \mathbf{a}_{i 2}$ ) independent of ( $\mathbf{x}_{i}, c_{i 1}, c_{i 2}$ ) and $\mathrm{E}\left(u_{i t 1} \mid \mathbf{a}_{i 2}\right)=\mathrm{E}\left(u_{i t 1} \mid a_{i t 2}\right)=\rho_{1} a_{i t 2}$. The second-to-last equality holds under the common assumption that $\left\{\left(u_{i t 1}, a_{i t 2}\right): t=1, \ldots, T\right\}$ is serially independent.

The preceding methods assume normality of the errors in the selection equation and, implicitly, the unobserved heterogeneity. Kyriazidou (1997) and Honoré and Kyriazidou (2000b) have proposed methods that do not require distributional assumptions. Dustmann and Rochina-Barrachina (2007) apply Wooldridge's (1995a) and Kyriazidou's (1997) methods to the problem of estimating a wage offer equation with selection into the work force.

Semykina and Wooldridge (in press) show how to extend sample selection corrections for panel data to the case of endogenous explanatory variables. As in the case of cross section corrections, the endogenous explanatory variables need not be observed in all time periods. As shown formally by Semykina and Wooldridge (in press), a simple test for sample selection bias-without taking a stand on endogeneity of the explanatory variables-is obtained by adding inverse Mills ratio terms to the unbal-\\
anced panel and using fixed effect 2SLS. Corrections are available by adding time averages of the exogenous variables (both exogenous explanatory variables and the instruments) and applying pooled 2SLS with inverse Mills ratio terms (which, as usual, are estimated in a first stage). See Semykina and Wooldridge (in press) for details.

\subsection*{19.9.3 Attrition}
We now turn specifically to testing and correcting for attrition in a linear, unobserved effects panel data model. General attrition, where units may reenter the sample after leaving, is complicated. We analyze a common special case. At $t=1$ a random sample is obtained from the relevant population-people, for concreteness. In $t=2$ and beyond, some people drop out of the sample for reasons that may not be entirely random. We assume that, once a person drops out, he or she is out forever: attrition is an absorbing state. Any panel data set with attrition can be set up in this way by ignoring any subsequent observations on units after they initially leave the sample. In Section 19.9.2 we discussed one way to test for attrition bias when we assume that attrition is an absorbing state: include $s_{i, t+1}$ as an additional explanatory variable in a fixed effects analysis, or use the number of subsequent periods in the sample.

One method for correcting for attrition bias is closely related to the corrections for incidental truncation covered in the previous subsection. Write the model for a random draw from the population as in equation (19.95), where we assume that $\left(\mathbf{x}_{i t}, y_{i t}\right)$ is observed for all $i$ when $t=1$. Let $s_{i t}$ denote the selection indicator for each time period, where $s_{i t}=1$ if $\left(\mathbf{x}_{i t}, y_{i t}\right)$ are observed. Because we ignore units once they initially leave the sample, $s_{i t}=1$ implies $s_{i r}=1$ for $r<t$.

The sequential nature of attrition makes first differencing a natural choice to remove the unobserved effect:\\
$\Delta y_{i t}=\Delta \mathbf{x}_{i t} \boldsymbol{\beta}+\Delta u_{i t}, \quad t=2, \ldots, T$.

Conditional on $s_{i, t-1}=1$, write a (reduced-form) selection equation for $t \geq 2$ as\\
$s_{i t}=1\left[\mathbf{w}_{i t} \boldsymbol{\delta}_{t}+v_{i t}>0\right], \quad v_{i t} \mid\left\{\Delta \mathbf{x}_{i t}, \mathbf{w}_{i t}, s_{i, t-1}=1\right\} \sim \operatorname{Normal}(0,1)$,\\
where $\mathbf{w}_{i t}$ must contain variables observed at time $t$ for all units with $s_{i, t-1}=1$. Good candidates for $\mathbf{w}_{i t}$ include the variables in $\mathbf{x}_{i, t-1}$ and any variables in $\mathbf{x}_{i t}$ that are observed at time $t$ when $s_{i, t-1}=1$ (for example, if $\mathbf{x}_{i t}$ contains lags of variables or a variable such as age). In general, the dimension of $\mathbf{w}_{i t}$ can grow with $t$. For example, if equation (19.95) is dynamically complete, then $y_{i, t-2}$ is orthogonal to $\Delta u_{i t}$, and so it can be an element of $\mathbf{w}_{i t}$. Since $y_{i, t-1}$ is correlated with $u_{i, t-1}$, it should not be included in $\mathbf{w}_{i t}$.

Under assumption (19.108), selection in time $t$, conditional on being in the sample at time $t-1\left(s_{i, t-1}=1\right)$, follows a probit model:\\
$\mathrm{P}\left(s_{i t}=1 \mid \mathbf{w}_{i t}, s_{i, t-1}=1\right)=\Phi\left(\mathbf{w}_{i t} \boldsymbol{\delta}_{t}\right), \quad t=2, \ldots, T$.

Therefore, we can estimate $\boldsymbol{\delta}_{t}, t=2, \ldots, T$, by a sequence of probits where, in each time period, one uses the units still in the sample in the previous time period.

In order to justify a two-step Heckman correction, we need to assume that, for those still in the sample at $t-1, \mathbf{x}_{i t}$ does not affect selection in time $t$ once we condition on $\mathbf{w}_{i t}$. A natural way to state this condition is\\
$\mathrm{P}\left(s_{i t}=1 \mid \Delta \mathbf{x}_{i t}, \mathbf{w}_{i t}, s_{i, t-1}=1\right)=\mathrm{P}\left(s_{i t}=1 \mid \mathbf{w}_{i t}, s_{i, t-1}=1\right)$.

Technically, assumption (19.111) is not strong enough. An assumption that is sufficient-stated in terms of the errors in the first-differenced equation-is that, conditional on $s_{i, t-1}=1$,\\
$\left(\Delta u_{i t}, v_{i t}\right)$ is independent of $\left(\Delta \mathbf{x}_{i t}, \mathbf{w}_{i t}\right)$,\\
which, of course, implies that $v_{i t}$ is independent of ( $\Delta \mathbf{x}_{i t}, \mathbf{w}_{i t}$ ) and implies condition (19.111) when $v_{i t}$ has a standard normal distribution. We also impose the standard linear functional form assumption:\\
$\mathrm{E}\left(\Delta u_{i t} \mid v_{i t}, s_{i, t-1}=1\right)=\rho_{t} v_{i t}$,\\
in which case\\
$\mathrm{E}\left(\Delta y_{i t} \mid \Delta \mathbf{x}_{i t}, \mathbf{w}_{i t}, s_{i t}=1\right)=\Delta \mathbf{x}_{i t} \boldsymbol{\beta}+\rho_{t} \lambda\left(\mathbf{w}_{i t} \boldsymbol{\delta}_{t}\right), \quad t=2, \ldots, T$,\\
where $\lambda\left(\mathbf{w}_{i t} \boldsymbol{\delta}_{t}\right)$ is the inverse Mills ratio.\\
Notice how, because $s_{i, t-1}=1$ when $s_{i t}=1$, we do not have to condition on $s_{i, t-1}$ in equation (19.113). It now follows from equation (19.113) that pooled OLS of $\Delta y_{i t}$ on $\Delta \mathbf{x}_{i t}, d 2_{t} \hat{\lambda}_{i t}, \ldots, d T_{t} \hat{\lambda}_{i t}, t=2, \ldots, T$, where the $\hat{\lambda}_{i t}$ are from the $T-1$ cross section probits in equation (19.110), is consistent for $\boldsymbol{\beta}_{1}$ and the $\rho_{t}$. A joint test of $\mathrm{H}_{0}: \rho_{t}=0$, $t=2, \ldots, T$, is a fairly simple test for attrition bias, although nothing guarantees serial independence of the errors.

There are two potential problems with this approach. For one, assumption (19.112) is restrictive because it means that $\mathbf{x}_{i t}$ does not affect attrition once the elements in $\mathbf{w}_{i t}$ have been controlled for. This is a very strong assumption. In fact, in one scenario where this assumption fails, it is actually harmful to apply the Heckman approach. To see how this result can occur, suppose that selection is actually a function of changes in the covariates in the sense that\\
$\mathrm{P}\left(s_{i t}=1 \mid \Delta \mathbf{x}_{i t}, \Delta u_{i t}\right)=\mathrm{P}\left(s_{i t}=1 \mid \Delta \mathbf{x}_{i t}\right)$.

Under this condition, pooled OLS estimation of equation (19.108), using the unbalanced panel, is consistent for $\boldsymbol{\beta}$. By contrast, the Heckman method just described is generally inconsistent if $\Delta \mathbf{x}_{i t}$ is not included in $\mathbf{w}_{i t}$ : it is very unlikely that\\
$\mathrm{E}\left(\Delta u_{i t} \mid \Delta \mathbf{x}_{i t}, \mathbf{w}_{i t}, v_{i t}, s_{i, t-1}=1\right)=\mathrm{E}\left(\Delta u_{i t} \mid v_{i t}, s_{i, t-1}=1\right)$\\
(because $v_{i t}$ generally cannot be assumed to be independent of $\Delta \mathbf{x}_{i t}$ if $\Delta \mathbf{x}_{i t}$ is not included in $\mathbf{w}_{i t}$ ).

A second shortcoming of the Heckman procedure just described is that it requires strict exogeneity of $\left\{\mathbf{x}_{i t}\right\}$ in the original (levels) equation. Fortunately, in some cases we can apply an instrumental variables approach that is consistent when $\left\{\mathbf{x}_{i t}\right\}$ is not strictly exogenous or always observed at time $t$.

Let $\mathbf{z}_{i t}$ be a vector of variables such that $\mathbf{z}_{i t}$ is redundant in the selection equation (possibly because $\mathbf{w}_{i t}$ contains $\mathbf{z}_{i t}$ ) and that $\mathbf{z}_{i t}$ is exogenous in the sense that equation (19.112) holds with $\mathbf{z}_{i t}$ in place of $\Delta \mathbf{x}_{i t}$; for example, $\mathbf{z}_{i t}$ should contain $\mathbf{x}_{i r}$ for $r<t$. Now, using an argument similar to the cross section case in Section 19.6.2, we can estimate the equation\\
$\Delta y_{i t}=\Delta \mathbf{x}_{i t} \boldsymbol{\beta}+\rho_{2} d 2_{t} \hat{\lambda}_{i t}+\cdots+\rho_{T} d T_{t} \hat{\lambda}_{i t}+$ error $_{i t}$\\
by instrumental variables with instruments $\left(\mathbf{z}_{i t}, d 2_{t} \hat{\lambda}_{i t}, \ldots, d T_{t} \hat{\lambda}_{i t}\right)$, using the selected sample. For example, the pooled 2SLS estimator on the selected sample is consistent and asymptotically normal, and attrition bias can be tested by a joint test of $\mathrm{H}_{0}$ : $\rho_{t}=0, t=2, \ldots, T$. Under $\mathrm{H}_{0}$, only serial correlation and heteroskedasticity adjustments are possibly needed. If $\mathrm{H}_{0}$ fails we have the usual generated regressors problem for estimating the asymptotic variance. Other IV procedures, such as GMM, can also be used, but they too must account for the generated regressors problem. The panel bootstrap is a simple alternative to the analytical formulas.

Example 19.9 (Dynamic Model with Attrition): Consider the model\\
$y_{i t}=\mathbf{g}_{i t} \gamma+\eta_{1} y_{i, t-1}+c_{i}+u_{i t}, \quad t=1, \ldots, T$,\\
where we assume that $\left(y_{i 0}, \mathbf{g}_{i 1}, y_{i 1}\right)$ are all observed for a random sample from the population. Assume that $\mathrm{E}\left(u_{i t} \mid \mathbf{g}_{i}, y_{i, t-1}, \ldots, y_{i 0}, c_{i}\right)=0$, so that $\mathbf{g}_{i t}$ is strictly exogenous. Then the explanatory variables in the probit at time $t$, $\mathbf{w}_{i t}$, can include $\mathbf{g}_{i, t-1}$, $y_{i, t-2}$, and further lags of these. After estimating the selection probit for each $t$, and differencing, we can estimate\\
$\Delta y_{i t}=\Delta \mathbf{g}_{i t} \boldsymbol{\beta}+\eta_{1} \Delta y_{i, t-1}+\rho_{3} d 3_{t} \hat{\lambda}_{i t}+\cdots+\rho_{T} d T_{t} \hat{\lambda}_{i t}+$ error $_{i t}$\\
by pooled 2SLS on the selected sample starting at $t=3$, using instruments $\left(\mathbf{g}_{i, t-1}, \mathbf{g}_{i, t-2}, y_{i, t-2}, y_{i, t-3}\right)$. As usual, there are other possibilities for the instruments.

Although the focus in this section has been on pure attrition, where units disappear entirely from the sample, the methods can also be used in the context of incidental truncation without strictly exogenous explanatory variables. For example, suppose we are interested in the population of men who are employed at $t=0$ and $t=1$, and we would like to estimate a dynamic wage equation with an unobserved effect. Problems arise if men become unemployed in future periods. Such events can be treated as an attrition problem if all subsequent time periods are dropped once a man first becomes unemployed. This approach loses information but makes the econometrics relatively straightforward, especially because, in the preceding general model, $\mathbf{x}_{i t}$ will always be observed at time $t$ and so can be included in the labor force participation probit (assuming that men do not leave the sample entirely). Things become much more complicated if we are interested in the wage offer for all working age men at $t=1$ because we have to deal with the sample selection problem into employment at $t=0$ and $t=1$.

The methods for attrition and selection just described apply only to linear models, and it is difficult to extend them to general nonlinear models. An alternative approach is based on inverse probability weighting (IPW), which we described in the cross section case in Section 19.8.

Moffitt, Fitzgerald, and Gottschalk (1999) (MFG) propose IPW to estimate linear panel data models under possibly nonrandom attrition. (MFG propose a different set of weights, analogous to those studied by Horowitz and Manski (1998), to solve missing data problems. The weights we use require estimation of only one attrition model, rather than two as in MFG.) IPW must be used with care to solve the attrition problem. As before, we assume that we have a random sample from the population at $t=1$. We are interested in some feature, such as the conditional mean, or maybe the entire conditional distribution, of $y_{i t}$ given $\mathbf{x}_{i t}$. Ideally, at each $t$ we would observe ( $y_{i t}, \mathbf{x}_{i t}$ ) for any unit that was in the random sample at $t=1$. Instead, we observe ( $y_{i t}, \mathbf{x}_{i t}$ ) only if $s_{i t}=1$. We can easily solve the attrition problem if we assume that, conditional on observables in the first time period, say, $\mathbf{z}_{i 1},\left(y_{i t}, \mathbf{x}_{i t}\right)$ is independent of $s_{i t}$ :


\begin{equation*}
\mathrm{P}\left(s_{i t}=1 \mid y_{i t}, \mathbf{x}_{i t}, \mathbf{z}_{i 1}\right)=\mathrm{P}\left(s_{i t}=1 \mid \mathbf{z}_{i 1}\right), \quad t=2, \ldots, T . \tag{19.118}
\end{equation*}


As in the cross section case, assumption (19.118) has been called "selection on observables" in the econometrics literature and "ignorable selection" or "unconfoundedness" in the statistics literature. (The Heckman method we just covered is more like the "selection on unobservables" that we discussed in Section 19.8.) Con-\\
dition (19.118) is a strong assumption in that it maintains that, in every time period, $\mathbf{z}_{i 1}$ is such a good predictor of attrition that the distribution of $s_{i t}$ given $\left\{\mathbf{z}_{i 1},\left(\mathbf{x}_{i t}, y_{i t}\right)\right\}$ does not depend on $\left(\mathbf{x}_{i t}, y_{i t}\right)$. If we have such an observed set of variables $\mathbf{z}_{i 1}$, it is not too surprising that we can account for nonrandom attrition for a large class of estimation problems.

As in the cross section case, IPW estimation typically involves two steps. First, for each $t$, we estimate a probit or logit of $s_{i t}$ on $\mathbf{z}_{i 1}$. (A crucial point is that the same cross section units-namely, all units appearing in the first time period-are used in the probit or logit for each time period.) Let $\hat{p}_{i t}$ be the fitted probabilities, $t=2, \ldots$, $T, i=1, \ldots, N$. In the second step, the objective function for ( $i, t$ ) is weighted by $1 / \hat{p}_{i t}$. For general M-estimation, the objective function is\\
$\sum_{i=1}^{N} \sum_{t=1}^{T}\left(s_{i t} / \hat{p}_{i t}\right) q_{t}\left(\mathbf{w}_{i t}, \boldsymbol{\theta}\right)$,\\
where $\mathbf{w}_{i t} \equiv\left(y_{i t}, \mathbf{x}_{i t}\right)$ and $q_{t}\left(\mathbf{w}_{i t}, \boldsymbol{\theta}\right)$ is the objective function in each time period. As usual, the selection indicator $s_{i t}$ chooses the observations where we actually observe data. (For $t=1, s_{i t}=\hat{p}_{i t}=1$ for all $i$.) For least squares, $q_{t}\left(\mathbf{w}_{i t}, \boldsymbol{\theta}\right)$ is simply the squared residual function; for partial MLE, $q_{t}\left(\mathbf{w}_{i t}, \boldsymbol{\theta}\right)$ is the log-likelihood function.

The argument for why IPW works is similar to the pure cross section case. Let $\boldsymbol{\theta}_{\mathrm{o}}$ denote the value of $\boldsymbol{\theta}$ that solves the population problem $\min _{\boldsymbol{\theta} \in \boldsymbol{\Theta}} \sum_{t=1}^{T} \mathrm{E}\left[q_{t}\left(\mathbf{w}_{i t}, \boldsymbol{\theta}\right)\right]$. Let $\boldsymbol{\delta}_{t}^{\mathrm{o}}$ denote the true values of the selection response parameters in each time period, so that $\mathrm{P}\left(s_{i t}=1 \mid \mathbf{z}_{i 1}\right)=p_{t}\left(\mathbf{z}_{i 1}, \boldsymbol{\delta}_{t}^{\mathrm{o}}\right) \equiv p_{i t}^{\mathrm{o}}$. Now, under standard regularity conditions, we can replace $p_{i t}^{\mathrm{o}}$ with $\hat{p}_{i t} \equiv p_{t}\left(\mathbf{z}_{i 1}, \hat{\boldsymbol{\delta}}_{t}\right)$ without affecting the consistency argument. So, apart from regularity conditions, it is sufficient to show that $\boldsymbol{\theta}_{\mathrm{o}}$ minimizes $\sum_{t=1}^{T} \mathrm{E}\left[\left(s_{i t} / p_{i t}^{\mathrm{o}}\right) q_{t}\left(\mathbf{w}_{i t}, \boldsymbol{\theta}\right)\right]$ over $\boldsymbol{\Theta}$. But, from iterated expectations,

$$
\begin{aligned}
\mathrm{E}\left[\left(s_{i t} / p_{i t}^{\mathrm{o}}\right) q_{t}\left(\mathbf{w}_{i t}, \boldsymbol{\theta}\right)\right] & =\mathrm{E}\left\{\mathrm{E}\left[\left(s_{i t} / p_{i t}^{\mathrm{o}}\right) q_{t}\left(\mathbf{w}_{i t}, \boldsymbol{\theta}\right) \mid \mathbf{w}_{i t}, \mathbf{z}_{i 1}\right]\right\} \\
& =\mathrm{E}\left\{\left[\mathrm{E}\left(s_{i t} \mid \mathbf{w}_{i t}, \mathbf{z}_{i 1}\right) / p_{i t}^{\mathrm{o}}\right] q_{t}\left(\mathbf{w}_{i t}, \boldsymbol{\theta}\right)\right\}=\mathrm{E}\left[q_{t}\left(\mathbf{w}_{i t}, \boldsymbol{\theta}\right)\right]
\end{aligned}
$$

because $\mathrm{E}\left(s_{i t} \mid \mathbf{w}_{i t}, \mathbf{z}_{i 1}\right)=\mathrm{P}\left(s_{i t}=1 \mid \mathbf{z}_{i 1}\right)$ by assumption (19.118). Therefore, the probability limit of the weighted objective function is identical to that of the unweighted function if we had no attrition problem. Using this simple analogy argument and standard two-step estimation results from Chapter 12 shows that the inverse probability weighting produces a consistent, $\sqrt{N}$-asymptotically normal estimator. The methods for adjusting the asymptotic variance matrix of two step M-estimatorsdescribed in Subsection 12.5.2-can be applied to the IPW M-estimator from equation (19.119). The panel bootstrap that accounts for the two estimation steps is also attractive.

MFG propose an IPW scheme where the conditioning variables in the attrition probits change across time. In particular, at time $t$ an attrition probit is estimated restricting attention to those units still in the sample at time $t-1$. (Out of this group, some are lost to attrition at time $t$, and some are not.) If we assume that attrition is an absorbing state, we can include in the conditioning variables, $\mathbf{z}_{i t}$, all values of $y$ and $\mathbf{x}$ dated at time $t-1$ and earlier (as well as other variables observed for all units in the sample at $t-1$ ). This approach is appealing because the ignorability assumption is much more plausible if we can condition on both recent responses and covariates. (That is, $\mathrm{P}\left(s_{i t}=1 \mid \mathbf{w}_{i t}, \mathbf{w}_{i, t-1}, \ldots, \mathbf{w}_{i 1}, s_{i, t-1}=1\right)=\mathrm{P}\left(s_{i t}=1 \mid \mathbf{w}_{i, t-1}, \ldots\right.$, $\mathbf{w}_{i 1}, s_{i, t-1}=1$ ) is more likely than assumption (19.64).) Unfortunately, obtaining the fitted probabilities in this way and using them in an IPW procedure does not generally produce consistent estimators. The problem is that the selection models at each time period are not representative of the population that was originally sampled at $t=1$. Letting $p_{i t}^{\mathrm{o}}=\mathrm{P}\left(s_{i t}=1 \mid \mathbf{w}_{i, t-1}, \ldots, \mathbf{w}_{i 1}, s_{i, t-1}=1\right)$, we can no longer use the iterated expectations argument to conclude that $\mathrm{E}\left[\left(s_{i t} / p_{i t}^{\mathrm{o}}\right) q_{t}\left(\mathbf{w}_{i t}, \boldsymbol{\theta}\right)\right]=\mathrm{E}\left[q_{t}\left(\mathbf{w}_{i t}, \boldsymbol{\theta}\right)\right]$. Only if $\mathrm{E}\left[q_{t}\left(\mathbf{w}_{i t}, \boldsymbol{\theta}\right)\right]=\mathrm{E}\left[q_{t}\left(\mathbf{w}_{i t}, \boldsymbol{\theta}\right) \mid s_{i, t-1}=1\right]$ for all $\boldsymbol{\theta}$ does the argument work, but this assumption essentially requires that $\mathbf{w}_{i t}$ be independent of $s_{i, t-1}$.

It is possible to allow the covariates in the selection probabilities to increase in richness over time, but the MFG procedure must be modified. For the case where attrition is an absorbing state, we can extend to the M-estimation case the nonlinear regression results of Robins, Rotnitzky, and Zhao (1995) (RRZ). It turns out that in some cases the probabilities to be used in IPW estimation can be constructed sequentially:


\begin{equation*}
p_{i t}\left(\delta_{\mathrm{t}}^{\mathrm{o}}\right) \equiv \pi_{i 2}\left(\gamma_{2}^{\mathrm{o}}\right) \pi_{i 3}\left(\gamma_{3}^{\mathrm{o}}\right) \cdots \pi_{i t}\left(\gamma_{t}^{\mathrm{o}}\right), \quad t=2, \ldots, T, \tag{19.120}
\end{equation*}


where


\begin{equation*}
\pi_{i t}\left(\gamma_{t}^{\mathrm{o}}\right) \equiv \mathrm{P}\left(s_{i t}=1 \mid \mathbf{z}_{i t}, s_{i, t-1}=1\right) . \tag{19.121}
\end{equation*}


In other words, as in the MFG procedure, we estimate probit models at each time $t$, restricted to units that are in the sample at $t-1$. The covariates in the probit are essentially everything we can observe for units in the sample at time $t-1$ that might affect attrition. For $t=2, \ldots, T$, let $\hat{\pi}_{i t}$ denote the fitted selection probabilities. Then we construct the probability weights as the product $\hat{p}_{i t} \equiv \hat{\pi}_{i 2} \hat{\pi}_{i 3} \cdots \hat{\pi}_{i t}$ and use the objective function (19.119). Naturally, this method only works under certain assumptions. The key ignorability of selection condition can be stated as


\begin{equation*}
\mathrm{P}\left(s_{i t}=1 \mid \mathbf{v}_{i 1}, \ldots, \mathbf{v}_{i T}, s_{i, t-1}=1\right)=\mathrm{P}\left(s_{i t}=1 \mid \mathbf{z}_{i t}, s_{i, t-1}=1\right), \tag{19.122}
\end{equation*}


where $\mathbf{v}_{i t} \equiv\left(\mathbf{w}_{i t}, \mathbf{z}_{i t}\right)$. Now, we must include future values of $\mathbf{w}_{i t}$ and $\mathbf{z}_{i t}$ in the conditioning set on the left-hand side. Assumption (19.122) is fairly strong, but it does allow for attrition to be strongly related to past outcomes on $y$ and $\mathbf{x}$ (which can be included in $\mathbf{z}_{i t}$ ).

It is easy to see how assumption (19.122) leads to equation (19.120). First, with $\mathbf{v}_{i}=\left(\mathbf{v}_{i 1}, \mathbf{v}_{i 2}, \ldots, \mathbf{v}_{i T}\right)$, we can always write


\begin{align*}
\mathrm{P}\left(s_{i t}\right. & \left.=1 \mid \mathbf{v}_{i}\right)=\mathrm{P}\left(s_{i t}=1 \mid \mathbf{v}_{i}, s_{i, t-1}=1\right) \cdots \mathrm{P}\left(s_{i 2}=1 \mid \mathbf{v}_{i}, s_{i 1}=1\right) \mathrm{P}\left(s_{i 1}=1 \mid \mathbf{v}_{i}\right) \\
& =\mathrm{P}\left(s_{i t}=1 \mid \mathbf{v}_{i}, s_{i, t-1}=1\right) \cdots \mathrm{P}\left(s_{i 2}=1 \mid \mathbf{v}_{i}\right) \tag{19.123}
\end{align*}


because we assume $s_{i 1} \equiv 1$. Under assumption (19.122), $\mathrm{P}\left(s_{i t}=1 \mid \mathbf{v}_{i}, s_{i, t-1}=1\right)= \mathrm{P}\left(s_{i t}=1 \mid \mathbf{z}_{i t}, s_{i, t-1}=1\right)$ for $t \geq 2$, and so\\
$p_{i t}^{o} \equiv \mathrm{P}\left(s_{i t}=1 \mid \mathbf{v}_{i}\right)=\mathrm{P}\left(s_{i t}=1 \mid \mathbf{z}_{i t}, s_{i, t-1}=1\right) \cdots \mathrm{P}\left(s_{i 2}=1 \mid \mathbf{z}_{i 2}\right)$,\\
which is exactly equation (19.121).\\
As in the cross section case, we can use the "surprising" efficiency result described in Section 13.10.2 to show that, under assumption (19.122), estimating the attrition probabilities is actually more efficient than using the known probabilities (if we could). The key is that we have fully specified $\mathrm{D}\left(s_{i 1}, s_{i 2}, \ldots, s_{i T} \mid \mathbf{v}_{i}\right)$. Because $\left\{\mathbf{w}_{i t}: t=\right. 1, \ldots, T\}$ is contained in $\mathbf{v}_{i}$, the conditional independence assumption (13.66) (with an appropriate change in notation) holds. To realize the efficiency gain in the computed variance matrix estimator, we partial out from the weighted M-estimator score the score from the first-stage MLE. Define\\
$k\left(\mathbf{s}_{i}, \mathbf{z}_{i}, \mathbf{w}_{i}, \boldsymbol{\gamma}, \boldsymbol{\theta}\right) \equiv \sum_{t=1}^{T}\left[s_{i t} / p_{i t}\left(\boldsymbol{\delta}_{t}\right)\right] \mathbf{r}_{t}\left(\mathbf{w}_{i t}, \boldsymbol{\theta}\right)$,\\
where $\mathbf{r}_{t}\left(\mathbf{w}_{i t}, \boldsymbol{\theta}\right)=\nabla_{\boldsymbol{\theta}} q\left(\mathbf{w}_{i t}, \boldsymbol{\theta}\right)$ is the $\mathrm{P} \times 1$ score of the unweighted objective function. Further, let $\mathbf{d}_{i}(\boldsymbol{\delta})$ denote the score of the conditional log likelihood

$$
\sum_{t=2}^{T} s_{i, t-1}\left\{s_{i t} \log \left(\pi_{t}\left(\mathbf{z}_{i t}, \gamma_{t}\right)+\left(1-s_{i t}\right) \log \left[1-\pi_{t}\left(\mathbf{z}_{i t}, \gamma_{t}\right)\right]\right\}\right.
$$

where $\boldsymbol{\delta}$ denotes the vector of $\gamma_{t}, t=2, \ldots, T$. Typically, in each time period we would estimate a logit or probit. Then, $\mathbf{d}_{i}(\boldsymbol{\delta})$ is just the long column vector of stacked scores for each logit or probit, selected at time $t$ so that we use only the units in the sample at $t-1$. That is,\\
$\mathbf{d}_{i}(\boldsymbol{\delta})^{\prime}=\left[s_{i 1} \mathbf{g}_{2}\left(s_{i 2}, \mathbf{z}_{i 2} ; \gamma_{2}\right), s_{i 2} \mathbf{g}_{3}\left(s_{i 3}, \mathbf{z}_{i 3} ; \gamma_{3}\right), \ldots, s_{i, T-1} \mathbf{g}_{T}\left(s_{i T}, \mathbf{z}_{i T} ; \gamma_{T}\right)\right]$,\\
where $\mathbf{g}_{t}\left(s_{i t}, \mathbf{z}_{i t} ; \boldsymbol{\gamma}_{t}\right)$ is just the gradient (row vector) for the binary response model at time $t$, using units in the sample at $t-1$. Then, as in the cross section case, let\\
$\hat{\mathbf{e}}_{i} \equiv \hat{\mathbf{k}}_{i}-\left(N^{-1} \sum_{i=1}^{N} \hat{\mathbf{k}}_{i} \hat{\mathbf{d}}_{i}^{\prime}\right)\left(N^{-1} \sum_{i=1}^{N} \hat{\mathbf{d}}_{i} \hat{\mathbf{d}}_{i}^{\prime}\right)^{-1} \hat{\mathbf{d}}_{i}$\\
be the $\mathrm{P} \times 1$ residuals from the multivariate regression of $\hat{\mathbf{k}}_{i}$ on $\hat{\mathbf{d}}_{i}, i=1, \ldots, N$, where all hatted quantities are evaluated at $\hat{\boldsymbol{\delta}}$ or $\hat{\boldsymbol{\theta}}_{w}$ (where the latter is the IPW M-estimator). Further, define\\
$\hat{\mathbf{D}} \equiv N^{-1} \sum_{i=1}^{N} \hat{\mathbf{e}}_{i} \hat{\mathbf{e}}_{i}^{\prime}$\\
and\\
$\hat{\mathbf{A}} \equiv N^{-1} \sum_{i=1}^{N} \sum_{t=1}^{T}\left[s_{i} / p_{i t}\left(\hat{\boldsymbol{\delta}}_{t}\right)\right] H_{t}\left(\mathbf{w}_{i t}, \hat{\boldsymbol{\theta}}_{w}\right)$.\\
The asymptotic variance of $\hat{\boldsymbol{\theta}}_{w}$ is consistently estimated as $\hat{\mathbf{A}}^{-1} \hat{\mathbf{D}} \hat{\mathbf{A}}^{-1} / N$, which allows for general serial dependence in the scores over time as well as violation of the information matrix equality in each time period. (So, for example, with nonlinear regression, or QMLE in the linear exponential family, the conditional variances may be arbitrary.)

An alternative, as usual with large $N$ and small $T$, is to include both steps of the two-step procedure in a panel bootstrapping routine. The obtained standard errors will properly reflect the increased efficiency from estimating the $\gamma_{t}$.

RRZ (1995) show that, in the case of nonlinear regression where the explanatory variables are always observed - a leading case occurs when the covariates are observed in the first period and do not change over time-condition (19.122) can be relaxed somewhat, and the IPW method still produces consistent, asymptotically normal estimators. When modeling a feature of $\mathrm{D}\left(\mathbf{y}_{i t} \mid \mathbf{x}_{i t}\right)$ when the covariates are not always observed, IPW suffers from the same problem we discussed in Section 19.8. Namely, if the feature of $\mathrm{D}\left(\mathbf{y}_{i t} \mid \mathbf{x}_{i t}\right)$ is correctly specified (and the objective function has been properly chosen), and if selection is a function of the covariates in the sense that\\
$\mathrm{P}\left(s_{i t}=1 \mid \mathbf{x}_{i t}, y_{i t}, s_{i, t-1}=1\right)=\mathrm{P}\left(s_{i t}=1 \mid \mathbf{x}_{i t}, s_{i, t-1}=1\right)$,\\
then the unweighted pooled M-estimator is consistent. If $\mathbf{x}_{i t}$ is not fully observed at time $t$, then it cannot be included in $\mathbf{z}_{i t}$, and condition (19.122) is very likely to fail.

As in the case of the Heckman approach applied to a first-differenced linear equation, the weighting is not innocuous: it would actually lead to inconsistent estimation. Thus, in the case of attrition in the leading case where interest lies in a feature of $\mathrm{D}\left(\mathbf{y}_{i t} \mid \mathbf{x}_{i t}\right)$, and some elements of $\mathbf{x}_{i t}$ are unobserved in time $t$ for those leaving the sample, one must be careful in interpreting important differences between the weighted and unweighted estimators: the problem might be with the weighted estimator. (Of course, both estimators might be inconsistent, too.)

\section*{Problems}
19.1. In the setup of Section 19.2.1, suppose that the costs, $r_{i}$, do not vary across $i$. What features of $\mathrm{D}\left(y_{i} \mid \mathbf{x}_{i}\right)$ can we consistently estimate?\\
19.2. Some occupations, such as major league baseball, are characterized by salary floors. This situation can be described by\\
$y=\exp (\mathbf{x} \boldsymbol{\beta}+u), \quad u \mid \mathbf{x} \sim \operatorname{Normal}\left(0, \sigma^{2}\right)$\\
$w=\max (f, y)$,\\
where $f>0$ is the common salary floor (the minimum wage), $y$ is the person's true worth (productivity), and $\mathbf{x}$ contains human capital and demographic variables.\\
a. Find the log-likelihood function for a random draw $i$ from the population.\\
b. How would you estimate $\mathrm{E}(y \mid \mathbf{x})$ ?\\
c. Is $\mathrm{E}(w \mid \mathbf{x})$ of much interest in this application? Explain.\\
19.3. Let $y$ be the percentage of annual income invested in a pension plan, and assume that current law caps this percentage at 10 percent. Therefore, in a sample of data, we observe $y_{i}$ between zero and 10 , with pileups at the two corners, zero and 10 .\\
a. What model would you use for $y$ that recognizes the pileups at zero and 10 ?\\
b. Explain the conceptual difference between the outcomes $y=0$ and $y=10$. In particular, which limit can be viewed as a form of data censoring?\\
c. How would you estimate the partial effects on the expected contribution percentage assuming that the current law always will be in effect?\\
d. Suppose you want to ask, What would be the effect on $\mathrm{E}(y \mid \mathbf{x})$, for any value of $\mathbf{x}$, if the contribution cap were increased from 10 to 11 ? How would you estimate the effect? [Hint: In the general two-limit Tobit model, call the upper bound, in general,\\
$a_{2}$ (which is known, not a parameter to estimate), and then take the partial derivative with respect to $a_{2}$.]\\
e. If there are no observations at $y=10$, what does the estimated model reduce to? Does this finding seem sensible?\\
19.4. a. Suppose you are hired to explain fire damage to buildings in terms of building and neighborhood characteristics. If you use cross section data on reported fires, is there a sample selection problem due to the fact that most buildings do not catch fire during the year?\\
b. If you want to estimate the relationship between contributions to a $401(\mathrm{k})$ plan and the match rate of the plan-the rate at which the employer matches employee contributions-is there a sample selection problem if you only use a sample of workers already enrolled in a 401(k) plan?\\
19.5. In example 19.4, suppose that $I Q$ is an indicator of abil, and $K W W$ is another indicator (see Section 5.3.2). Find assumptions under which IV on the selected sample is valid.\\
19.6. Let $f\left(\cdot \mid \mathbf{x}_{i} ; \boldsymbol{\theta}\right)$ denote the density of $y_{i}$ given $\mathbf{x}_{i}$ for a random draw from the population. Find the conditional density of $y_{i}$ given $\left(\mathbf{x}_{i}, s_{i}=1\right)$ when the selection rule is $s_{i}=1\left[a_{1}\left(\mathbf{x}_{i}\right)<y_{i}<a_{2}\left(\mathbf{x}_{i}\right)\right]$, where $a_{1}(\mathbf{x})$ and $a_{2}(\mathbf{x})$ are known functions of $\mathbf{x}$. In the Hausman and Wise (1977) example, $a_{2}(\mathbf{x})$ was a function of family size because the poverty income level depends on family size.\\
19.7. Suppose in Section 19.6.1 we replace assumption 19.1d with\\
$\mathrm{E}\left(u_{1} \mid v_{2}\right)=\gamma_{1} v_{2}+\gamma_{2}\left(v_{2}^{2}-1\right)$.\\
(We subtract unity from $v_{2}^{2}$ to ensure that the second term has zero expectation.)\\
a. Using the fact that $\operatorname{Var}\left(v_{2} \mid v_{2}>-a\right)=1-\lambda(a)[\lambda(a)+a]$, show that\\
$\mathrm{E}\left(y_{1} \mid \mathbf{x}, y_{2}=1\right)=\mathbf{x}_{1} \boldsymbol{\beta}_{1}+\gamma_{1} \lambda\left(\mathbf{x} \boldsymbol{\delta}_{2}\right)-\gamma_{2} \lambda\left(\mathbf{x} \boldsymbol{\delta}_{2}\right) \mathbf{x} \boldsymbol{\delta}_{2}$.\\[0pt]
[Hint: Take $a=\mathbf{x} \boldsymbol{\delta}_{2}$ and use the fact that $\mathrm{E}\left(v_{2}^{2} \mid v_{2}>-a\right)=\operatorname{Var}\left(v_{2} \mid v_{2}>-a\right)+ \left[\mathrm{E}\left(v_{2} \mid v_{2}>-a\right)\right]^{2}$.]\\
b. Explain how to correct for sample selection in this case.\\
c. How would you test for the presence of sample selection bias?\\
19.8. Consider the following alternative to procedure 19.2 when $y_{2}$ is always observed. First, run the OLS regression of $y_{2}$ on $\mathbf{z}$ and obtain the fitted values, $\hat{y}_{2}$.

Next, get the inverse Mills ratio, $\hat{\lambda}_{3}$, from the probit of $y_{3}$ on $\mathbf{z}$. Finally, run the OLS regression $y_{1}$ on $\mathbf{z}_{1}, \hat{y}_{2}, \hat{\lambda}_{3}$ using the selected sample.\\
a. Find a set of sufficient conditions that imply consistency of the proposed procedure. (Do not worry about regularity conditions.)\\
b. Show that the assumptions from part a are more restrictive than those in procedure 19.2, and give some examples that are covered by procedure 19.2 but not by the alternative procedure.\\
19.9. Apply procedure 19.4 to the data in MROZ.RAW. Use a constant, exper, and exper $^{2}$ as elements of $\mathbf{z}_{1}$; take $y_{2}=$ educ. The other elements of $\mathbf{z}$ should include age, kidslt6, kidsge6, nwifeinc, motheduc, fatheduc, and huseduc.\\
19.10. Consider the model\\
$y_{1}=\mathbf{z} \boldsymbol{\delta}_{1}+v_{1}$\\
$y_{2}=\mathbf{z} \boldsymbol{\delta}_{2}+v_{2}$\\
$y_{3}=\max \left(0, \alpha_{31} y_{1}+\alpha_{32} y_{2}+\mathbf{z}_{3} \boldsymbol{\delta}_{3}+u_{3}\right)$,\\
where $\left(\mathbf{z}, y_{2}, y_{3}\right)$ are always observed and $y_{1}$ is observed when $y_{3}>0$. The first two equations are reduced-form equations, and the third equation is of primary interest. For example, take $y_{1}=\log \left(\right.$ wage $\left.^{\mathrm{o}}\right), y_{2}=e d u c$, and $y_{3}=h o u r s$, and then education and $\log \left(\right.$ wage $\left.^{\mathrm{o}}\right)$ are possibly endogenous in the labor supply function. Assume that $\left(v_{1}, v_{2}, u_{3}\right)$ are jointly zero-mean normal and independent of $\mathbf{z}$.\\
a. Find a simple way to consistently estimate the parameters in the third equation allowing for arbitrary correlations among $\left(v_{1}, v_{2}, u_{3}\right)$. Be sure to state any identification assumptions needed.\\
b. Now suppose that $y_{2}$ is observed only when $y_{3}>0$; for example, $y_{1}= \log \left(\right.$ wage $\left.^{\mathrm{o}}\right), y_{2}=\log \left(\right.$ benefits $\left.^{\mathrm{o}}\right), y_{3}=$ hours. Now derive a multistep procedure for estimating the third equation under the same assumptions as in part a.\\
c. How can we estimate the average partial effects?\\
19.11. Consider the following conditional moment restrictions problem with a selected sample. In the population, $\mathrm{E}\left[\mathbf{r}\left(\mathbf{w}, \boldsymbol{\theta}_{\mathrm{o}}\right) \mid \mathbf{x}\right]=\mathbf{0}$. Let $s$ be the selection indicator, and assume that\\
$\mathrm{E}\left[\mathbf{r}\left(\mathbf{w}, \boldsymbol{\theta}_{\mathrm{o}}\right) \mid \mathbf{x}, s\right]=\mathbf{0}$.\\
Sufficient is that $s=f(\mathbf{x})$ for a nonrandom function $f$.\\
a. Let $\mathbf{Z}_{i}$ be a $G \times L$ matrix of functions of $\mathbf{x}_{i}$. Show that $\boldsymbol{\theta}_{\mathrm{o}}$ satisfies $\mathrm{E}\left[s_{i} \mathbf{Z}_{i}^{\prime} \mathbf{r}\left(\mathbf{w}_{i}, \boldsymbol{\theta}_{\mathrm{o}}\right)\right]=\mathbf{0}$\\
b. Write down the objective function for the system nonlinear 2SLS estimator based on the selected sample. Argue that, under the appropriate rank condition, the estimator is consistent and $\sqrt{N}$-asymptotically normal.\\
c. Write down the objective function for a minimum chi-square estimator using the selected sample. Use the estimates from part $b$ to estimate the weighting matrix. Argue that the estimator is consistent and $\sqrt{N}$-asymptotically normal.\\
19.12. Consider model (19.56), where selection is ignorable in the sense that $\mathrm{E}\left(u_{1} \mid \mathbf{z}, u_{3}\right)=0$. However, data are missing on $y_{2}$ when $y_{3}=0$, and $\mathrm{E}\left(y_{2} \mid \mathbf{z}, y_{3}\right) \neq \mathrm{E}\left(y_{2} \mid \mathbf{z}\right)$.\\
a. Find $\mathrm{E}\left(y_{1} \mid \mathbf{z}, y_{3}\right)$.\\
b. If, in addition to assumption 19.2, $\left(v_{2}, v_{3}\right)$ is independent of $\mathbf{z}$ and $\mathrm{E}\left(v_{2} \mid v_{3}\right)= \gamma_{2} v_{3}$, find $\mathrm{E}\left(y_{1} \mid \mathbf{z}, y_{3}=1\right)$.\\
c. Suggest a two-step method for consistently estimating $\boldsymbol{\delta}_{1}$ and $\alpha_{1}$.\\
d. Does this method generally work if $\mathrm{E}\left(u_{1} \mid \mathbf{z}, y_{3}\right) \neq 0$ ?\\
e. Would you bother with the method from part c if $\mathrm{E}\left(u_{1} \mid \mathbf{z}, y_{2}, y_{3}\right)=0$ ? Explain.\\
19.13. In Section 17.6 we discussed two-part models for a corner solution outcome, say, $y$. These models have sometimes been studied in the context of incidental truncation.\\
a. Suppose you have a parametric model for the distribution of $y$ conditional on $\mathbf{x}$ and $y>0$. (Cragg's model and the lognormal model from Section 17.6 are examples.) If you estimate the parameters of this model by conditional MLE, using only the observations for which $y_{i}>0$, do the parameter estimates suffer from sample selection bias? Explain.\\
b. If instead you specify only $\mathrm{E}(y \mid \mathbf{x}, y>0)=\exp (\mathbf{x} \boldsymbol{\beta})$ and estimate $\boldsymbol{\beta}$ by nonlinear least squares using observations for which $y_{i}>0$, do the estimates suffer from sample selection bias?\\
c. In addition to the specification from part b , suppose that $\mathrm{P}(y=0 \mid \mathbf{x})= 1-\Phi(\mathbf{x} \boldsymbol{\gamma})$. How would you estimate $\boldsymbol{\gamma}$ ?\\
d. Given the assumptions in parts b and c , how would you estimate $\mathrm{E}(y \mid \mathbf{x})$ ?\\
e. Given your answers to the first four parts, do you think viewing estimation of twopart models as an incidental truncation problem is appropriate?\\
19.14. Consider equation (19.30) under the assumption $\mathrm{E}(u \mid \mathbf{z}, s)=\mathrm{E}(u \mid s)= (1-s) \alpha_{0}+s \alpha_{1}$. The first equality is the assumption; the second is unrestrictive, as it simply allows the mean of $u$ to differ in the selected and unselected subpopulations.\\
a. Show that 2SLS estimation using the selected subsample consistently estimates the slope parameters, $\beta_{2}, \ldots, \beta_{K}$. What is the plim of the intercept estimator? (Hint: Replace $u$ with $(1-s) \alpha_{0}+s \alpha_{1}+e$, where $\mathrm{E}(e \mid \mathbf{z}, s)=0$.)\\
b. Show that $\mathrm{E}(u \mid \mathbf{z}, s)=\mathrm{E}(u \mid s)$ if ( $u, s$ ) is independent of $\mathbf{z}$. Does independence of $s$ and $\mathbf{z}$ seem reasonable?\\
19.15. Suppose that $y$ given $\mathbf{x}$ follows a standard censored Tobit, where $y$ is a corner solution response. However, there is at least one element of $\mathbf{x}$ that we can observe only when $y>0$. (An example is seen when $y$ is quantity demanded of a good or service, and one element of $\mathbf{x}$ is price, derived as total expenditure on the good divided by $y$ whenever $y>0$.)\\
a. Explain why we cannot use standard censored Tobit maximum likelihood estimation to estimate $\boldsymbol{\beta}$ and $\sigma^{2}$. What method can we use instead?\\
b. How is it that we can still estimate $\mathrm{E}(y \mid \mathbf{x})$, even though we do not observe some elements of $\mathbf{x}$ when $y=0$ ?\\
19.16. Consider a standard linear model with an endogenous explanatory variable:\\
$y_{1}=\mathbf{z}_{1} \boldsymbol{\delta}_{1}+\alpha_{1} y_{2}+u_{1}$\\
$y_{2}=\mathbf{z} \boldsymbol{\delta}_{2}+v_{2}$,\\
but where $y_{2}$ is censored from above. Let $r_{2}$ be the censoring variable (which may be random) and define $w_{2}=\min \left(r_{2}, y_{2}\right)$. Assume that ( $u_{1}, v_{2}$ ) is independent of ( $\mathbf{z}, r_{2}$ ) and $E\left(u_{1} \mid v_{2}\right)=\rho_{1} v_{2}$.\\
a. Show that\\
$\mathrm{E}\left(y_{1} \mid \mathbf{z}, r_{2}, v_{2}\right)=\mathbf{z}_{1} \boldsymbol{\delta}_{1}+\alpha_{1} y_{2}+\rho_{1} v_{2}$.\\
b. Define $s_{2}=1\left[y_{2}<w_{2}\right]$, so $s_{2}=1$ means $y_{2}$ is observed. How come\\
$\mathrm{E}\left(y_{1} \mid \mathbf{z}, r_{2}, v_{2}, s_{2}\right)=\mathrm{E}\left(y_{1} \mid \mathbf{z}, r_{2}, v_{2}\right)=\mathbf{z}_{1} \boldsymbol{\delta}_{1}+\alpha_{1} y_{2}+\rho_{1} v_{2} ?$\\
c. Propose a two-step estimator that consistently estimates $\boldsymbol{\delta}_{1}$ and $\alpha_{1}$ (along with $\rho_{1}$ ) where only the data with $y_{i 2}$ uncensored are used in the second step. How would you obtain valid standard errors?\\
19.17. Consider an unobserved effects panel data model with binary censoring:\\
$y_{i t}=\mathbf{x}_{i t} \boldsymbol{\beta}+c_{i}+u_{i t}$\\
$w_{i t}=1\left[y_{i t}>r_{i t}\right]$,\\
where $r_{i t}$ is the known censoring value for unit $i$ in time $t$. Our goal is to estimate $\boldsymbol{\beta}$, which, in the absense of censoring, we could achieve by fixed effects, first difference, random effects (under stronger assumptions), or some variation on these.\\
a. Maintain that the censoring values are strictly exogenous conditional on $c_{i}$ (and $\left.\mathbf{x}_{i}=\left(\mathbf{x}_{i 1}, \ldots, \mathbf{x}_{i T}\right)\right)$ along with normality:\\
$\mathrm{D}\left(u_{i t} \mid \mathbf{x}_{i}, \mathbf{r}_{i}, c_{i}\right)=\operatorname{Normal}\left(0, \sigma_{u}^{2}\right)$,\\
where $\mathbf{r}_{i}=\left(r_{i 1}, r_{i 2}, \ldots, r_{i T}\right)$. Explain what kinds of censoring this assumption rules out.\\
b. Allow censoring to depend on $c_{i}$ using the Chamberlain-Mundlak device:\\
$c_{i}=\psi+\overline{\mathbf{x}}_{i} \boldsymbol{\xi}+\eta \bar{r}_{i}+a_{i}$\\
$\mathrm{D}\left(a_{i} \mid \mathbf{x}_{i}, \mathbf{r}_{i}\right)=\operatorname{Normal}\left(0, \sigma_{a}^{2}\right)$,\\
where, as usual, the overbar means time average. In other words, heterogeneity may be (partially) correlated with the average censoring value. Using this assumption with that in part a, show that $\mathrm{P}\left(w_{i t}=1 \mid \mathbf{x}_{i}, \mathbf{r}_{i}\right)$ follows a probit model and obtain the coefficients on the various explanatory variables.\\
c. Explain how to consistently estimate $\boldsymbol{\beta}$ without further assumptions (other than that both $\left\{\mathbf{x}_{i t}\right\}$ and $\left\{r_{i t}\right\}$ have sufficient time variation).\\
d. What additional assumption would allow you to also estimate $\sigma_{u}^{2}$ and $\sigma_{a}^{2}$ ? Explain how you would use this assumption in estimation.\\
19.18. Suppose that in the population, $\mathrm{D}(y \mid \mathbf{x})$ follows a $\operatorname{Tobit}\left(\mathbf{x} \boldsymbol{\beta}, \sigma^{2}\right)$ distribution. However, we observe $y$ only when it is strictly positive. (Such a scenario is common for on-site sampling, where only individuals known to participate in an activity are sampled. An example is $y$ is annual days spent visiting national parks, and surveys are taken at national parks.)\\
a. How would you estimate $\boldsymbol{\beta}$ and $\sigma^{2}$ given a random sample from the subpopulation with $y>0$ ?\\
b. How would you estimate partial effects of the $x_{j}$ on $\mathrm{E}(y \mid \mathbf{x})$ ? How and why does this approach differ from the usual truncated regression setup?\\
c. Can you estimate a hurdle model of the kind covered in Section 17.6? Explain.\\
19.19. Consider a count random variable, $y_{i}$, with conditional probability density $f_{y}(\cdot \mid \mathbf{x})$. Suppose that $r_{i} \geq 0$ is a right censoring point, and assume that $\mathrm{P}\left(y_{i}=y \mid \mathbf{x}_{i}\right.$, $\left.r_{i}\right)=\mathrm{P}\left(y_{i}=y \mid \mathbf{x}_{i}\right)$. For a random draw $i$ from the population, we observe the random variable $w_{i}=\min \left(y_{i}, r_{i}\right)$.\\
a. Let $F_{y}(\cdot \mid \mathbf{x})$ denote the conditional cdf of $y_{i}$ given $\mathbf{x}_{i}=\mathbf{x}$. Find the density of $w_{i}$ given $\left(\mathbf{x}_{i}, r_{i}\right)$ for $w=0,1, \ldots, r_{i}$.\\
b. Find the density from part a if $y_{i}$ has a conditional Poisson distribution with mean $\exp \left(\mathbf{x}_{i} \boldsymbol{\beta}\right)$.\\
c. If you estimate $\boldsymbol{\beta}$ by maximum likelihood using the right censored data, is the resulting estimator generally robust to failure of the Poisson distribution? Explain.\\
d. Use your answer to part a to discuss the costs of data censoring for a count variable.\\
e. What other underlying population distribution might you use, and why?

\section*{\begin{center}

\end{document}
